% Dérivées et sens de variation — Mathématiques Tle A — Cours signé OnBuch+
% Programme officiel MINESEC. Compiler avec : tectonic MA03-derivees-a.tex
\newcommand{\DOCMATIERE}{Mathématiques}
\newcommand{\DOCNIVEAU}{Tle A}
\newcommand{\DOCMODULE}{Relations et opérations fondamentales dans l'ensemble des nombres réels}
\newcommand{\DOCLECON}{A03}
\newcommand{\DOCTITRE}{Dérivées et sens de variation}
% OnBuch+ — préambule « cours signés » ENRICHI (compilé avec tectonic / XeLaTeX)
% Système visuel « Pop » d'OnBuch+ : fond crème (jamais blanc), bordures encre
% épaisses, ombres « dures » décalées, titres Archivo Black en capitales.
% Version MATHÉMATIQUES : graphes (pgfplots/tikz), boîtes pédagogiques
% (définition, propriété, méthode, attention, le savais-tu, exercice résolu,
% exercice, TP, bilan), page de garde et signature OnBuch+ ancrée en bas.
%
% Macros attendues dans main.tex AVANT \input{preamble} :
%   \DOCMATIERE  \DOCNIVEAU  \DOCMODULE  \DOCLECON (numéro)  \DOCTITRE
\documentclass[11pt]{article}

\usepackage{fontspec}
\usepackage[french]{babel}
\usepackage[a4paper,margin=2cm,top=3.4cm,bottom=2.8cm,headheight=1.4cm,footskip=1.3cm]{geometry}
\usepackage{amsmath,amssymb,amsfonts}
\usepackage{mathtools} % \xmapsto, \coloneqq... souvent utilisés par les modèles
\usepackage{cancel} % simplifications barrées
% \abs{z} (module, valeur absolue) et \norm{u} (norme), utilisés par les modèles
\providecommand{\abs}{}\let\abs\relax\DeclarePairedDelimiter{\abs}{\lvert}{\rvert}
\providecommand{\norm}{}\let\norm\relax\DeclarePairedDelimiter{\norm}{\lVert}{\rVert}
\usepackage[table]{xcolor} % [table] : fournit aussi \rowcolors (alternance de lignes)
\usepackage{pagecolor}
\usepackage{tikz}
\usetikzlibrary{arrows.meta,positioning,calc,shapes.geometric,shapes.misc,shapes.arrows,decorations.pathmorphing,decorations.markings,patterns,shadows,mindmap,trees,fit,backgrounds,matrix,intersections}
\usepackage{pgfplots}
\usepackage{pgf-pie} % diagrammes circulaires \pie (statistiques)
\pgfplotsset{compat=1.18}
\usepgfplotslibrary{fillbetween,groupplots} % aires entre courbes (calcul intégral)
\usepackage{enumitem}
\usepackage{fancyhdr}
% [most] charge toutes les bibliothèques tcolorbox (listings, minted, raster,
% documentation...) et gonfle inutilement la mémoire TeX sur un document long
% (~300 boîtes) : on ne charge que ce qui est réellement utilisé.
\usepackage[skins,breakable]{tcolorbox}
\usepackage{graphicx}
\usepackage{colortbl}
\usepackage{booktabs}
\usepackage{tabularx}
\usepackage{diagbox} % case d'en-tête diagonale des tableaux à double entrée
% Colonnes extensibles centrée / gauche / droite (C, L, R), souvent utilisées par les modèles
\newcolumntype{C}{>{\centering\arraybackslash}X}
\newcolumntype{L}{>{\raggedright\arraybackslash}X}
\newcolumntype{R}{>{\raggedleft\arraybackslash}X}
\usepackage{array}
\usepackage{multicol}
\usepackage{siunitx}
% parse-numbers=false : accepte aussi \SI{2,5 \times 10^{-3}}{...}
\sisetup{locale=FR,output-decimal-marker={,},group-minimum-digits=4,parse-numbers=false}
\DeclareSIUnit\ohm{\ensuremath{\symup{\Omega}}} % U+2126 absent de la police
\usepackage{qrcode}
% \degree : commande fréquemment utilisée par les modèles pour un angle ou une
% température en dehors de \SI{}{} (non fournie par les paquets chargés).
\providecommand{\degree}{^{\circ}}
% \qed (fin de démonstration), utilisé par les modèles sans amsthm
\providecommand{\qed}{\hfill$\square$}
% \barre : double barre des valeurs interdites dans un tableau de variations (tkz-tab)
\providecommand{\barre}{\ensuremath{\|}}
% environnement proof (amsthm non chargé) utilisé par les modèles
\ifdefined\proof\else\newenvironment{proof}[1][Démonstration]{\par\noindent\textit{#1.}\ }{\qed\par}\fi
\usepackage{lastpage}
\usepackage{hyperref}
\hypersetup{hidelinks}

% ============================================================
% Palette « Pop » OnBuch+ — mêmes hex que la classe P (plus_ui.dart)
% ============================================================
\definecolor{poporange}{HTML}{F59321}
\definecolor{popdark}{HTML}{E07A0C}
\definecolor{popcream}{HTML}{FFF6E8}
\definecolor{popink}{HTML}{1C1714}
\definecolor{poppurple}{HTML}{7A5AE0}
\definecolor{popgreen}{HTML}{2E9E6A}
\definecolor{poppink}{HTML}{E0457B}
\definecolor{popblue}{HTML}{2F7DE1}
\definecolor{popgold}{HTML}{C98A12}
\definecolor{popmuted}{HTML}{8A7F76}
\definecolor{popline}{HTML}{EADFD2}
\definecolor{popgray}{HTML}{8A7F76}
\colorlet{popgrey}{popgray}
% Alias tolérants (le modèle nomme parfois les couleurs différemment)
\colorlet{popwhite}{white}
\colorlet{popblack}{popink}
\colorlet{popred}{poppink}
\colorlet{popyellow}{popgold}
% Teintes claires pour fonds de tableaux / zones
\colorlet{popblueL}{popblue!10!white}
\colorlet{poporangeL}{poporange!14!white}
\colorlet{popgreenL}{popgreen!12!white}
\colorlet{poppurpleL}{poppurple!12!white}
\colorlet{poppinkL}{poppink!12!white}
\colorlet{popgoldL}{popgold!14!white}

\pagecolor{popcream}

% ============================================================
% Typographie
% ============================================================
\defaultfontfeatures{Ligatures=TeX}
\setmainfont{PlusJakartaSans-Medium.ttf}[Path=../../../fonts/,BoldFont=PlusJakartaSans-Bold.ttf]
% Maths en sans-serif (Fira Math) avec chiffres et lettres droites de Plus
% Jakarta, pour que les formules se fondent dans le texte.
\usepackage[math-style=ISO,bold-style=ISO]{unicode-math}
\setmathfont{FiraMath-Regular.otf}[Path=../../../fonts/]
\setmathfont{PlusJakartaSans-Medium.ttf}[Path=../../../fonts/,range={up/num,up/{Latin,latin},\mathup}]
\usepackage{icomma} % virgule décimale sans espace en mode math
% Caractères mathématiques Unicode absents des polices de texte (Plus Jakarta,
% Archivo Black) : rendus par leur équivalent mathématique, en texte comme en
% formule — sinon ils s'affichent en carré vide (ex. « Arithmétique dans ℤ »).
\usepackage{newunicodechar}
\newunicodechar{ℝ}{\ensuremath{\mathbb{R}}}
\newunicodechar{ℤ}{\ensuremath{\mathbb{Z}}}
\newunicodechar{ℂ}{\ensuremath{\mathbb{C}}}
\newunicodechar{ℕ}{\ensuremath{\mathbb{N}}}
\newunicodechar{ℚ}{\ensuremath{\mathbb{Q}}}
\newunicodechar{𝔻}{\ensuremath{\mathbb{D}}}
\newunicodechar{α}{\ensuremath{\alpha}}
\newunicodechar{β}{\ensuremath{\beta}}
\newunicodechar{γ}{\ensuremath{\gamma}}
\newunicodechar{δ}{\ensuremath{\delta}}
\newunicodechar{ε}{\ensuremath{\varepsilon}}
\newunicodechar{θ}{\ensuremath{\theta}}
\newunicodechar{λ}{\ensuremath{\lambda}}
\newunicodechar{μ}{\ensuremath{\mu}}
\newunicodechar{σ}{\ensuremath{\sigma}}
\newunicodechar{τ}{\ensuremath{\tau}}
\newunicodechar{φ}{\ensuremath{\varphi}}
\newunicodechar{ψ}{\ensuremath{\psi}}
\newunicodechar{ω}{\ensuremath{\omega}}
\newunicodechar{Σ}{\ensuremath{\Sigma}}
% majuscules produites par \MakeUppercase dans les titres (α-aminés -> Α)
\newunicodechar{Α}{\ensuremath{\alpha}}
\newunicodechar{Β}{\ensuremath{\beta}}
\newunicodechar{Γ}{\ensuremath{\Gamma}}
\newunicodechar{Δ}{\ensuremath{\Delta}}
\newunicodechar{Ε}{\ensuremath{\varepsilon}}
\newunicodechar{Θ}{\ensuremath{\Theta}}
\newunicodechar{Λ}{\ensuremath{\Lambda}}
\newunicodechar{Μ}{\ensuremath{\mu}}
\newunicodechar{Τ}{\ensuremath{\tau}}
\newunicodechar{Φ}{\ensuremath{\Phi}}
\newunicodechar{Ψ}{\ensuremath{\Psi}}
\newunicodechar{Ω}{\ensuremath{\Omega}}
\newunicodechar{↦}{\ensuremath{\mapsto}}
\newunicodechar{∈}{\ensuremath{\in}}
\newunicodechar{∉}{\ensuremath{\notin}}
\newunicodechar{∧}{\ensuremath{\wedge}}
\newunicodechar{⇔}{\ensuremath{\Leftrightarrow}}
\newunicodechar{⟺}{\ensuremath{\Longleftrightarrow}}
\newunicodechar{⇒}{\ensuremath{\Rightarrow}}
\newunicodechar{⟹}{\ensuremath{\Longrightarrow}}
\newunicodechar{∘}{\ensuremath{\circ}}
\newunicodechar{∪}{\ensuremath{\cup}}
\newunicodechar{∩}{\ensuremath{\cap}}
\newunicodechar{≡}{\ensuremath{\equiv}}
\newunicodechar{⊂}{\ensuremath{\subset}}
\newunicodechar{⊥}{\ensuremath{\perp}}
\newunicodechar{∃}{\ensuremath{\exists}}
\newunicodechar{∀}{\ensuremath{\forall}}
\newunicodechar{∛}{\ensuremath{\sqrt[3]{\,}}}
\newunicodechar{∜}{\ensuremath{\sqrt[4]{\,}}}
\newunicodechar{⃗}{\ensuremath{{}^{\to}}}
\newunicodechar{ⁿ}{\ifmmode^{n}\else\textsuperscript{\ensuremath{n}}\fi}
\newunicodechar{ˣ}{\ifmmode^{x}\else\textsuperscript{\ensuremath{x}}\fi}
\newunicodechar{ᵇ}{\ifmmode^{b}\else\textsuperscript{\ensuremath{b}}\fi}
\newunicodechar{ʳ}{\ifmmode^{r}\else\textsuperscript{\ensuremath{r}}\fi}
\newunicodechar{ᵘ}{\ifmmode^{u}\else\textsuperscript{\ensuremath{u}}\fi}
\newunicodechar{ʸ}{\ifmmode^{y}\else\textsuperscript{\ensuremath{y}}\fi}
\newunicodechar{ᵃ}{\ifmmode^{a}\else\textsuperscript{\ensuremath{a}}\fi}
\newunicodechar{ᵉ}{\ifmmode^{e}\else\textsuperscript{\ensuremath{e}}\fi}
\newunicodechar{ᵏ}{\ifmmode^{k}\else\textsuperscript{\ensuremath{k}}\fi}
\newunicodechar{ᵖ}{\ifmmode^{p}\else\textsuperscript{\ensuremath{p}}\fi}
\newunicodechar{ᴺ}{\ifmmode^{N}\else\textsuperscript{\ensuremath{N}}\fi}
\newunicodechar{ᶜ}{\ifmmode^{c}\else\textsuperscript{\ensuremath{c}}\fi}
\newunicodechar{ʰ}{\ifmmode^{h}\else\textsuperscript{\ensuremath{h}}\fi}
\newunicodechar{ᵠ}{\ifmmode^{q}\else\textsuperscript{\ensuremath{q}}\fi}
\newunicodechar{ₙ}{\ifmmode_{n}\else\textsubscript{\ensuremath{n}}\fi}
\newunicodechar{ₐ}{\ifmmode_{a}\else\textsubscript{\ensuremath{a}}\fi}
\newunicodechar{ᵢ}{\ifmmode_{i}\else\textsubscript{\ensuremath{i}}\fi}
\newunicodechar{ᵣ}{\ifmmode_{r}\else\textsubscript{\ensuremath{r}}\fi}
\newunicodechar{ₖ}{\ifmmode_{k}\else\textsubscript{\ensuremath{k}}\fi}
\newunicodechar{ᵦ}{\ifmmode_{\beta}\else\textsubscript{\ensuremath{\beta}}\fi}
\newunicodechar{ₓ}{\ifmmode_{x}\else\textsubscript{\ensuremath{x}}\fi}
\newfontfamily\popdisplay{ArchivoBlack-Regular.ttf}[Path=../../../fonts/]
\newfontfamily\poplogo{SpaceGrotesk-Bold.ttf}[Path=../../../fonts/]
\newfontfamily\popsemi{PlusJakartaSans-SemiBold.ttf}[Path=../../../fonts/]
\newfontfamily\popxbold{PlusJakartaSans-ExtraBold.ttf}[Path=../../../fonts/]
\newfontfamily\popxboldls{PlusJakartaSans-ExtraBold.ttf}[Path=../../../fonts/,LetterSpace=3.0]

\setlist[enumerate]{leftmargin=1.7em,itemsep=5pt,topsep=4pt}
\setlist[itemize]{leftmargin=1.7em,itemsep=4pt,topsep=4pt,label={\color{poporange}\textbullet}}
\setlength{\parindent}{0pt}
\setlength{\parskip}{5pt}
\renewcommand{\arraystretch}{1.25}

% pgfplots : style Pop par défaut
\pgfplotsset{
  every axis/.append style={axis line style={popink,line width=1pt},tick style={popink},
    grid style={popline,line width=0.5pt},label style={font=\small},tick label style={font=\footnotesize}},
  popcurve/.style={poporange,line width=1.8pt,no markers,smooth},
}
% Même style disponible dans un \draw TikZ simple (hors pgfplots)
\tikzset{popcurve/.style={poporange,line width=1.8pt}}

% ============================================================
% Pill / squiggle
% ============================================================
\newcommand{\poppill}[3]{%
  \tikz[baseline=-0.55ex]\node[draw=popink,line width=1.2pt,rounded corners=2.6pt,%
    fill=#2,inner xsep=6.5pt,inner ysep=3pt]{%
    \popxboldls\fontsize{7}{7}\selectfont\color{#3}\MakeUppercase{#1}};%
}
\newcommand{\popsquiggle}[1]{%
  \tikz[baseline=-0.4ex]{
    \draw[#1,line width=2.6pt,line cap=round]
      plot [smooth] coordinates {(0,0.09) (0.34,0.01) (0.68,0.21) (1.02,0.01) (1.36,0.10)};
  }%
}
% Mot-clé surligné (marqueur orange sous le texte)
\newcommand{\cle}[1]{{\popxbold\color{popdark}#1}}

% ============================================================
% En-tête / pied
% ============================================================
\pagestyle{fancy}
\fancyhf{}
\renewcommand{\headrulewidth}{0pt}
\renewcommand{\footrulewidth}{0pt}
\lhead{\raisebox{-0.15ex}{\poplogo\fontsize{13}{13}\selectfont\color{popink}OnBuch\color{poporange}+}}
\rhead{\poppill{\DOCMATIERE\ \textbullet\ \DOCNIVEAU\ \textbullet\ Leçon \DOCLECON}{white}{popink}}
\lfoot{\popxbold\fontsize{7.6}{7.6}\selectfont\color{popmuted}\MakeUppercase{Cours signé OnBuch+}\ \textbullet\ {\popsemi\DOCTITRE}}
\rfoot{\tikz[baseline=-0.6ex]\node[rounded corners=8.5pt,fill=popink,minimum height=17pt,minimum width=17pt,inner xsep=5pt,inner sep=0]{\popxbold\fontsize{8}{8}\selectfont\color{popcream}\thepage\,/\,\pageref*{LastPage}};}

% ============================================================
% Titres
% ============================================================
\newcommand{\chaptitre}[2]{%
  \begin{tcolorbox}[enhanced,colback=poporange,colframe=popink,boxrule=2.6pt,arc=11pt,
      drop shadow=popink,left=15pt,right=15pt,top=12pt,bottom=11pt,before skip=3pt,after skip=18pt]
    \poppill{#2}{popink}{popcream}\par\vspace{9pt}
    {\raggedright\hyphenpenalty=10000\exhyphenpenalty=10000\popdisplay\fontsize{22}{24}\selectfont\color{popcream}\MakeUppercase{#1}\par}
  \end{tcolorbox}
}
% Section numérotée : pastille orange avec le numéro + titre Archivo
\newcounter{popsec}
\newcommand{\coursec}[1]{%
  \stepcounter{popsec}\setcounter{popsub}{0}%
  \par\vspace{16pt}\noindent
  \begin{minipage}{\linewidth}
  \tikz[baseline=-0.7ex]\node[draw=popink,line width=1.6pt,fill=poporange,rounded corners=4pt,minimum size=22pt,inner sep=0]{\popdisplay\fontsize{12}{12}\selectfont\color{popcream}\thepopsec};%
  \hspace{8pt}{\popdisplay\fontsize{15}{17}\selectfont\color{popink}\MakeUppercase{#1}}\par
  \vspace{4pt}\noindent{\color{poporange}\rule{36pt}{3.6pt}}
  \end{minipage}\par\nopagebreak\vspace{8pt}
}
\newcounter{popsub}
\newcommand{\courssub}[1]{%
  \stepcounter{popsub}%
  \par\vspace{9pt}\noindent{\popxbold\fontsize{11.5}{13}\selectfont\color{popink}{\color{poporange}\thepopsec.\thepopsub}\hspace{6pt}#1}\par\nopagebreak\vspace{4pt}
}

% ============================================================
% Boîtes pédagogiques — même recette : fond blanc, bordure épaisse colorée,
% ombre dure encre, badge plein. Titre optionnel [#1].
% ============================================================
\tcbset{popbase/.style={breakable,enhanced,colback=white,boxrule=2.3pt,arc=9pt,
  shadow={3pt}{-3pt}{0pt}{popink},left=14pt,right=14pt,top=11pt,bottom=11pt,before skip=11pt,after skip=15pt,
  fontupper=\popsemi\fontsize{10.6}{14.5}\selectfont\color{popink},
  fontlower=\popsemi\fontsize{10.6}{14.5}\selectfont\color{popink}}}
% Coche dessinée (le glyphe ✓ n'existe pas dans les polices du document)
\newcommand{\popcheck}[1]{\tikz[baseline=-0.5ex]\draw[#1,line width=1.6pt,line cap=round,line join=round](-0.13,0.0)--(-0.04,-0.09)--(0.13,0.1);}
\newcommand{\popboxhead}[3]{\poppill{#1}{#2}{white}\ifx\relax#3\relax\else\hspace{6pt}{\popxbold\fontsize{10.6}{12}\selectfont\color{popink}#3}\fi\par\vspace{7pt}}

\newtcolorbox{aretenir}[1][]{popbase,colframe=popgreen,before upper={\popboxhead{À retenir}{popgreen}{#1}}}
\newtcolorbox{exemplebox}[1][]{popbase,colframe=poporange,before upper={\popboxhead{Exemple}{poporange}{#1}}}
\newtcolorbox{definition}[1][]{popbase,colframe=popblue,before upper={\popboxhead{Définition}{popblue}{#1}}}
\newtcolorbox{propriete}[1][]{popbase,colframe=poppurple,before upper={\popboxhead{Propriété}{poppurple}{#1}}}
\newtcolorbox{methode}[1][]{popbase,colframe=popgold,before upper={\popboxhead{Méthode}{popgold}{#1}}}
\newtcolorbox{attention}[1][]{popbase,colframe=poppink,before upper={\popboxhead{Attention}{poppink}{#1}}}
\newtcolorbox{experience}[1][]{popbase,colframe=popblue,colback=popblueL,before upper={\popboxhead{Expérience}{popblue}{#1}}}
% « Le savais-tu ? » : carte violette pleine (contexte, culture, Cameroun)
\newtcolorbox{savaistu}[1][]{popbase,colback=poppurple,colframe=popink,
  fontupper=\popsemi\fontsize{10.4}{14.2}\selectfont\color{white},
  before upper={\poppill{Le savais-tu\,?}{white}{poppurple}\ifx\relax#1\relax\else\hspace{6pt}{\popxbold\color{white}#1}\fi\par\vspace{7pt}}}
% Exercice résolu : énoncé en haut, \tcblower puis la solution sur fond crème
\newcounter{popexo}\newcounter{popexr}
\newtcolorbox{exoresolu}[1][]{popbase,colframe=poporange,colbacklower=popcream,
  segmentation style={popink,line width=1.2pt,dashed},
  before upper={\stepcounter{popexr}\popboxhead{Exercice résolu \thepopexr}{poporange}{#1}},
  before lower={\poppill{Solution}{popink}{popcream}\par\vspace{6pt}}}
% Exercice (non résolu en place) — la correction est dans la partie Corrigés
\newtcolorbox{exercice}[1][]{popbase,colframe=popink,
  before upper={\stepcounter{popexo}\popboxhead{Exercice \thepopexo}{popink}{#1}}}
% Corrigé
\newtcolorbox{corrige}[1][]{popbase,colframe=popgreen,colback=popgreenL,
  before upper={\popboxhead{Corrigé}{popgreen}{#1}}}
% Objectifs / prérequis / bilan
\newtcolorbox{objectifs}{popbase,colframe=popink,colback=poporangeL,
  before upper={\popboxhead{Objectifs de la leçon}{popink}{}}}
\newtcolorbox{prerequis}{popbase,colframe=popink,colback=popblueL,
  before upper={\popboxhead{Prérequis}{popblue}{}}}
\newtcolorbox{bilan}{popbase,colframe=popink,colback=white,boxrule=2.8pt,
  before upper={\popboxhead{Fiche bilan}{poporange}{L'essentiel à connaître pour l'examen}}}
% Figure encadrée avec légende
\newtcolorbox{popfigure}[1][]{enhanced,breakable=false,colback=white,colframe=popline,boxrule=1.4pt,arc=8pt,
  left=10pt,right=10pt,top=10pt,bottom=8pt,before skip=10pt,after skip=12pt,halign=center,
  lower separated=false,halign lower=center,
  fontlower=\popsemi\fontsize{9}{11.5}\selectfont\color{popmuted},
  #1}
\newcommand{\legende}[1]{\par\vspace{6pt}{\popsemi\fontsize{9}{11.5}\selectfont\color{popmuted}#1}}

% ============================================================
% Signature OnBuch+
% ============================================================
\newcommand{\poplogoinline}[1]{{\poplogo\fontsize{#1}{#1}\selectfont\color{popink}OnBuch\color{poporange}+}}
\newcommand{\signatureonbuch}{%
  \begin{tcolorbox}[enhanced,colback=popink,colframe=popink,boxrule=2.2pt,arc=10pt,
      drop shadow=poporange,left=14pt,right=14pt,top=10pt,bottom=10pt,before skip=0pt,after skip=0pt]
    \tikz[baseline=-0.7ex]\node[circle,fill=popgreen,draw=popcream,line width=1.3pt,minimum size=22pt,inner sep=0]{\popcheck{white}};%
    \hspace{9pt}\begin{minipage}[c]{0.62\linewidth}
      {\popxbold\fontsize{12}{13}\selectfont\color{popcream}Signé\ \textbullet\ {\poplogo OnBuch\color{poporange}+}}\par\vspace{2pt}
      {\raggedright\popsemi\fontsize{8}{10.5}\selectfont\color{popline}\MakeUppercase{Équipe pédagogique OnBuch+}\\\MakeUppercase{Conforme au programme officiel MINESEC}\par}
    \end{minipage}\hfill
    {\poplogo\fontsize{22}{22}\selectfont\color{popcream}OnBuch\color{poporange}+}
  \end{tcolorbox}%
}

% ============================================================
% Page de garde
% #1 titre · #2 matière · #3 niveau · #4 module · #5 n° leçon · #6 durée ·
% #7 description / objectifs (texte ou itemize)
% ============================================================
\newcommand{\pagegarde}[7]{%
  \thispagestyle{empty}
  \newgeometry{a4paper,margin=1.7cm,top=1.6cm,bottom=1.4cm}%
  \begin{tikzpicture}[remember picture,overlay]
    \fill[poporange!18] ([xshift=-3.2cm,yshift=-3.6cm]current page.north east) circle (4.6cm);
    \fill[poppurple!14] ([xshift=2.4cm,yshift=7.6cm]current page.south west) circle (3.8cm);
  \end{tikzpicture}%
  {\poplogo\fontsize{30}{30}\selectfont\color{popink}OnBuch\color{poporange}+}\hfill
  \raisebox{0.6ex}{\poppill{Cours signé \textbullet\ Premium}{popink}{popcream}}\par
  \vspace{1pt}\popsquiggle{poppurple}\par
  \vspace{8pt}
  \begin{tcolorbox}[enhanced,colback=poporange,colframe=popink,boxrule=2.8pt,arc=13pt,
      drop shadow=popink,left=18pt,right=18pt,top=12pt,bottom=13pt,after skip=10pt]
    \poppill{#2 \textbullet\ #3}{popink}{popcream}\hspace{5pt}\poppill{#4}{white}{popink}\par\vspace{8pt}
    {\popdisplay\fontsize{14}{14}\selectfont\color{popink}LEÇON #5}\par\vspace{4pt}
    {\raggedright\hyphenpenalty=10000\exhyphenpenalty=10000\popdisplay\fontsize{25}{27}\selectfont\color{popcream}\MakeUppercase{#1}\par}
    \vspace{8pt}
    {\popxbold\fontsize{9.5}{9.5}\selectfont\color{popink}\MakeUppercase{Durée conseillée : #6}}
  \end{tcolorbox}
  \begin{tcolorbox}[enhanced,colback=white,colframe=popink,boxrule=2.2pt,arc=10pt,
      drop shadow=popink,left=16pt,right=16pt,top=10pt,bottom=10pt,after skip=8pt]
    \poppill{Dans cette leçon}{poppurple}{white}\par\vspace{5pt}
    \popsemi\fontsize{9.3}{12.6}\selectfont\color{popink}#7
  \end{tcolorbox}
  \noindent\begin{minipage}[t]{0.487\linewidth}
    \begin{tcolorbox}[enhanced,colback=popgreen,colframe=popink,boxrule=2.2pt,arc=10pt,
        drop shadow=popink,left=11pt,right=11pt,top=10pt,bottom=10pt,height=3.1cm,valign=center]
      \tcbox[colback=white,colframe=popink,boxrule=1.3pt,arc=5pt,boxsep=0pt,nobeforeafter,
        left=3pt,right=3pt,top=3pt,bottom=3pt,valign=center]{\qrcode[height=1.9cm]{https://play.google.com/store/apps/details?id=com.luvvix.onbuch&hl=fr}}%
      \hspace{8pt}\begin{minipage}[c]{0.5\linewidth}
        {\popdisplay\fontsize{11}{12.5}\selectfont\color{white}\MakeUppercase{Télécharge OnBuch}}\par\vspace{4pt}
        {\raggedright\popsemi\fontsize{8.4}{11}\selectfont\color{white}Gratuit \textbullet\ Google Play\\Tuteur IA, annales, concours, résultats\par}
      \end{minipage}
    \end{tcolorbox}
  \end{minipage}\hfill
  \begin{minipage}[t]{0.487\linewidth}
    \begin{tcolorbox}[enhanced,colback=poppurple,colframe=popink,boxrule=2.2pt,arc=10pt,
        drop shadow=popink,left=11pt,right=11pt,top=10pt,bottom=10pt,height=3.1cm,valign=center]
      \tikz[baseline=-0.7ex]\node[circle,fill=white,draw=popink,line width=1.3pt,minimum size=1.6cm,inner sep=0]{\popdisplay\fontsize{24}{24}\selectfont\color{poppurple}+};%
      \hspace{8pt}\begin{minipage}[c]{0.6\linewidth}
        {\popdisplay\fontsize{11}{12.5}\selectfont\color{white}\MakeUppercase{Passe à OnBuch+}}\par\vspace{4pt}
        {\raggedright\popsemi\fontsize{8.4}{11}\selectfont\color{white}Tous les cours signés, Tuteur IA illimité, fiches PDF — onglet OnBuch+ de l'app\par}
      \end{minipage}
    \end{tcolorbox}
  \end{minipage}
  \vspace{8pt}
  \popcontenu
  \vfill
  \signatureonbuch
  \restoregeometry
  \newpage
}

% Rangée « Ce cours contient » (page de garde)
\newcommand{\poptile}[3]{% #1 couleur #2 gros texte #3 légende
  \begin{tcolorbox}[enhanced,colback=#1,colframe=popink,boxrule=2pt,arc=9pt,shadow={2.5pt}{-2.5pt}{0pt}{popink},
      left=6pt,right=6pt,top=9pt,bottom=9pt,halign=center,valign=center,height=1.75cm,width=\linewidth,nobeforeafter]
    {\popdisplay\fontsize{12.5}{13}\selectfont\color{white}\MakeUppercase{#2}}\par\vspace{3pt}
    {\centering\popsemi\fontsize{7.8}{9.5}\selectfont\color{white}#3\par}
  \end{tcolorbox}}
\newcommand{\popcontenu}{%
  \par\noindent{\popxboldls\fontsize{7.5}{7.5}\selectfont\color{popmuted}CE COURS SIGNÉ CONTIENT}\par\vspace{7pt}
  \noindent\begin{minipage}[t]{0.235\linewidth}\poptile{popblue}{Cours}{complet, illustré, pas à pas}\end{minipage}\hfill
  \begin{minipage}[t]{0.235\linewidth}\poptile{poporange}{Méthodes}{et exercices résolus}\end{minipage}\hfill
  \begin{minipage}[t]{0.235\linewidth}\poptile{poppink}{Exercices}{type examen + corrigés}\end{minipage}\hfill
  \begin{minipage}[t]{0.235\linewidth}\poptile{popgreen}{Fiche bilan}{l'essentiel à retenir}\end{minipage}\par}

% Sommaire
\newtcolorbox{sommaire}{enhanced,colback=white,colframe=popink,boxrule=2.2pt,arc=10pt,
  drop shadow=popink,left=18pt,right=18pt,top=13pt,bottom=13pt,before skip=6pt,after skip=20pt,
  before upper={\poppill{Sommaire}{poppurple}{white}\par\vspace{9pt}\popsemi\fontsize{10.6}{14.5}\selectfont\color{popink}}}

% Signature de fin de cours — poussée tout en bas de la dernière page
\newcommand{\findecours}{%
  \par\vfill\nopagebreak
  \begin{center}\popsquiggle{poporange}\end{center}\vspace{4pt}
  \signatureonbuch
}
\AtBeginDocument{\def\llbracket{\mathopen{[\mkern-3.5mu[}}\def\rrbracket{\mathclose{]\mkern-3.5mu]}}} % intervalles d'entiers (glyphe ⟦ absent de la police)
% Glyphes absents de Fira Math : redéfinis à partir de glyphes disponibles
\AtBeginDocument{%
  \def\setminus{\mathbin{\backslash}}\let\smallsetminus\setminus
  \def\vdots{\mathord{\vbox{\baselineskip3.2pt\lineskiplimit0pt\kern4pt\hbox{.}\hbox{.}\hbox{.}}}}%
  \def\ddots{\mathinner{\mkern1mu\raise6.5pt\hbox{.}\mkern2mu\raise3.6pt\hbox{.}\mkern2mu\raise0.7pt\hbox{.}\mkern1mu}}%
  \def\triangle{\mathord{\scalebox{0.9}{$\symup{\Delta}$}}}%
  \def\nabla{\mathord{\raisebox{\depth}{\scalebox{1}[-1]{$\symup{\Delta}$}}}}%
  \def\checkmark{\popcheck{popdark}}%
  \def\star{\mathbin{\ast}}\def\bigstar{\mathord{\ast}}%
  \def\diamond{\mathbin{\scalebox{0.6}{$\Diamond$}}}%
}

%% FIGFIX-BEGIN v5 — correctifs globaux des figures (docs/onbuch-generation/tools/figfix)
\usepackage{adjustbox}\usepackage{etoolbox}\tcbuselibrary{raster}
% toute figure TikZ/pgfplots plus large que la ligne est réduite à la largeur disponible
\BeforeBeginEnvironment{tikzpicture}{\begin{adjustbox}{max width=\linewidth}}
\AfterEndEnvironment{tikzpicture}{\end{adjustbox}}
% légendes pgfplots lisibles par-dessus les courbes
\pgfplotsset{every axis/.append style={legend style={fill=white,fill opacity=0.92,text opacity=1,draw=black!15,inner sep=3pt}}}
% étiquettes d'arêtes (syntaxe edge["..."]) avec fond blanc
\tikzset{every edge quotes/.append style={fill=white,inner sep=1.2pt}}
%% FIGFIX-END
\begin{document}

\pagegarde{Dérivées et sens de variation}{Mathématiques}{Tle A}{Relations et opérations fondamentales dans l'ensemble des nombres réels}{A03}{5 h}{Cette leçon introduit le calcul des dérivées des fonctions usuelles et composées, puis établit le lien fondamental entre le signe de la dérivée et le sens de variation d'une fonction. Elle aboutit à la construction de tableaux de variation complets et à la résolution de problèmes d'optimisation concrets.
\vspace{4pt}
\begin{multicols}{2}\small
\begin{itemize}
\item À la fin de cette leçon, je sais rappeler la définition du nombre dérivé et son interprétation graphique
\item calculer la dérivée de toutes les fonctions usuelles (puissances, racine, inverse, trigonométriques)
\item calculer la dérivée d'une somme, d'un produit, d'un quotient et d'une fonction composée
\item étudier le signe d'une dérivée sur un intervalle
\item déduire le sens de variation d'une fonction du signe de sa dérivée
\item dresser un tableau de variation complet (domaine, signe de f', variations, limites ou valeurs aux bornes)
\end{itemize}\end{multicols}}

\begin{sommaire}
\begin{multicols}{2}
\begin{enumerate}
\item Rappels : nombre dérivé et tangente
\item Dérivées des fonctions usuelles
\item Opérations sur les dérivées
\item Dérivée d'une fonction composée
\item Signe de la dérivée et sens de variation
\item Tableau de variation et optimisation
\item Activité d'intégration
\item Méthodes et exercices résolus
\item Exercices
\item Corrigés des exercices
\item Fiche bilan
\end{enumerate}
\end{multicols}
\end{sommaire}

\begin{objectifs}
\begin{itemize}
\item À la fin de cette leçon, je sais rappeler la définition du nombre dérivé et son interprétation graphique
\item calculer la dérivée de toutes les fonctions usuelles (puissances, racine, inverse, trigonométriques)
\item calculer la dérivée d'une somme, d'un produit, d'un quotient et d'une fonction composée
\item étudier le signe d'une dérivée sur un intervalle
\item déduire le sens de variation d'une fonction du signe de sa dérivée
\item dresser un tableau de variation complet (domaine, signe de f', variations, limites ou valeurs aux bornes)
\item repérer les extremums d'une fonction à partir de son tableau de variation
\item modéliser et résoudre un problème d'optimisation simple à l'aide de la dérivée
\end{itemize}
\end{objectifs}

\begin{prerequis}
\begin{itemize}
\item Nombre dérivé en un point et tangente à une courbe (Première)
\item Fonctions de référence : x ↦ xⁿ, x ↦ √x, x ↦ 1/x, fonctions trigonométriques
\item Signe d'un binôme, d'un trinôme et d'un produit/quotient (Seconde/Première)
\item Lecture graphique du sens de variation d'une fonction
\item Factorisation et résolution d'équations du second degré
\end{itemize}
\end{prerequis}

\newpage

\coursec{Rappels : nombre dérivé et tangente}

\textbf{Situation d'accroche.} Au marché Mokolo, Maman Ngo Bell vend ses beignets depuis des années. Elle a remarqué un phénomène troublant : lorsqu'elle augmente le prix de \SI{25}{FCFA}, elle vend certes moins de beignets, mais chaque beignet rapporte davantage. Alors elle se pose la question que tout commerçant se pose : \emph{à quel prix mon bénéfice quotidien sera-t-il maximal ?} Pour répondre, il faut comprendre comment une quantité (le bénéfice) \cle{varie} lorsqu'une autre quantité (le prix) \cle{change}. Le bénéfice est-il en train d'augmenter ou de diminuer lorsque le prix passe de 100 à 125 FCFA ? Et surtout, à quel moment cesse-t-il d'augmenter ? C'est exactement à ce type de questions que répond l'outil mathématique central de cette leçon : la \cle{dérivée}.

\textbf{Problématique.} Comment mesurer précisément la vitesse de variation d'une fonction en un point, et comment exploiter cette information pour décrire le comportement global de la fonction ?

Dans cette première section, nous revenons sur les fondements vus en Première : le nombre dérivé, son interprétation géométrique et la tangente à une courbe. Ces rappels sont indispensables : tout le reste de la leçon repose sur eux.

\courssub{Le taux d'accroissement : mesurer une variation moyenne}

Avant de parler de vitesse de variation \emph{en un point}, commençons par une idée plus simple : la variation \emph{moyenne} entre deux points.

Imaginons que tu relèves la distance parcourue par un taxi-brousse entre Yaoundé et Douala. Si le véhicule parcourt \SI{240}{km} en \SI{4}{h}, sa vitesse moyenne est de $\frac{240}{4} = \SI{60}{km/h}$. Cette vitesse moyenne ne dit rien de la vitesse instantanée à un moment précis : le chauffeur a pu rouler vite sur l'autoroute et lentement dans les embouteillages. De la même façon, pour une fonction $f$, la variation moyenne entre deux abscisses $a$ et $a + h$ se mesure par le \cle{taux d'accroissement}.

\begin{definition}[Taux d'accroissement]
Soit $f$ une fonction définie sur un intervalle $I$, $a$ un réel de $I$ et $h$ un réel non nul tel que $a + h \in I$. On appelle \cle{taux d'accroissement} (ou taux de variation) de $f$ entre $a$ et $a + h$ le réel :
\[ \tau(h) = \frac{f(a+h) - f(a)}{h}. \]
\end{definition}

Géométriquement, ce taux n'est rien d'autre que le \cle{coefficient directeur} de la droite $(AM)$, appelée \cle{sécante}, où $A\big(a \,;\, f(a)\big)$ et $M\big(a+h \,;\, f(a+h)\big)$ sont deux points de la courbe représentative de $f$. En effet, le coefficient directeur de la droite passant par deux points est le rapport de la différence des ordonnées sur la différence des abscisses.

\begin{exemplebox}[Un taux d'accroissement concret]
Soit $f(x) = x^2$. Le taux d'accroissement de $f$ entre $1$ et $1 + h$ est :
\[ \tau(h) = \frac{(1+h)^2 - 1^2}{h} = \frac{1 + 2h + h^2 - 1}{h} = \frac{2h + h^2}{h} = 2 + h. \]
Ainsi, entre $1$ et $1{,}1$ (c'est-à-dire $h = 0{,}1$), le taux vaut $2{,}1$ ; entre $1$ et $1{,}01$, il vaut $2{,}01$ ; entre $1$ et $1{,}001$, il vaut $2{,}001$. On devine que lorsque $h$ se rapproche de $0$, le taux se rapproche de $2$.
\end{exemplebox}

\courssub{Le nombre dérivé : une variation instantanée}

L'idée décisive, due à Newton et Leibniz au XVII\textsuperscript{e} siècle, consiste à faire tendre $h$ vers $0$ : le point $M$ glisse alors le long de la courbe vers le point $A$, et la sécante $(AM)$ pivote pour se rapprocher d'une position limite. Si le taux d'accroissement admet une limite finie, cette limite mesure la variation \emph{instantanée} de $f$ en $a$.

\begin{definition}[Nombre dérivé en un point]
Soit $f$ une fonction définie sur un intervalle $I$ et $a \in I$. On dit que $f$ est \cle{dérivable en $a$} si le taux d'accroissement $\dfrac{f(a+h) - f(a)}{h}$ admet une limite finie lorsque $h$ tend vers $0$. Cette limite, notée $f'(a)$, est appelée \cle{nombre dérivé} de $f$ en $a$ :
\[ f'(a) = \lim_{h \to 0} \frac{f(a+h) - f(a)}{h}. \]
\end{definition}

\begin{attention}[Bien comprendre la limite]
Dans la définition, on ne peut \textbf{jamais} remplacer directement $h$ par $0$ : on obtiendrait $\frac{0}{0}$, une forme indéterminée. Tout le travail consiste à \emph{simplifier} l'expression du taux (factoriser $h$, utiliser une identité remarquable, etc.) \emph{avant} de faire tendre $h$ vers $0$. C'est exactement ce que nous avons fait dans l'exemple précédent : $\tau(h) = 2 + h$ pour $h \neq 0$, et cette expression a une limite évidente quand $h \to 0$.
\end{attention}

\begin{exoresolu}[Calcul de $f'(2)$ pour $f(x) = x^2 - 3x$ par la définition]
Soit $f$ la fonction définie sur $\mathbb{R}$ par $f(x) = x^2 - 3x$. Montrer que $f$ est dérivable en $2$ et déterminer $f'(2)$.
\tcblower
\textbf{Solution détaillée.} On applique la définition avec $a = 2$.
\begin{align*}
f(2 + h) &= (2+h)^2 - 3(2+h) = 4 + 4h + h^2 - 6 - 3h = h^2 + h - 2, \\
f(2) &= 2^2 - 3 \times 2 = 4 - 6 = -2.
\end{align*}
Le taux d'accroissement est donc, pour $h \neq 0$ :
\[ \frac{f(2+h) - f(2)}{h} = \frac{h^2 + h - 2 - (-2)}{h} = \frac{h^2 + h}{h} = \frac{h(h+1)}{h} = h + 1. \]
Or $\displaystyle\lim_{h \to 0} (h + 1) = 1$, qui est un réel fini. Donc $f$ est dérivable en $2$ et :
\[ f'(2) = 1. \]
\textbf{Interprétation.} Au voisinage de $x = 2$, la fonction $f$ varie approximativement comme une fonction affine de coefficient $1$ : si $x$ augmente de $0{,}01$ à partir de $2$, alors $f(x)$ augmente d'environ $0{,}01$.
\end{exoresolu}

\courssub{Interprétation géométrique : la tangente à la courbe}

Revenons à notre sécante $(AM)$ qui pivote lorsque $M$ se rapproche de $A$. La position limite de cette sécante est une droite qui « épouse » la courbe au point $A$ : c'est la \cle{tangente}.

\begin{definition}[Tangente en un point]
Soit $f$ une fonction dérivable en $a$, de courbe représentative $\mathcal{C}_f$, et $A\big(a \,;\, f(a)\big)$. La \cle{tangente} à $\mathcal{C}_f$ au point $A$ est la droite passant par $A$ et de coefficient directeur $f'(a)$. Ainsi :
\[ \text{le nombre dérivé } f'(a) \text{ est le coefficient directeur de la tangente en } A. \]
\end{definition}

La figure suivante illustre ce phénomène pour $f(x) = x^2$ au point d'abscisse $1$ : les sécantes, de coefficients directeurs $3$, $2{,}5$ et $2{,}1$, se rapprochent de la tangente, de coefficient directeur $f'(1) = 2$.

\begin{popfigure}
\begin{center}
\begin{tikzpicture}
\begin{axis}[
    width=0.95\linewidth, height=7.5cm,
    axis lines=middle,
    xlabel={$x$}, ylabel={$y$},
    xmin=-0.6, xmax=2.8, ymin=-0.8, ymax=6.2,
    xtick={0,1,2}, ytick={1,2,4},
    legend style={at={(0.02,0.98)}, anchor=north west, font=\small},
    clip=false
]
\addplot[popcurve, domain=-0.5:2.5, samples=200]{x^2};
\addlegendentry{$\mathcal{C}_f : y = x^2$}
\addplot[poporange, thick, domain=-0.3:2.6, samples=2]{2*x - 1};
\addlegendentry{Tangente en $1$ : $y = 2x - 1$}
\addplot[popgreen, thick, dashed, domain=0.7:2.6, samples=2]{3*x - 2};
\addlegendentry{Sécante ($h=2$) : pente $3$}
\addplot[poppurple, thick, dashed, domain=0.5:2.6, samples=2]{2.5*x - 1.5};
\addlegendentry{Sécante ($h=1$) : pente $2{,}5$}
\addplot[popgold, thick, dashed, domain=0.3:2.6, samples=2]{2.1*x - 1.1};
\addlegendentry{Sécante ($h=0{,}1$) : pente $2{,}1$}
\addplot[only marks, mark=*, popdark] coordinates {(1,1)};
\node[popdark, above right, font=\small] at (axis cs:1,1) {$A(1\,;\,1)$};
\end{axis}
\end{tikzpicture}
\end{center}
\legende{Lorsque $h$ tend vers $0$, les sécantes passant par $A(1\,;\,1)$ pivotent et se rapprochent de la tangente, dont le coefficient directeur est $f'(1) = 2$.}
\end{popfigure}

\begin{propriete}[Équation de la tangente]
Soit $f$ une fonction dérivable en $a$. La tangente à la courbe $\mathcal{C}_f$ au point $A\big(a \,;\, f(a)\big)$ a pour équation :
\[ y = f'(a)\,(x - a) + f(a). \]
\end{propriete}

\begin{experience}[Démonstration guidée]
Démontrons cette formule, elle est formatrice et souvent demandée dans les raisonnements du Baccalauréat.
\begin{enumerate}
\item La tangente $T$ est une droite non verticale : elle admet une équation de la forme $y = mx + p$.
\item Par définition, son coefficient directeur est $m = f'(a)$. Donc $T : y = f'(a)\,x + p$.
\item La tangente passe par le point $A\big(a \,;\, f(a)\big)$ : ses coordonnées vérifient l'équation, donc $f(a) = f'(a) \times a + p$, d'où $p = f(a) - f'(a)\,a$.
\item En remplaçant : $y = f'(a)\,x + f(a) - f'(a)\,a = f'(a)(x - a) + f(a)$. \quad $\blacksquare$
\end{enumerate}
\end{experience}

\begin{exemplebox}[Équation d'une tangente]
Reprenons $f(x) = x^2 - 3x$ en $a = 2$. Nous avons calculé $f(2) = -2$ et $f'(2) = 1$. La tangente à $\mathcal{C}_f$ au point d'abscisse $2$ a donc pour équation :
\[ y = 1 \times (x - 2) + (-2) \quad \text{c'est-à-dire} \quad y = x - 4. \]
Vérification rapide : au point de contact, $x = 2$ donne $y = -2 = f(2)$. La droite passe bien par le point de tangence.
\end{exemplebox}

\begin{methode}[Déterminer l'équation d'une tangente en $a$]
\begin{enumerate}
\item Calculer $f(a)$ (ordonnée du point de contact).
\item Calculer $f'(a)$, soit par la définition, soit (comme nous le verrons dans les sections suivantes) à l'aide des formules de dérivation.
\item Écrire $y = f'(a)(x - a) + f(a)$, puis développer et réduire si nécessaire.
\item Vérifier mentalement que la droite obtenue passe bien par le point $\big(a \,;\, f(a)\big)$.
\end{enumerate}
\end{methode}

\begin{attention}[Pièges fréquents sur la tangente]
\begin{itemize}
\item \textbf{Confondre $f(a)$ et $f'(a)$} : $f(a)$ est une \emph{valeur de la fonction} (une ordonnée), $f'(a)$ est une \emph{pente} (un coefficient directeur). Ce sont deux nombres de nature différente.
\item \textbf{Oublier les parenthèses} : l'équation est $y = f'(a)(x - a) + f(a)$, et non $y = f'(a)\,x - a + f(a)$.
\item \textbf{Croire que la tangente ne coupe la courbe qu'en un seul point} : c'est faux en général ! La tangente à $y = x^2$ en $1$ ne recoupe pas la parabole, mais pour d'autres fonctions, la tangente peut traverser la courbe (pense à $x \mapsto x^3$ en $0$). La tangente est définie par une propriété \emph{locale} : elle approche la courbe au mieux \emph{au voisinage} de $A$.
\end{itemize}
\end{attention}

\courssub{De la dérivabilité en un point à la fonction dérivée}

Jusqu'ici, nous avons travaillé en un point $a$ \emph{fixé}. Mais rien n'empêche de faire varier $a$ : si une fonction est dérivable en \emph{tout} point d'un intervalle, chaque point possède son nombre dérivé, et on obtient une nouvelle fonction.

\begin{definition}[Fonction dérivée]
Soit $f$ une fonction définie sur un intervalle $I$. On dit que $f$ est \cle{dérivable sur $I$} si elle est dérivable en tout réel $a$ de $I$. Dans ce cas, la fonction qui à tout $x$ de $I$ associe le nombre dérivé $f'(x)$ est appelée \cle{fonction dérivée} de $f$ (ou simplement \emph{dérivée} de $f$) et se note $f'$ :
\[ f' : x \longmapsto f'(x). \]
\end{definition}

\begin{exemplebox}[La dérivée de la fonction carré, par la définition]
Soit $f(x) = x^2$. Pour un réel $a$ quelconque et $h \neq 0$ :
\[ \frac{f(a+h) - f(a)}{h} = \frac{(a+h)^2 - a^2}{h} = \frac{2ah + h^2}{h} = 2a + h \xrightarrow[h \to 0]{} 2a. \]
La fonction carré est donc dérivable sur $\mathbb{R}$ et, pour tout réel $x$ : $f'(x) = 2x$. On retrouve au passage que $f'(1) = 2$, cohérent avec notre figure. Dans la section suivante, nous établirons de telles formules pour toutes les fonctions usuelles : tu n'auras plus besoin de repasser par la définition à chaque fois.
\end{exemplebox}

\courssub{Dérivabilité et continuité}

Terminons par un lien fondamental entre deux notions que tu connais : la continuité et la dérivabilité. Lequel des deux concepts est le plus fort ?

\begin{propriete}[Dérivable implique continue]
Si une fonction $f$ est dérivable en un réel $a$, alors $f$ est continue en $a$.
\end{propriete}

\begin{experience}[Démonstration guidée]
Supposons $f$ dérivable en $a$. Pour $h \neq 0$, écrivons l'astuce suivante :
\[ f(a + h) - f(a) = \frac{f(a+h) - f(a)}{h} \times h. \]
Lorsque $h \to 0$ : le premier facteur tend vers $f'(a)$ (limite finie, par hypothèse de dérivabilité) et le second tend vers $0$. Le produit tend donc vers $f'(a) \times 0 = 0$. Ainsi $\displaystyle\lim_{h \to 0} f(a+h) = f(a)$, ce qui signifie exactement que $f$ est continue en $a$. \quad $\blacksquare$
\end{experience}

\begin{attention}[La réciproque est fausse !]
Une fonction \textbf{continue} en un point n'est \textbf{pas nécessairement dérivable} en ce point. Le contre-exemple classique — à connaître absolument pour le Bac — est la fonction \cle{valeur absolue} $f(x) = |x|$ en $0$.
\begin{itemize}
\item Elle est \textbf{continue en $0$} : $\displaystyle\lim_{x \to 0} |x| = 0 = f(0)$.
\item Mais son taux d'accroissement en $0$ est $\dfrac{|h| - 0}{h} = \dfrac{|h|}{h}$, qui vaut $1$ si $h > 0$ et $-1$ si $h < 0$. Les limites à droite ($1$) et à gauche ($-1$) sont différentes : le taux n'a \textbf{pas de limite} en $0$. Donc $f$ n'est \textbf{pas dérivable en $0$}.
\end{itemize}
Géométriquement, la courbe présente un \cle{point anguleux} en $0$ : il n'existe pas de tangente unique, mais deux « demi-tangentes » de pentes $-1$ et $1$. Retiens la hiérarchie : \textbf{dérivable $\Rightarrow$ continue}, mais \textbf{continue $\nRightarrow$ dérivable}.
\end{attention}

\begin{popfigure}
\begin{center}
\begin{tikzpicture}
\begin{axis}[
    width=0.85\linewidth, height=6.5cm,
    axis lines=middle,
    xlabel={$x$}, ylabel={$y$},
    xmin=-2.6, xmax=2.6, ymin=-0.7, ymax=2.6,
    xtick={-2,-1,1,2}, ytick={1,2},
    legend style={at={(0.02,0.98)}, anchor=north west, font=\small}
]
\addplot[popcurve, domain=-2.3:2.3, samples=200]{abs(x)};
\addlegendentry{$y = |x|$}
\addplot[poporange, thick, dashed, domain=0:2.3, samples=2]{x};
\addlegendentry{Demi-tangente à droite : pente $1$}
\addplot[popgreen, thick, dashed, domain=-2.3:0, samples=2]{-x};
\addlegendentry{Demi-tangente à gauche : pente $-1$}
\addplot[only marks, mark=*, popdark] coordinates {(0,0)};
\node[popdark, below right, font=\small] at (axis cs:0.05,-0.05) {Point anguleux};
\end{axis}
\end{tikzpicture}
\end{center}
\legende{La courbe de $x \mapsto |x|$ est continue en $0$ mais y présente un point anguleux : pas de tangente unique, donc pas de dérivabilité en $0$.}
\end{popfigure}

\begin{savaistu}[Pourquoi ce mot, « dérivée » ?]
Le terme vient du latin \emph{derivare}, « détourner un cours d'eau » : la fonction dérivée « découle » de la fonction initiale, comme un canal découle d'une rivière. Leibniz et Newton ont introduit ce concept indépendamment dans les années 1680 : Newton pour étudier les vitesses des planètes, Leibniz pour les tangentes aux courbes. Leur querelle de priorité a divisé les mathématiciens européens pendant des décennies ! Aujourd'hui, la dérivée est partout : c'est elle qui permet d'optimiser le réseau d'antennes d'Orange ou de MTN, de prévoir la propagation d'une épidémie, ou encore... de maximiser le bénéfice de Maman Ngo Bell, comme nous le ferons à la fin de cette leçon.
\end{savaistu}

\begin{aretenir}[L'essentiel de la section]
\begin{itemize}
\item Le \cle{taux d'accroissement} de $f$ entre $a$ et $a+h$ est $\dfrac{f(a+h)-f(a)}{h}$ : c'est la pente de la sécante.
\item $f$ est \cle{dérivable en $a$} si ce taux admet une limite finie quand $h \to 0$ ; cette limite est le nombre dérivé $f'(a)$.
\item $f'(a)$ est le \cle{coefficient directeur de la tangente} à la courbe au point d'abscisse $a$.
\item Équation de la tangente : $y = f'(a)(x - a) + f(a)$.
\item Si $f$ est dérivable en tout point d'un intervalle $I$, on définit la \cle{fonction dérivée} $f' : x \mapsto f'(x)$.
\item \cle{Dérivable en $a$} $\Rightarrow$ \cle{continue en $a$}, mais la réciproque est fausse (contre-exemple : $x \mapsto |x|$ en $0$).
\end{itemize}
\end{aretenir}

\begin{exercice}[À toi de jouer]
\begin{enumerate}
\item Soit $g(x) = 2x^2 + x$. En utilisant la définition, montrer que $g$ est dérivable en $3$ et calculer $g'(3)$.
\item Déterminer l'équation réduite de la tangente à la courbe de $g$ au point d'abscisse $3$.
\item La fonction $h(x) = |x - 1|$ est-elle dérivable en $1$ ? Justifier en étudiant le taux d'accroissement à gauche et à droite.
\end{enumerate}
\end{exercice}

\begin{corrige}[Exercice : À toi de jouer]
\textbf{1.} Pour $h \neq 0$ :
\[ \frac{g(3+h) - g(3)}{h} = \frac{2(3+h)^2 + (3+h) - 21}{h} = \frac{2h^2 + 13h}{h} = 2h + 13 \xrightarrow[h \to 0]{} 13. \]
La limite est finie, donc $g$ est dérivable en $3$ et $g'(3) = 13$.

\textbf{2.} On a $g(3) = 2 \times 9 + 3 = 21$ et $g'(3) = 13$. La tangente a pour équation $y = 13(x - 3) + 21$, soit $y = 13x - 18$.

\textbf{3.} Le taux d'accroissement de $h$ en $1$ est $\dfrac{|1 + h - 1| - 0}{h} = \dfrac{|h|}{h}$, qui vaut $1$ si $h > 0$ et $-1$ si $h < 0$. Les limites à gauche et à droite diffèrent, donc le taux n'a pas de limite : $h$ n'est pas dérivable en $1$ (elle y est pourtant continue : c'est un point anguleux, comme la valeur absolue en $0$).
\end{corrige}

Dans la section suivante, nous quitterons les calculs « à la main » par la définition pour établir les \textbf{dérivées des fonctions usuelles} : un formulaire puissant qui te fera gagner un temps précieux au Baccalauréat.

\coursec{Dérivées des fonctions usuelles}

Dans la section précédente, nous avons appris à calculer le \cle{nombre dérivé} d'une fonction en un point $a$ à l'aide du taux d'accroissement :
\[ f'(a) = \lim_{h \to 0} \frac{f(a+h) - f(a)}{h}. \]
Cette méthode est fondamentale, mais avouons-le : refaire le calcul de la limite à chaque fois que l'on rencontre une fonction serait aussi fastidieux que de refaire le trajet Douala--Yaoundé à pied à chaque voyage ! La bonne nouvelle, c'est que pour toutes les fonctions que l'on rencontre au lycée (puissances, inverse, racine carrée, sinus, cosinus), le calcul de la limite a été fait \emph{une fois pour toutes}. Il suffit alors d'apprendre les résultats dans un tableau et de savoir les utiliser. C'est l'objet de cette section : établir rigoureusement ce tableau, comprendre d'où viennent les formules, puis s'entraîner à les appliquer.

\courssub{Le tableau des dérivées usuelles}

\begin{propriete}[Dérivées des fonctions usuelles]
Les fonctions usuelles sont dérivables sur leur ensemble de dérivabilité, et l'on a :
\begin{center}
\begin{tabularx}{\linewidth}{|c|X|c|c|}
\hline
\rowcolor{popblueL} \textbf{Fonction $f(x)$} & \textbf{Dérivée $f'(x)$} & \textbf{Ensemble de définition} & \textbf{Ensemble de dérivabilité} \\
\hline
$k$ (constante) & $0$ & $\mathbb{R}$ & $\mathbb{R}$ \\
\hline
$x$ & $1$ & $\mathbb{R}$ & $\mathbb{R}$ \\
\hline
$x^n$ \ ($n \in \mathbb{N}^*$) & $n\,x^{n-1}$ & $\mathbb{R}$ & $\mathbb{R}$ \\
\hline
$\dfrac{1}{x}$ & $-\dfrac{1}{x^2}$ & $\mathbb{R}^*$ & $\mathbb{R}^*$ \\
\hline
$\sqrt{x}$ & $\dfrac{1}{2\sqrt{x}}$ & $[0\,;+\infty[$ & $]0\,;+\infty[$ \\
\hline
$\sin x$ & $\cos x$ & $\mathbb{R}$ & $\mathbb{R}$ \\
\hline
$\cos x$ & $-\sin x$ & $\mathbb{R}$ & $\mathbb{R}$ \\
\hline
\end{tabularx}
\end{center}
\end{propriete}

\begin{aretenir}[Les deux idées directrices]
\begin{itemize}
\item Dériver une puissance $x^n$ : on \cle{fait descendre l'exposant} devant et on le \cle{diminue de 1} : $(x^n)' = n x^{n-1}$.
\item Les fonctions $1/x$ et $\sqrt{x}$ ont des dérivées \emph{négatives} ou \emph{singulières} qu'il faut mémoriser à part : elles ne rentrent pas directement dans la formule des puissances entières.
\end{itemize}
\end{aretenir}

\courssub{Premières dérivées : constante et identité}

Commençons par les plus simples. Si $f(x) = k$ est constante, alors pour tout $h \neq 0$ :
\[ \frac{f(a+h)-f(a)}{h} = \frac{k-k}{h} = 0 \xrightarrow[h\to 0]{} 0. \]
Une fonction constante ne varie pas : sa pente est nulle partout. C'est cohérent avec l'image graphique : la courbe est une droite horizontale, sa tangente en tout point est elle-même.

Si $f(x) = x$, alors :
\[ \frac{f(a+h)-f(a)}{h} = \frac{(a+h)-a}{h} = \frac{h}{h} = 1 \xrightarrow[h\to 0]{} 1. \]
La droite d'équation $y = x$ a pour pente $1$ : sa dérivée vaut $1$ partout. Logique !

\courssub{Dérivée de la fonction puissance : $(x^n)' = n\,x^{n-1}$}

C'est le résultat le plus utilisé de tout le chapitre. Sa démonstration est formatrice et peut être demandée sous forme guidée au Baccalauréat.

\begin{propriete}[Théorème : dérivée de $x^n$]
Pour tout entier $n \in \mathbb{N}^*$, la fonction $f : x \mapsto x^n$ est dérivable sur $\mathbb{R}$ et :
\[ \boxed{\,f'(x) = n\,x^{n-1}\,} \]
\end{propriete}

\begin{experience}[Démonstration guidée par récurrence]
Nous admettons provisoirement (la démonstration rigoureuse sera faite dans la section « Opérations sur les dérivées ») la formule du produit : si $u$ et $v$ sont dérivables, alors $(uv)' = u'v + uv'$.

\textbf{Initialisation.} Pour $n = 1$ : $f(x) = x$, donc $f'(x) = 1 = 1 \times x^{0}$. La formule est vraie au rang $1$.

\textbf{Hérédité.} Supposons la formule vraie au rang $n$, c'est-à-dire $(x^n)' = n x^{n-1}$. Montrons-la au rang $n+1$. Écrivons :
\[ x^{n+1} = x^n \times x. \]
Posons $u(x) = x^n$ et $v(x) = x$. Par hypothèse de récurrence, $u'(x) = n x^{n-1}$, et $v'(x) = 1$. La formule du produit donne :
\begin{align*}
\left(x^{n+1}\right)' &= u'(x)\,v(x) + u(x)\,v'(x) \\
&= n x^{n-1} \times x + x^n \times 1 \\
&= n x^{n} + x^n = (n+1)\,x^n.
\end{align*}
La formule est donc vraie au rang $n+1$.

\textbf{Conclusion.} Par le principe de récurrence, pour tout $n \in \mathbb{N}^*$, $(x^n)' = n x^{n-1}$. $\blacksquare$
\end{experience}

\begin{exemplebox}[Application immédiate]
\begin{itemize}
\item $(x^2)' = 2x$ ; \quad $(x^3)' = 3x^2$ ; \quad $(x^4)' = 4x^3$ ;
\item $(x^{10})' = 10 x^9$ ; \quad $(x^{2025})' = 2025\,x^{2024}$.
\end{itemize}
Observe le mécanisme : l'exposant « descend » en facteur, puis on lui retire $1$.
\end{exemplebox}

\courssub{Dérivée de la fonction inverse : $(1/x)' = -1/x^2$}

\begin{propriete}[Théorème : dérivée de l'inverse]
La fonction $f : x \mapsto \dfrac{1}{x}$ est dérivable sur $\mathbb{R}^*$ et :
\[ f'(x) = -\frac{1}{x^2}. \]
\emph{Démonstration exigible : elle se fait par le taux d'accroissement.}
\end{propriete}

\begin{experience}[Démonstration guidée par le taux d'accroissement]
Soit $a \neq 0$ fixé. Pour $h$ assez petit (de sorte que $a + h \neq 0$) :
\begin{align*}
\frac{f(a+h)-f(a)}{h} &= \frac{1}{h}\left(\frac{1}{a+h} - \frac{1}{a}\right) \\
&= \frac{1}{h} \times \frac{a - (a+h)}{a(a+h)} \\
&= \frac{1}{h} \times \frac{-h}{a(a+h)} \\
&= \frac{-1}{a(a+h)}.
\end{align*}
Quand $h \to 0$, le dénominateur $a(a+h)$ tend vers $a^2$. Donc :
\[ f'(a) = \lim_{h \to 0} \frac{-1}{a(a+h)} = -\frac{1}{a^2}. \]
Ceci étant vrai pour tout $a \neq 0$, on a bien $f'(x) = -\dfrac{1}{x^2}$ sur $\mathbb{R}^*$. $\blacksquare$
\end{experience}

\begin{attention}[Le signe moins est essentiel]
L'erreur classique consiste à écrire $(1/x)' = 1/x^2$ en oubliant le signe moins. Retiens : la fonction inverse est \emph{décroissante} sur chacun des intervalles $]-\infty\,;0[$ et $]0\,;+\infty[$, donc sa dérivée doit être \emph{négative}. Le signe moins est une garantie de cohérence.
\end{attention}

\courssub{Dérivée de la fonction racine carrée}

\begin{propriete}[Théorème : dérivée de $\sqrt{x}$]
La fonction $f : x \mapsto \sqrt{x}$ est dérivable sur $]0\,;+\infty[$ et :
\[ f'(x) = \frac{1}{2\sqrt{x}}. \]
\end{propriete}

\begin{experience}[Démonstration guidée par le taux d'accroissement]
Soit $a > 0$. L'astuce consiste à multiplier par la \cle{quantité conjuguée} $\sqrt{a+h}+\sqrt{a}$ :
\begin{align*}
\frac{f(a+h)-f(a)}{h} &= \frac{\sqrt{a+h}-\sqrt{a}}{h} \\
&= \frac{(\sqrt{a+h}-\sqrt{a})(\sqrt{a+h}+\sqrt{a})}{h\left(\sqrt{a+h}+\sqrt{a}\right)} \\
&= \frac{(a+h)-a}{h\left(\sqrt{a+h}+\sqrt{a}\right)} \\
&= \frac{h}{h\left(\sqrt{a+h}+\sqrt{a}\right)} = \frac{1}{\sqrt{a+h}+\sqrt{a}}.
\end{align*}
Quand $h \to 0$, $\sqrt{a+h} \to \sqrt{a}$, donc :
\[ f'(a) = \frac{1}{\sqrt{a}+\sqrt{a}} = \frac{1}{2\sqrt{a}}. \quad \blacksquare \]
\end{experience}

\begin{attention}[La racine carrée n'est PAS dérivable en 0]
La fonction $\sqrt{x}$ est \emph{définie} en $0$, mais pas \emph{dérivable} en $0$. En effet, le taux d'accroissement en $0$ vaut :
\[ \frac{\sqrt{0+h}-\sqrt{0}}{h} = \frac{\sqrt{h}}{h} = \frac{1}{\sqrt{h}} \xrightarrow[h \to 0^+]{} +\infty. \]
La limite est infinie : ce n'est pas un nombre réel, donc $f'(0)$ n'existe pas. Graphiquement, la courbe de $\sqrt{x}$ admet en $0$ une \cle{demi-tangente verticale}. Dans un exercice, écris toujours : « $f$ est dérivable sur $]0\,;+\infty[$ », jamais sur $[0\,;+\infty[$.
\end{attention}

\courssub{Dérivées des fonctions sinus et cosinus}

\begin{propriete}[Théorème admis : dérivées de sin et cos]
Les fonctions sinus et cosinus sont dérivables sur $\mathbb{R}$ et :
\[ (\sin x)' = \cos x \qquad \text{et} \qquad (\cos x)' = -\sin x. \]
\emph{Ce résultat est admis en Terminale A ; la démonstration rigoureuse repose sur la limite $\lim_{h \to 0} \frac{\sin h}{h} = 1$, hors programme.}
\end{propriete}

\textbf{Justification graphique.} On peut se convaincre du résultat en observant les courbes. La dérivée en un point, c'est la \emph{pente de la tangente}. Regardons la courbe de $\sin x$ :
\begin{itemize}
\item En $x = 0$, la courbe de sinus monte avec une pente maximale égale à $1$ : or $\cos 0 = 1$. 
\item En $x = \dfrac{\pi}{2}$, la courbe de sinus atteint son sommet : la tangente est horizontale, pente nulle, et $\cos\dfrac{\pi}{2} = 0$. 
\item En $x = \pi$, la courbe descend avec pente $-1$, et $\cos \pi = -1$. 
\end{itemize}
À chaque instant, la pente de la courbe de $\sin x$ vaut exactement $\cos x$ : c'est bien le contenu de la formule $(\sin)' = \cos$.

\begin{popfigure}
\begin{center}
\begin{tikzpicture}
\begin{axis}[
    width=0.85\linewidth, height=6cm,
    axis lines=middle,
    xlabel={$x$}, ylabel={$y$},
    xtick={-3.1416,-1.5708,1.5708,3.1416},
    xticklabels={$-\pi$,$-\frac{\pi}{2}$,$\frac{\pi}{2}$,$\pi$},
    ytick={-1,1},
    domain=-3.6:3.6, samples=200,
    legend style={at={(0.5,-0.22)}, anchor=north, legend columns=2}
]
\addplot[popcurve, popblue, very thick] {sin(deg(x))};
\addlegendentry{$y = \sin x$}
\addplot[popcurve, poporange, very thick, dashed] {cos(deg(x))};
\addlegendentry{$y = \cos x = (\sin x)'$}
\end{axis}
\end{tikzpicture}
\end{center}
\legende{La courbe en pointillés orange ($\cos x$) donne en chaque point la pente de la courbe bleue ($\sin x$) : quand $\sin x$ est à son maximum, $\cos x$ vaut $0$ ; quand $\sin x$ croît le plus vite, $\cos x$ vaut $1$.}
\end{popfigure}

\begin{attention}[Attention au signe pour le cosinus]
$(\cos x)' = -\sin x$ : le signe moins est là aussi fréquemment oublié. Moyen mnémotechnique : « \emph{co}sinus se dérive en \emph{contraire} du sinus ». Vérifie graphiquement : en $x = 0$, la courbe de $\cos x$ est à son maximum (pente nulle) et $-\sin 0 = 0$ ; juste après $0$, $\cos x$ décroît (pente négative) et $-\sin x < 0$ pour $x > 0$ petit. C'est cohérent.
\end{attention}

\courssub{Complément utile : dérivée de $x^a$ pour $a \in \mathbb{Q}^*$}

\begin{propriete}[Complément (admis) : puissances rationnelles]
Pour tout rationnel non nul $a \in \mathbb{Q}^*$, la fonction $x \mapsto x^a$ est dérivable sur $]0\,;+\infty[$ et :
\[ \left(x^a\right)' = a\,x^{a-1}. \]
Ce résultat généralise la formule $(x^n)' = n x^{n-1}$ et unifie toutes les formules du tableau.
\end{propriete}

\begin{savaistu}[Une seule formule pour toutes les gouverner]
La formule $(x^a)' = a x^{a-1}$ contient toutes les autres ! En effet :
\begin{itemize}
\item $\dfrac{1}{x} = x^{-1}$, donc $\left(\dfrac{1}{x}\right)' = -1 \times x^{-2} = -\dfrac{1}{x^2}$ ;
\item $\sqrt{x} = x^{1/2}$, donc $(\sqrt{x})' = \dfrac{1}{2} x^{-1/2} = \dfrac{1}{2\sqrt{x}}$.
\end{itemize}
Au Baccalauréat, tu peux utiliser directement les formules du tableau ; cette généralisation est un excellent outil de vérification et un atout pour les fonctions comme $x^{3/2}$ ou $\dfrac{1}{x^3} = x^{-3}$.
\end{savaistu}

\courssub{Entraînement au calcul direct}

\begin{exoresolu}[Dériver trois fonctions usuelles]
\textbf{Énoncé.} Déterminer l'ensemble de dérivabilité et la fonction dérivée de chacune des fonctions suivantes :
\[ f(x) = x^5, \qquad g(x) = \frac{1}{x^2}, \qquad h(x) = 3\sqrt{x}. \]
\tcblower
\textbf{Solution détaillée.}

\textbf{1. Fonction $f(x) = x^5$.} C'est une fonction puissance, définie et dérivable sur $\mathbb{R}$. On applique $(x^n)' = n x^{n-1}$ avec $n = 5$ :
\[ f'(x) = 5\,x^{5-1} = 5x^4. \]

\textbf{2. Fonction $g(x) = \dfrac{1}{x^2}$.} Deux rédactions possibles.

\emph{Méthode 1 (complément puissances rationnelles).} On écrit $g(x) = x^{-2}$, définie sur $\mathbb{R}^*$ et dérivable sur $\mathbb{R}^*$. Alors :
\[ g'(x) = -2\,x^{-2-1} = -2x^{-3} = -\frac{2}{x^3}. \]

\emph{Méthode 2 (vérification par le taux d'accroissement en un point).} Vérifions en $a = 1$ que $g'(1) = -2$ :
\[ \frac{g(1+h)-g(1)}{h} = \frac{1}{h}\left(\frac{1}{(1+h)^2} - 1\right) = \frac{1-(1+h)^2}{h(1+h)^2} = \frac{-2h-h^2}{h(1+h)^2} = \frac{-2-h}{(1+h)^2} \xrightarrow[h\to 0]{} -2. \]
La formule est confirmée : $g'(1) = -2 = -\dfrac{2}{1^3}$. 

\textbf{3. Fonction $h(x) = 3\sqrt{x}$.} La fonction est définie sur $[0\,;+\infty[$ mais dérivable seulement sur $]0\,;+\infty[$ (la racine carrée n'est pas dérivable en $0$). Le coefficient $3$ se conserve par dérivation (nous le justifierons dans la section sur les opérations) :
\[ h'(x) = 3 \times \frac{1}{2\sqrt{x}} = \frac{3}{2\sqrt{x}}. \]
\end{exoresolu}

\begin{methode}[Réflexes pour dériver une fonction usuelle]
\begin{enumerate}
\item \textbf{Identifier} la forme de la fonction : puissance, inverse, racine, sinus, cosinus.
\item \textbf{Préciser l'ensemble de dérivabilité} avant tout calcul : $\mathbb{R}$ pour les puissances entières, sin et cos ; $\mathbb{R}^*$ pour l'inverse ; $]0\,;+\infty[$ pour la racine.
\item \textbf{Réécrire si besoin} sous forme de puissance : $\dfrac{1}{x^2} = x^{-2}$, $\sqrt{x} = x^{1/2}$.
\item \textbf{Appliquer la formule} en faisant descendre l'exposant et en le diminuant de $1$.
\item \textbf{Vérifier la cohérence du signe} : fonction croissante $\Rightarrow$ dérivée positive.
\end{enumerate}
\end{methode}

\begin{exemplebox}[Tableau synthétique à compléter en classe]
Recopie et complète ce tableau sans regarder le cours, puis vérifie :
\begin{center}
\begin{tabularx}{\linewidth}{|X|X|X|}
\hline
\rowcolor{popblueL} \textbf{Fonction} & \textbf{Dérivée} & \textbf{Dérivable sur} \\
\hline
$x^7$ & $\ldots$ & $\ldots$ \\
\hline
$\dfrac{1}{x}$ & $\ldots$ & $\ldots$ \\
\hline
$\sqrt{x}$ & $\ldots$ & $\ldots$ \\
\hline
$\sin x$ & $\ldots$ & $\ldots$ \\
\hline
$\cos x$ & $\ldots$ & $\ldots$ \\
\hline
$\dfrac{1}{x^3} = x^{-3}$ & $\ldots$ & $\ldots$ \\
\hline
\end{tabularx}
\end{center}
\emph{Réponses :} $7x^6$ sur $\mathbb{R}$ ; $-\dfrac{1}{x^2}$ sur $\mathbb{R}^*$ ; $\dfrac{1}{2\sqrt{x}}$ sur $]0\,;+\infty[$ ; $\cos x$ sur $\mathbb{R}$ ; $-\sin x$ sur $\mathbb{R}$ ; $-\dfrac{3}{x^4}$ sur $\mathbb{R}^*$.
\end{exemplebox}

\begin{savaistu}[Lagrange et la notation $f'$]
La notation $f'(x)$ que nous utilisons est due à \textbf{Joseph-Louis Lagrange} (1736--1813), mathématicien italien naturalisé français, l'un des plus grands esprits du XVIII\textsuperscript{ème} siècle. Avant lui, Leibniz notait la dérivée $\dfrac{\mathrm{d}f}{\mathrm{d}x}$ (notation toujours utilisée en physique) et Newton la notait $\dot{f}$ (encore employée pour les dérivées par rapport au temps). Lagrange est aussi l'inventeur du mot « dérivée » lui-même, et le théorème des accroissements finis, que tu rencontreras peut-être, porte son nom. Des lycées français portent souvent son nom... et ses idées irriguent tout ton programme de Terminale !
\end{savaistu}

\begin{exercice}[À toi de jouer]
Dériver chacune des fonctions suivantes en précisant l'ensemble de dérivabilité :
\[ a(x) = x^8 \ ; \quad b(x) = \frac{1}{x^5} \ ; \quad c(x) = 2\sqrt{x} \ ; \quad d(x) = \sin x + \cos x \ ; \quad e(x) = x^{3/2}. \]
\end{exercice}

\begin{corrige}[Exercice : dérivées usuelles]
\begin{itemize}
\item $a(x) = x^8$ : dérivable sur $\mathbb{R}$, $a'(x) = 8x^7$.
\item $b(x) = \dfrac{1}{x^5} = x^{-5}$ : dérivable sur $\mathbb{R}^*$, $b'(x) = -5x^{-6} = -\dfrac{5}{x^6}$.
\item $c(x) = 2\sqrt{x}$ : dérivable sur $]0\,;+\infty[$, $c'(x) = 2 \times \dfrac{1}{2\sqrt{x}} = \dfrac{1}{\sqrt{x}}$.
\item $d(x) = \sin x + \cos x$ : dérivable sur $\mathbb{R}$, $d'(x) = \cos x - \sin x$.
\item $e(x) = x^{3/2}$ : dérivable sur $]0\,;+\infty[$, $e'(x) = \dfrac{3}{2}x^{1/2} = \dfrac{3\sqrt{x}}{2}$.
\end{itemize}
\end{corrige}

Tu maîtrises désormais les dérivées des fonctions de base. Mais les fonctions rencontrées en pratique sont rarement « pures » : ce sont des sommes, des produits, des quotients ou des composées de fonctions usuelles. C'est précisément l'objet des deux sections suivantes : apprendre à dériver des combinaisons de fonctions.

\coursec{Opérations sur les dérivées}

Dans la section précédente, tu as appris à dériver les fonctions usuelles : $x^n$, $\sqrt{x}$, $\dfrac{1}{x}$, $\sin x$, $\cos x$, etc. C'est un bon début, mais les fonctions que tu rencontreras au Baccalauréat sont rarement aussi simples. Pense à une fonction comme $f(x) = (x^2+1)(3x-2)$ ou $g(x) = \dfrac{2x+1}{x-3}$ : aucune ligne du tableau des dérivées usuelles ne s'applique directement !

Heureusement, il existe des \cle{règles de dérivation} qui permettent de dériver une fonction « compliquée » à partir des dérivées de ses morceaux. L'idée est exactement celle d'un cuisinier : si tu sais préparer le riz et la sauce séparément, il te faut encore apprendre comment les combiner. C'est l'objet de cette section : apprendre à dériver les \cle{sommes}, les \cle{produits}, les \cle{inverses} et les \cle{quotients} de fonctions.

\courssub{Dérivée d'une somme et d'un produit par une constante}

\textbf{Intuition.} Imaginons deux taxis-moto qui partent de la gare routière de Bafoussam. Le premier avance à une vitesse de $u'(x)$ km/h, le second à $v'(x)$ km/h. Si l'on s'intéresse à la distance totale parcourue par les deux, il est naturel que la vitesse de croissance de cette somme soit la somme des vitesses. La dérivée, qui mesure une vitesse de variation, se comporte exactement ainsi.

\begin{propriete}[Dérivée d'une somme]
Soient $u$ et $v$ deux fonctions dérivables sur un intervalle $I$. Alors la fonction $u+v$ est dérivable sur $I$ et :
\[ (u+v)' = u' + v'. \]
\end{propriete}

\begin{propriete}[Dérivée d'un produit par une constante]
Soit $u$ une fonction dérivable sur un intervalle $I$ et $k$ un nombre réel. Alors la fonction $ku$ est dérivable sur $I$ et :
\[ (ku)' = k\,u'. \]
\end{propriete}

Ces deux règles se combinent : la dérivation est dite \cle{linéaire}. Ainsi, pour tous réels $a$ et $b$ :
\[ (au + bv)' = a\,u' + b\,v'. \]

\begin{exemplebox}[Dériver un polynôme terme à terme]
Soit $f(x) = 4x^3 - 5x^2 + 7x - 9$. On dérive terme par terme :
\begin{align*}
f'(x) &= 4 \times 3x^2 - 5 \times 2x + 7 \times 1 - 0 \\
&= 12x^2 - 10x + 7.
\end{align*}
Remarque que la constante $-9$ disparaît : une constante ne varie pas, sa dérivée est nulle.
\end{exemplebox}

\begin{attention}[La constante multiplicative reste, la constante additive disparaît]
Deux situations à ne pas confondre :
\begin{itemize}
\item Dans $3x^2$, le $3$ \emph{multiplie} : on le conserve, $(3x^2)' = 3 \times 2x = 6x$.
\item Dans $x^2 + 3$, le $3$ \emph{s'ajoute} : il disparaît, $(x^2+3)' = 2x$.
\end{itemize}
C'est l'erreur la plus fréquente en Première et en Terminale : relis-toi toujours sur ce point.
\end{attention}

\courssub{Dérivée d'un produit : la formule $(uv)' = u'v + uv'$}

\textbf{Intuition.} Un commerçant de Mokolo vend $u(x)$ sacs d'arachide par jour, à $v(x)$ FCFA le sac. Sa recette est le produit $u \times v$. Si le nombre de sacs augmente \emph{et} que le prix augmente, la recette augmente pour \emph{deux raisons} : on vend plus de sacs (au prix actuel), et chaque sac rapporte plus (pour la quantité actuelle). La variation totale de la recette est donc la somme de deux contributions : $u'v + uv'$. C'est exactement le contenu de la formule du produit.

\begin{propriete}[Dérivée d'un produit]
Soient $u$ et $v$ deux fonctions dérivables sur un intervalle $I$. Alors la fonction $uv$ est dérivable sur $I$ et :
\[ (uv)' = u'v + uv'. \]
\end{propriete}

Cette démonstration est formatrice et sa rédaction peut être exigée : suivons-la pas à pas.

\begin{experience}[Démonstration guidée de $(uv)' = u'v + uv'$]
On revient à la définition du nombre dérivé par le taux d'accroissement. Soit $a \in I$ et $h \neq 0$ tel que $a+h \in I$.

\textbf{Étape 1.} Écris le taux d'accroissement du produit :
\[ \frac{(uv)(a+h) - (uv)(a)}{h} = \frac{u(a+h)v(a+h) - u(a)v(a)}{h}. \]

\textbf{Étape 2.} L'astuce centrale : on ajoute et on retranche la même quantité $u(a)v(a+h)$ au numérateur (cela ne change rien, mais cela va faire apparaître les taux de $u$ et de $v$) :
\[ \frac{u(a+h)v(a+h) - u(a)v(a+h) + u(a)v(a+h) - u(a)v(a)}{h}. \]

\textbf{Étape 3.} Factorise par paires :
\[ \frac{u(a+h) - u(a)}{h} \times v(a+h) \;+\; u(a) \times \frac{v(a+h) - v(a)}{h}. \]

\textbf{Étape 4.} Fais tendre $h$ vers $0$ :
\begin{itemize}
\item $\dfrac{u(a+h)-u(a)}{h} \to u'(a)$ et $\dfrac{v(a+h)-v(a)}{h} \to v'(a)$ par définition ;
\item $v(a+h) \to v(a)$ car $v$, étant dérivable, est continue en $a$.
\end{itemize}

\textbf{Conclusion.} Le taux d'accroissement du produit tend vers $u'(a)v(a) + u(a)v'(a)$, ce qui prouve que $uv$ est dérivable en $a$ et que $(uv)'(a) = u'(a)v(a) + u(a)v'(a)$.
\end{experience}

\begin{attention}[La dérivée d'un produit n'est PAS le produit des dérivées]
C'est le piège numéro un de ce chapitre : $(uv)' \neq u'v'$. Vérifions-le sur un exemple simple. Prenons $u(x) = x$ et $v(x) = x$, donc $(uv)(x) = x^2$.
\begin{itemize}
\item Avec la vraie règle : $(x^2)' = 2x$.
\item Avec la fausse règle : $u'(x) \times v'(x) = 1 \times 1 = 1$.
\end{itemize}
Or $2x \neq 1$ ! La bonne formule donne bien $(x \cdot x)' = 1 \cdot x + x \cdot 1 = 2x$. Retiens : \textbf{on dérive chaque facteur à tour de rôle, en laissant l'autre intact, puis on additionne}.
\end{attention}

\begin{methode}[Identifier la structure avant de choisir la formule]
Avant de dériver, prends cinq secondes pour analyser la fonction :
\begin{enumerate}
\item \textbf{Repère l'opération principale} (celle qui se fait en dernier) : est-ce une somme, un produit, un quotient, une puissance ?
\item \textbf{Nomme les morceaux} : pose $u(x) = \ldots$ et $v(x) = \ldots$, et calcule $u'$ et $v'$ séparément.
\item \textbf{Applique la formule adaptée} en remplaçant méthodiquement.
\item \textbf{Simplifie et factorise} le résultat : une dérivée factorisée est indispensable pour étudier son signe ensuite.
\end{enumerate}
\end{methode}

\begin{exoresolu}[Dériver $f(x) = (x^2+1)(3x-2)$]
\textbf{Énoncé.} Soit $f$ la fonction définie sur $\mathbb{R}$ par $f(x) = (x^2+1)(3x-2)$. Calculer $f'(x)$ et le mettre sous forme développée.
\tcblower
\textbf{Solution.} La fonction $f$ est un \emph{produit} : on pose
\[ u(x) = x^2 + 1 \quad \text{et} \quad v(x) = 3x - 2. \]
On calcule séparément : $u'(x) = 2x$ et $v'(x) = 3$.

On applique la formule $(uv)' = u'v + uv'$ :
\begin{align*}
f'(x) &= u'(x)\,v(x) + u(x)\,v'(x) \\
&= 2x(3x-2) + (x^2+1) \times 3 \\
&= 6x^2 - 4x + 3x^2 + 3 \\
&= 9x^2 - 4x + 3.
\end{align*}
\textbf{Vérification (excellente habitude)} : développons d'abord $f(x) = 3x^3 - 2x^2 + 3x - 2$, puis dérivons terme à terme : $f'(x) = 9x^2 - 4x + 3$. On retrouve le même résultat, la formule du produit est confirmée.
\end{exoresolu}

\courssub{Dérivée de l'inverse et d'un quotient}

\textbf{Intuition.} Dire « inverse », c'est dire « partage ». Si $v(x)$ représente le nombre de bénéficiaires d'une aide et que chacun reçoit $\dfrac{1}{v(x)}$ milliard de FCFA, alors plus il y a de bénéficiaires, moins chacun reçoit : quand $v$ augmente, $\dfrac{1}{v}$ diminue. Il est donc logique que la dérivée de $\dfrac{1}{v}$ soit \emph{négative} quand $v' > 0$ : le signe moins de la formule n'est pas un hasard.

\begin{propriete}[Dérivée de l'inverse]
Soit $v$ une fonction dérivable sur un intervalle $I$, telle que $v(x) \neq 0$ pour tout $x \in I$. Alors la fonction $\dfrac{1}{v}$ est dérivable sur $I$ et :
\[ \left(\frac{1}{v}\right)' = -\frac{v'}{v^2}. \]
\end{propriete}

\begin{exemplebox}[Application directe]
Soit $h(x) = \dfrac{1}{x^2 + 1}$, définie sur $\mathbb{R}$ (le dénominateur ne s'annule jamais). Avec $v(x) = x^2+1$, on a $v'(x) = 2x$, donc :
\[ h'(x) = -\frac{2x}{(x^2+1)^2}. \]
\end{exemplebox}

\begin{propriete}[Dérivée d'un quotient]
Soient $u$ et $v$ deux fonctions dérivables sur un intervalle $I$, avec $v(x) \neq 0$ pour tout $x \in I$. Alors la fonction $\dfrac{u}{v}$ est dérivable sur $I$ et :
\[ \left(\frac{u}{v}\right)' = \frac{u'v - uv'}{v^2}. \]
\end{propriete}

Cette formule se \emph{déduit} des deux précédentes, ce qui est un excellent exercice de logique : un quotient n'est rien d'autre qu'un produit par un inverse. Écrivons $\dfrac{u}{v} = u \times \dfrac{1}{v}$ et appliquons la formule du produit :
\begin{align*}
\left(u \times \frac{1}{v}\right)' &= u' \times \frac{1}{v} + u \times \left(-\frac{v'}{v^2}\right) \\
&= \frac{u'}{v} - \frac{uv'}{v^2} \\
&= \frac{u'v}{v^2} - \frac{uv'}{v^2} \quad \text{(on réduit au même dénominateur)} \\
&= \frac{u'v - uv'}{v^2}.
\end{align*}

\begin{attention}[Ordre des termes et condition d'existence]
Deux vigilances avec le quotient :
\begin{itemize}
\item Le numérateur est $u'v - uv'$ : on commence \textbf{toujours} par dériver le \emph{numérateur}. Inverser en $uv' - u'v$ change le signe de toute la dérivée, et donc tout le tableau de variations qui suivra !
\item La formule n'a de sens que si $v(x) \neq 0$ : avant tout calcul, précise l'ensemble de définition. Pour $g(x) = \dfrac{2x+1}{x-3}$, il faut exclure $x = 3$.
\end{itemize}
\end{attention}

\begin{exoresolu}[Dériver $g(x) = \dfrac{2x+1}{x-3}$]
\textbf{Énoncé.} Soit $g$ la fonction définie sur $\mathbb{R} \setminus \{3\}$ par $g(x) = \dfrac{2x+1}{x-3}$. Calculer $g'(x)$ et en déduire le signe de $g'$.
\tcblower
\textbf{Solution.} La fonction $g$ est un \emph{quotient} : on pose
\[ u(x) = 2x + 1 \quad \text{et} \quad v(x) = x - 3, \]
d'où $u'(x) = 2$ et $v'(x) = 1$.

On applique $\left(\dfrac{u}{v}\right)' = \dfrac{u'v - uv'}{v^2}$ :
\begin{align*}
g'(x) &= \frac{2(x-3) - (2x+1) \times 1}{(x-3)^2} \\
&= \frac{2x - 6 - 2x - 1}{(x-3)^2} \\
&= \frac{-7}{(x-3)^2}.
\end{align*}
\textbf{Signe :} pour tout $x \neq 3$, le dénominateur $(x-3)^2$ est strictement positif et le numérateur $-7$ est strictement négatif. Donc $g'(x) < 0$ sur chacun des intervalles $]-\infty\,;\,3[$ et $]3\,;\,+\infty[$ : la fonction $g$ est strictement décroissante sur chacun de ces intervalles. (Attention à ne pas dire « décroissante sur $\mathbb{R}\setminus\{3\}$ » : la décroissance s'énonce intervalle par intervalle.)
\end{exoresolu}

\courssub{Application aux polynômes et aux fractions rationnelles}

Les règles précédentes couvrent deux grandes familles de fonctions du programme de Terminale A :

\begin{itemize}
\item Les \cle{polynômes} : sommes de termes de la forme $ax^n$. On dérive terme à terme avec $(x^n)' = nx^{n-1}$. Un polynôme est dérivable sur $\mathbb{R}$ tout entier.
\item Les \cle{fractions rationnelles} : quotients de deux polynômes. On les dérive avec la formule du quotient, sur leur ensemble de définition (on exclut les valeurs qui annulent le dénominateur). La dérivée d'une fraction rationnelle est encore une fraction rationnelle.
\end{itemize}

\begin{exemplebox}[Un polynôme de degré 3]
Soit $p(x) = 2x^3 - \dfrac{3}{2}x^2 + x - 5$. Alors :
\[ p'(x) = 6x^2 - 3x + 1. \]
Ce trinôme a pour discriminant $\Delta = 9 - 24 = -15 < 0$ : il est du signe de $a = 6$, donc $p'(x) > 0$ sur $\mathbb{R}$. Le polynôme $p$ est strictement croissant sur $\mathbb{R}$ — voilà un avant-goût de la section 5 !
\end{exemplebox}

Pour choisir rapidement la bonne formule face à une fonction donnée, l'arbre de décision suivant te servira de réflexe :

\begin{popfigure}
\begin{center}
\begin{tikzpicture}[
  grow=right,
  level distance=4.2cm,
  level 1/.style={sibling distance=3.2cm},
  level 2/.style={sibling distance=1.6cm},
  every node/.style={font=\small},
  racine/.style={rectangle, rounded corners, draw=popdark, fill=popblueL, thick, align=center, inner sep=5pt},
  branche/.style={rectangle, rounded corners, draw=popblue, fill=popcream, align=center, inner sep=4pt},
  formule/.style={rectangle, draw=poporange, fill=poporangeL, align=center, inner sep=4pt},
  edge from parent/.style={draw=popmuted, thick, -{Stealth}}
]
\node[racine, align=center]{Quelle est l'opération\\principale de $f$ ?}
  child { node[branche, align=center]{Quotient\\$f = \dfrac{u}{v}$}
    child { node[formule, align=center]{$f' = \dfrac{u'v - uv'}{v^2}$\\si $u$ \text{non constante}} }
    child { node[formule, align=center]{$f' = -\dfrac{k\,v'}{v^2}$\\si $f = \dfrac{k}{v}$} }
  }
  child { node[branche, align=center]{Produit\\$f = u \times v$}
    child { node[formule, align=center]{$f' = u'v + uv'$\\si $u$ \text{non constante}} }
    child { node[formule, align=center]{$f' = k\,u'$\\si $f = k\,u$} }
  }
  child { node[branche, align=center]{Somme\\$f = u + v$}
    child { node[formule, align=center]{$f' = u' + v'$\\(\text{terme à terme})} }
  };
\end{tikzpicture}
\end{center}
\legende{Arbre de décision : identifier d'abord la structure de la fonction, puis choisir la formule de dérivation adaptée.}
\end{popfigure}

\begin{aretenir}[Les cinq formules à connaître par cœur]
Pour $u$ et $v$ dérivables sur $I$ (et $v \neq 0$ sur $I$ pour les deux dernières) :
\[ (u+v)' = u' + v' \qquad (ku)' = k\,u' \qquad (uv)' = u'v + uv' \]
\[ \left(\frac{1}{v}\right)' = -\frac{v'}{v^2} \qquad \left(\frac{u}{v}\right)' = \frac{u'v - uv'}{v^2} \]
Et le réflexe qui sauve des points : \textbf{la dérivée d'un produit n'est pas le produit des dérivées}.
\end{aretenir}

\begin{savaistu}[Pourquoi ces formules portent-elles le nom de Leibniz ?]
C'est Gottfried Wilhelm Leibniz (1646--1716), contemporain et rival de Newton, qui le premier formula proprement la règle du produit, avec la notation $\mathrm{d}(uv) = u\,\mathrm{d}v + v\,\mathrm{d}u$ que les physiciens utilisent encore. Une anecdote amusante : Leibniz confessa avoir d'abord cru, comme beaucoup d'élèves, que la dérivée d'un produit était le produit des dérivées, avant de se corriger en testant sa règle sur $x \cdot x = x^2$ — exactement la vérification que nous avons faite plus haut ! Même les plus grands mathématiciens sont passés par l'erreur que tu dois maintenant éviter.
\end{savaistu}

\begin{exercice}[À toi de jouer]
Dériver les fonctions suivantes en précisant leur ensemble de dérivabilité :
\begin{enumerate}
\item $f(x) = (3x - 5)(x^2 + 4x)$ ;
\item $g(x) = \dfrac{4}{x^2 + 2}$ ;
\item $h(x) = \dfrac{x^2 - x + 1}{x + 2}$.
\end{enumerate}
\end{exercice}

\begin{corrige}[Exercice]
\textbf{1.} Produit avec $u(x) = 3x-5$, $v(x) = x^2+4x$, donc $u'(x) = 3$ et $v'(x) = 2x+4$ :
\begin{align*}
f'(x) &= 3(x^2+4x) + (3x-5)(2x+4) \\
&= 3x^2 + 12x + 6x^2 + 12x - 10x - 20 \\
&= 9x^2 + 14x - 20, \quad \text{dérivable sur } \mathbb{R}.
\end{align*}
\textbf{2.} Inverse avec $v(x) = x^2+2$ (jamais nul), $v'(x) = 2x$ :
\[ g'(x) = -\frac{4 \times 2x}{(x^2+2)^2} = -\frac{8x}{(x^2+2)^2}, \quad \text{dérivable sur } \mathbb{R}. \]
\textbf{3.} Quotient avec $u(x) = x^2 - x + 1$, $v(x) = x+2$, $u'(x) = 2x-1$, $v'(x) = 1$ :
\begin{align*}
h'(x) &= \frac{(2x-1)(x+2) - (x^2 - x + 1)}{(x+2)^2} \\
&= \frac{2x^2 + 4x - x - 2 - x^2 + x - 1}{(x+2)^2} \\
&= \frac{x^2 + 4x - 3}{(x+2)^2}, \quad \text{dérivable sur } \mathbb{R} \setminus \{-2\}.
\end{align*}
\end{corrige}

Tu maîtrises désormais toutes les règles opératoires de la dérivation. Il reste un dernier outil de calcul — la dérivée d'une fonction composée, objet de la section suivante — avant de passer au cœur du chapitre : utiliser le signe de la dérivée pour déterminer le sens de variation d'une fonction et dresser son tableau de variations.

\coursec{Dérivée d'une fonction composée}

Dans la section précédente, tu as appris à dériver des sommes, des produits et des quotients de fonctions. Mais il reste un obstacle de taille : comment dériver une fonction comme $f(x) = (3x^2 - 5)^4$ ? On pourrait certes développer, mais c'est long et périlleux. Et que faire de $g(x) = \sqrt{x^2 + 1}$ ou de $h(x) = \cos\left(2x - \dfrac{\pi}{3}\right)$, où aucun développement n'est possible ? Ces fonctions ont un point commun : elles sont obtenues en \emph{enchaînant} deux fonctions, l'une après l'autre, comme deux machines montées en série dans un atelier. C'est la notion de \cle{fonction composée}, et il existe une règle de dérivation magnifique pour elle : la \cle{règle de dérivation en chaîne}. C'est l'un des outils les plus utilisés du programme de Terminale, au Baccalauréat comme dans la vie scientifique.

\courssub{Notion de fonction composée : la chaîne de fonctions}

\textbf{Une intuition concrète.}\par
Imagine un transporteur de Douala qui expédie du cacao vers Yaoundé. Le prix payé au planteur dépend du poids $x$ de la cargaison : c'est une première fonction $u$. Puis le bénéfice final du transporteur dépend du prix payé : c'est une seconde fonction $v$. Le bénéfice en fonction du poids, c'est donc $v$ appliquée au résultat de $u$ : on a \emph{composé} les deux fonctions.

\begin{definition}[Fonction composée]
Soit $u$ une fonction définie sur un intervalle $I$, à valeurs dans un intervalle $J$, et $v$ une fonction définie sur $J$. La \cle{fonction composée} de $u$ par $v$, notée $v \circ u$ (lire « $v$ rond $u$ »), est la fonction définie sur $I$ par :
\[ (v \circ u)(x) = v\big(u(x)\big). \]
On applique d'abord $u$, puis $v$ au résultat.
\end{definition}

\begin{attention}[L'ordre compte !]
Dans $v \circ u$, c'est $u$ qui agit \textbf{en premier}, même si elle s'écrit à droite. En général $v \circ u \neq u \circ v$. Par exemple, si $u(x) = x^2$ et $v(x) = x + 1$, alors $(v \circ u)(x) = x^2 + 1$ mais $(u \circ v)(x) = (x+1)^2$.
\end{attention}

\begin{exemplebox}[Lire une chaîne de fonctions]
Décomposons $f(x) = (3x^2 - 5)^4$. On pose $u(x) = 3x^2 - 5$ (la fonction « intérieure ») et $v(t) = t^4$ (la fonction « extérieure »). Alors :
\[ v\big(u(x)\big) = \big(3x^2 - 5\big)^4 = f(x), \quad \text{donc } f = v \circ u. \]
La chaîne est : $x \xrightarrow{\ u\ } 3x^2 - 5 \xrightarrow{\ v\ } (3x^2 - 5)^4$.
\end{exemplebox}

\begin{popfigure}
\begin{center}
\begin{tikzpicture}[>=Stealth, node distance=2.6cm]
\node[draw=popblue, fill=popblueL, rounded corners, minimum width=1.6cm, minimum height=0.9cm] (x) {$x$};
\node[draw=poporange, fill=poporangeL, rounded corners, minimum width=2.2cm, minimum height=0.9cm, right=of x] (u) {$u(x)$};
\node[draw=poppurple, fill=poppurpleL, rounded corners, minimum width=2.6cm, minimum height=0.9cm, right=of u] (v) {$v\big(u(x)\big)$};
\draw[->, very thick, popdark] (x) -- node[above, font=\small]{$u$} node[below, font=\small, popmuted]{dérivée $u'$} (u);
\draw[->, very thick, popdark] (u) -- node[above, font=\small]{$v$} node[below, font=\small, popmuted]{dérivée $v'$} (v);
\node[below=0.7cm of u, font=\small\itshape, popmuted] {la dérivée globale est le \textbf{produit} $u' \times (v' \circ u)$};
\end{tikzpicture}
\end{center}
\legende{La chaîne de composition : à chaque étape, la fonction se dérive ; la dérivée de la composée est le produit des dérivées le long de la chaîne.}
\end{popfigure}

\courssub{Théorème de dérivation des fonctions composées}

\textbf{Pourquoi un produit ? L'idée heuristique.}\par
Reprenons notre transporteur. Si le poids augmente de $1$ kg, le prix payé augmente de $u'$ francs. Et si le prix payé augmente de $1$ franc, le bénéfice varie de $v'$ francs. Donc si le poids augmente de $1$ kg, le bénéfice varie de $u' \times v'$ francs : les taux de variation se \emph{multiplient} le long de la chaîne. C'est exactement ce que dit le théorème.

\begin{experience}[Justification heuristique par les taux d'accroissement]
Soit $a$ un point où $u$ est dérivable, et supposons $v$ dérivable en $u(a)$. Pour $h \neq 0$ « petit », écrivons le taux d'accroissement de $v \circ u$ entre $a$ et $a + h$ :
\[ \frac{(v \circ u)(a+h) - (v \circ u)(a)}{h} = \frac{v\big(u(a+h)\big) - v\big(u(a)\big)}{u(a+h) - u(a)} \times \frac{u(a+h) - u(a)}{h}. \]
Quand $h$ tend vers $0$ :
\begin{enumerate}
\item le second facteur tend vers $u'(a)$, par définition du nombre dérivé ;
\item comme $u$ est continue en $a$ (toute fonction dérivable est continue), $u(a+h)$ tend vers $u(a)$, donc le premier facteur est un taux d'accroissement de $v$ qui tend vers $v'\big(u(a)\big)$.
\end{enumerate}
Le produit tend donc vers $u'(a) \times v'\big(u(a)\big)$. (Cette justification n'est pas une démonstration complète — il faudrait traiter le cas où $u(a+h) = u(a)$ — mais elle en donne l'idée essentielle. Le théorème est \textbf{admis} au programme.)
\end{experience}

\begin{propriete}[Théorème de dérivation des fonctions composées (admis)]
Soit $u$ une fonction dérivable sur un intervalle $I$, à valeurs dans $J$, et $v$ une fonction dérivable sur $J$. Alors $v \circ u$ est dérivable sur $I$ et :
\[ (v \circ u)' = u' \times (v' \circ u), \qquad \text{c'est-à-dire} \qquad (v \circ u)'(x) = u'(x) \times v'\big(u(x)\big). \]
On dérive la fonction \emph{extérieure} $v$ en laissant l'intérieur $u(x)$ inchangé, \textbf{puis on multiplie par la dérivée de l'intérieur} $u'$.
\end{propriete}

\begin{attention}[Le piège numéro 1 : oublier le facteur $u'$]
L'erreur la plus fréquente au Baccalauréat est d'écrire $(v \circ u)' = v' \circ u$, en oubliant de multiplier par $u'$. Retiens cette phrase : \emph{« on dérive l'extérieur, on garde l'intérieur, et on n'oublie jamais de multiplier par la dérivée de l'intérieur »}. Vérifie systématiquement : si $u$ n'est pas une simple fonction $x \mapsto x$, le facteur $u'$ doit apparaître.
\end{attention}

\courssub{Cas particuliers essentiels}

En appliquant le théorème avec les fonctions extérieures usuelles $v(t) = t^n$, $v(t) = \dfrac{1}{t}$ et $v(t) = \sqrt{t}$, on obtient les trois formules à connaître par cœur.

\begin{propriete}[Formules de dérivation des composées usuelles]
Soit $u$ une fonction dérivable sur un intervalle $I$ et $n$ un entier naturel non nul.
\begin{itemize}
\item \textbf{Puissance :} $\left(u^n\right)' = n\, u'\, u^{\,n-1}$.
\item \textbf{Inverse :} si $u$ ne s'annule pas sur $I$, $\left(\dfrac{1}{u}\right)' = -\dfrac{u'}{u^2}$.
\item \textbf{Racine carrée :} si $u(x) > 0$ sur $I$, $\left(\sqrt{u}\right)' = \dfrac{u'}{2\sqrt{u}}$.
\end{itemize}
\end{propriete}

\begin{exemplebox}[D'où viennent ces formules ?]
Pour la racine carrée, on applique le théorème avec $v(t) = \sqrt{t}$, dont la dérivée est $v'(t) = \dfrac{1}{2\sqrt{t}}$. Alors :
\[ \left(\sqrt{u}\right)'(x) = u'(x) \times v'\big(u(x)\big) = u'(x) \times \frac{1}{2\sqrt{u(x)}} = \frac{u'(x)}{2\sqrt{u(x)}}. \]
Les deux autres formules s'obtiennent de la même manière avec $v(t) = t^n$ et $v(t) = \dfrac{1}{t}$.
\end{exemplebox}

\textbf{Fonctions du type $x \mapsto f(ax + b)$.}\par
Un cas particulier extrêmement fréquent en physique (signaux, ondes) et au Baccalauréat : la fonction intérieure est \emph{affine}, $u(x) = ax + b$, de dérivée $u'(x) = a$.

\begin{propriete}[Dérivée de $x \mapsto f(ax+b)$]
Si $f$ est dérivable sur un intervalle contenant les valeurs $ax + b$, alors la fonction $x \mapsto f(ax + b)$ est dérivable et :
\[ \big(f(ax + b)\big)' = a\, f'(ax + b). \]
En particulier, pour $a$ et $b$ réels :
\[ \big(\sin(ax + b)\big)' = a\cos(ax + b) \qquad \text{et} \qquad \big(\cos(ax + b)\big)' = -a\sin(ax + b). \]
\end{propriete}

\begin{aretenir}[Réflexe à acquérir]
Quand tu dérives une fonction de $ax + b$, le facteur $a$ « sort » devant. Par exemple : $\big(\cos(3x)\big)' = -3\sin(3x)$ et $\big(\sin(2x - 1)\big)' = 2\cos(2x - 1)$. Ce facteur $a$ est le facteur $u'$ de la règle de la chaîne.
\end{aretenir}

\courssub{Méthode : décomposer en chaîne avant de dériver}

\begin{methode}[Dériver une fonction composée pas à pas]
Pour dériver une fonction $f$ qui n'est ni une somme, ni un produit simple :
\begin{enumerate}
\item \textbf{Identifier la chaîne} : repérer la fonction intérieure $u$ (ce qui est « à l'intérieur » de la parenthèse, de la racine, du cosinus...) et la fonction extérieure $v$.
\item \textbf{Écrire la décomposition} : $f = v \circ u$, avec $u(x) = \ldots$ et $v(t) = \ldots$.
\item \textbf{Dériver chaque maillon} : calculer $u'(x)$ et $v'(t)$ séparément.
\item \textbf{Assembler} : $f'(x) = u'(x) \times v'\big(u(x)\big)$, en remplaçant $t$ par $u(x)$ dans $v'$.
\item \textbf{Simplifier} et vérifier les conditions d'existence (dénominateur non nul, quantité sous la racine strictement positive).
\end{enumerate}
Avec l'entraînement, les étapes 1 et 2 se font mentalement, mais au début, écris-les toujours : c'est la meilleure garantie contre l'oubli du facteur $u'$.
\end{methode}

\begin{exoresolu}[Trois dérivations de fonctions composées]
Dériver les fonctions suivantes en précisant leur ensemble de dérivabilité :
\[ f(x) = (3x^2 - 5)^4, \qquad g(x) = \sqrt{x^2 + 1}, \qquad h(x) = \cos\left(2x - \frac{\pi}{3}\right). \]
\tcblower
\textbf{1. Dérivée de $f(x) = (3x^2 - 5)^4$.}

Décomposition : $u(x) = 3x^2 - 5$ et $v(t) = t^4$, donc $f = v \circ u$, dérivable sur $\mathbb{R}$ comme composée de fonctions dérivables sur $\mathbb{R}$.

Dérivées des maillons : $u'(x) = 6x$ et $v'(t) = 4t^3$.

Assemblage avec la formule $\left(u^n\right)' = n u' u^{n-1}$ :
\[ f'(x) = 4 \times u'(x) \times \big(u(x)\big)^3 = 4 \times 6x \times (3x^2 - 5)^3 = 24x\,(3x^2 - 5)^3. \]

\textbf{2. Dérivée de $g(x) = \sqrt{x^2 + 1}$.}

Décomposition : $u(x) = x^2 + 1$ et $v(t) = \sqrt{t}$. Comme $x^2 + 1 > 0$ pour tout réel $x$, $g$ est dérivable sur $\mathbb{R}$ tout entier.

Dérivées des maillons : $u'(x) = 2x$ et $v'(t) = \dfrac{1}{2\sqrt{t}}$.

Assemblage avec la formule $\left(\sqrt{u}\right)' = \dfrac{u'}{2\sqrt{u}}$ :
\[ g'(x) = \frac{2x}{2\sqrt{x^2 + 1}} = \frac{x}{\sqrt{x^2 + 1}}. \]

\textbf{3. Dérivée de $h(x) = \cos\left(2x - \dfrac{\pi}{3}\right)$.}

C'est une fonction de la forme $\cos(ax + b)$ avec $a = 2$ et $b = -\dfrac{\pi}{3}$, dérivable sur $\mathbb{R}$. La fonction intérieure $u(x) = 2x - \dfrac{\pi}{3}$ a pour dérivée $u'(x) = 2$. Donc :
\[ h'(x) = -2\sin\left(2x - \frac{\pi}{3}\right). \]
Remarque le facteur $2$ qui « sort » : c'est la dérivée de l'intérieur.
\end{exoresolu}

\courssub{Applications : polynômes, fractions rationnelles, fonctions irrationnelles}

La règle de la chaîne transforme radicalement notre capacité de calcul. Voyons-la à l'œuvre sur les trois grandes familles de fonctions du programme.

\begin{exoresolu}[Une fraction rationnelle]
Soit $f$ définie par $f(x) = \dfrac{1}{x^2 - 4x + 5}$. Montrer que $f$ est dérivable sur $\mathbb{R}$ et calculer $f'(x)$.
\tcblower
\textbf{Ensemble de dérivabilité.} Le dénominateur est un trinôme de discriminant $\Delta = (-4)^2 - 4 \times 1 \times 5 = 16 - 20 = -4 < 0$. Comme le coefficient de $x^2$ est positif, $x^2 - 4x + 5 > 0$ pour tout réel $x$ : le dénominateur ne s'annule jamais, donc $f$ est définie et dérivable sur $\mathbb{R}$.

\textbf{Décomposition.} On pose $u(x) = x^2 - 4x + 5$ et $v(t) = \dfrac{1}{t}$ : alors $f = v \circ u$.

\textbf{Calcul.} $u'(x) = 2x - 4$. Avec la formule $\left(\dfrac{1}{u}\right)' = -\dfrac{u'}{u^2}$ :
\[ f'(x) = -\frac{2x - 4}{(x^2 - 4x + 5)^2} = \frac{4 - 2x}{(x^2 - 4x + 5)^2}. \]
\textbf{Bonus.} Le dénominateur étant strictement positif, $f'(x)$ est du signe de $4 - 2x$ : $f$ est croissante sur $\left]-\infty\,;\,2\right]$ et décroissante sur $\left[2\,;\,+\infty\right[$. Nous exploiterons pleinement ce genre de conclusion dans la section suivante.
\end{exoresolu}

\begin{exemplebox}[Vérification graphique : $f(x) = (2x+1)^2$ et sa dérivée]
Avec $u(x) = 2x + 1$ et $v(t) = t^2$, la formule donne $f'(x) = 2 \times 2 \times (2x+1) = 4(2x+1) = 8x + 4$. La figure ci-dessous permet de \emph{vérifier} ce calcul : la dérivée s'annule en $x = -\dfrac{1}{2}$, exactement là où la courbe de $f$ admet une tangente horizontale (son minimum) ; elle est négative avant (la courbe descend) et positive après (la courbe monte).
\end{exemplebox}

\begin{popfigure}
\begin{center}
\begin{tikzpicture}
\begin{axis}[
  width=0.85\linewidth, height=6.5cm,
  axis lines=middle, xlabel={$x$}, ylabel={$y$},
  xmin=-2.2, xmax=1.6, ymin=-2, ymax=9,
  grid=both, grid style={popline!40},
  legend style={at={(0.02,0.98)}, anchor=north west, font=\small},
  samples=200
]
\addplot[popcurve, domain=-2.1:1.5]{(2*x+1)^2};
\addlegendentry{$f(x) = (2x+1)^2$}
\addplot[popcurve, poporange, domain=-1.1:0.55]{8*x+4};
\addlegendentry{$f'(x) = 8x + 4$}
\addplot[mark=*, popdark, mark size=1.5pt] coordinates {(-0.5,0)};
\node[font=\small, popdark] at (axis cs:-0.5,-1.1) {$-\frac{1}{2}$};
\end{axis}
\end{tikzpicture}
\end{center}
\legende{La dérivée $f'(x) = 8x + 4$ s'annule en $x = -\tfrac{1}{2}$, abscisse du minimum de $f$ : le signe de $f'$ donne le sens de variation de $f$.}
\end{popfigure}

\begin{attention}[Erreurs fréquentes à bannir]
\begin{itemize}
\item Écrire $\big((3x^2-5)^4\big)' = 4(3x^2-5)^3$ : il manque le facteur $u'(x) = 6x$.
\item Écrire $\left(\sqrt{u}\right)' = \dfrac{1}{2\sqrt{u}}$ : il manque $u'$ au numérateur.
\item Dériver $\cos(2x)$ en $-\sin(2x)$ : le facteur $a = 2$ a été oublié ; la bonne réponse est $-2\sin(2x)$.
\item Oublier les conditions d'existence : $\sqrt{u}$ n'est dérivable que là où $u > 0$ (et pas seulement $u \geq 0$, car la racine carrée n'est pas dérivable en $0$).
\end{itemize}
\end{attention}

\begin{savaistu}[La règle de la chaîne dans ta poche]
Les réseaux mobiles comme MTN ou Orange modulent des signaux de la forme $s(t) = A\sin(\omega t + \varphi)$, où $\omega$ est la pulsation. Pour étudier la vitesse de variation du signal, les ingénieurs dérivent : $s'(t) = A\omega\cos(\omega t + \varphi)$ — une application directe de la formule de cette leçon, avec le facteur $\omega$ qui « sort ». Historiquement, c'est Gottfried Wilhelm Leibniz qui, au XVII\textsuperscript{e} siècle, a formulé cette règle avec sa notation $\dfrac{\mathrm{d}y}{\mathrm{d}x} = \dfrac{\mathrm{d}y}{\mathrm{d}u} \times \dfrac{\mathrm{d}u}{\mathrm{d}x}$, où l'on « voit » littéralement les taux se multiplier le long de la chaîne.
\end{savaistu}

\begin{exercice}[À toi de jouer]
\begin{enumerate}
\item Décomposer en chaîne puis dériver : $f(x) = (5x - 2)^3$ ; $g(x) = \dfrac{1}{3x + 2}$ sur $\left]-\dfrac{2}{3}\,;\,+\infty\right[$ ; $k(x) = \sqrt{4x - 3}$.
\item Calculer la dérivée de $s(x) = \sin\left(3x + \dfrac{\pi}{4}\right)$ et de $c(x) = \cos(5x)$.
\item Soit $f(x) = (x^2 + x)^5$. Calculer $f'(x)$, puis déterminer l'équation de la tangente à la courbe de $f$ au point d'abscisse $0$.
\end{enumerate}
\end{exercice}

\begin{corrige}[Exercice : dérivées de fonctions composées]
\textbf{1.} $f = v \circ u$ avec $u(x) = 5x - 2$, $v(t) = t^3$ : $f'(x) = 3 \times 5 \times (5x-2)^2 = 15(5x-2)^2$.

$g = v \circ u$ avec $u(x) = 3x + 2$ (qui ne s'annule pas sur l'intervalle), $v(t) = \dfrac{1}{t}$ : $g'(x) = -\dfrac{3}{(3x+2)^2}$.

$k = v \circ u$ avec $u(x) = 4x - 3$, $v(t) = \sqrt{t}$ ; $k$ est dérivable là où $4x - 3 > 0$, c'est-à-dire sur $\left]\dfrac{3}{4}\,;\,+\infty\right[$ : $k'(x) = \dfrac{4}{2\sqrt{4x-3}} = \dfrac{2}{\sqrt{4x-3}}$.

\textbf{2.} $s'(x) = 3\cos\left(3x + \dfrac{\pi}{4}\right)$ et $c'(x) = -5\sin(5x)$.

\textbf{3.} Avec $u(x) = x^2 + x$, $u'(x) = 2x + 1$ : $f'(x) = 5(2x+1)(x^2+x)^4$. En $x = 0$ : $f(0) = 0$ et $f'(0) = 5 \times 1 \times 0 = 0$. La tangente au point d'abscisse $0$ a pour équation $y = f'(0)(x - 0) + f(0)$, soit $y = 0$ : c'est l'axe des abscisses.
\end{corrige}

Tu maîtrises désormais l'outil de dérivation le plus puissant du programme. Dans la prochaine section, nous mettrons ces dérivées au service d'une question fondamentale : comment le \emph{signe} de la dérivée révèle-t-il le \emph{sens de variation} de la fonction ? C'est le pont entre le calcul et la lecture des courbes — et la clé des problèmes d'optimisation du Baccalauréat.

\coursec{Signe de la dérivée et sens de variation}

Dans la section précédente, tu as appris à calculer la dérivée de n'importe quelle fonction usuelle ou composée. Mais à quoi sert ce calcul ? Voici la réponse, et c'est l'une des plus belles idées de tout le programme de Terminale : \cle{le signe de la dérivée révèle le sens de variation de la fonction}. Autrement dit, un simple calcul algébrique (la dérivée) permet de répondre à une question géométrique (la courbe monte-t-elle ou descend-elle ?). C'est ce pont entre l'algèbre et la géométrie que nous allons construire maintenant, et c'est l'outil central de tous les problèmes d'optimisation du Baccalauréat.

\courssub{Intuition : lire la pente des tangentes}

Reprenons l'image de la route. Quand tu parcours la courbe d'une fonction de gauche à droite, la tangente en chaque point indique la « pente » de la route à cet instant :

\begin{itemize}
\item si la tangente \emph{monte} (coefficient directeur positif), la fonction est en train de \emph{croître} ;
\item si la tangente \emph{descend} (coefficient directeur négatif), la fonction est en train de \emph{décroître} ;
\item si la tangente est \emph{horizontale} (coefficient directeur nul), la fonction est à un moment de « palier » : sommet de colline ou creux de vallée, le plus souvent.
\end{itemize}

Or le coefficient directeur de la tangente en un point d'abscisse $a$ est précisément $f'(a)$. Le signe de $f'(a)$ nous dit donc, localement, si la courbe monte ou descend. Observons cela sur un exemple concret.

\begin{popfigure}
\begin{center}
\begin{tikzpicture}
\begin{axis}[
  width=0.85\linewidth, height=7cm,
  axis lines=middle,
  xlabel={$x$}, ylabel={$y$},
  xmin=-0.5, xmax=4.5, ymin=-2, ymax=4.5,
  xtick={0,1,2,3,4}, ytick={-1,1,2,3,4},
  grid=both, grid style={popline!40},
  legend style={at={(0.02,0.98)}, anchor=north west, font=\small}
]
\addplot[popcurve,domain=-0.3:4.3,samples=200]{x^2-4*x+3};
\addlegendentry{$f(x)=x^2-4x+3$}
\addplot[poporange,thick,domain=0:1.2,samples=2]{-4*x+3};
\addlegendentry{tangente en $a=0$ : pente $-4$}
\addplot[popgreen,thick,domain=1.2:2.8,samples=2]{-1};
\addlegendentry{tangente en $a=2$ : pente $0$}
\addplot[poppurple,thick,domain=2.6:4.2,samples=2]{3*x-9.25};
\addlegendentry{tangente en $a=3{,}5$ : pente $3$}
\addplot[only marks,mark=*,popdark] coordinates {(0,3) (2,-1) (3.5,1.25)};
\end{axis}
\end{tikzpicture}
\end{center}
\legende{La parabole $f(x)=x^2-4x+3$, de dérivée $f'(x)=2x-4$. À gauche de $x=2$, les tangentes descendent ($f'<0$, $f$ décroît) ; au sommet, la tangente est horizontale ($f'(2)=0$) ; à droite, les tangentes montent ($f'>0$, $f$ croît). Vérifie : en $a=3{,}5$, $f(3{,}5)=1{,}25$ et $f'(3{,}5)=3$, d'où la tangente $y=3x-9{,}25$.}
\end{popfigure}

Sur cette figure, tout se lit : avant le sommet, la pente des tangentes est négative et la courbe descend ; après le sommet, elle est positive et la courbe monte ; exactement au sommet, la tangente est horizontale. Généralisons cette observation.

\courssub{Le théorème fondamental (admis)}

\begin{propriete}[Lien entre le signe de la dérivée et le sens de variation]
Soit $f$ une fonction \cle{dérivable} sur un intervalle $I$.
\begin{enumerate}
\item Si $f'(x) > 0$ pour tout $x \in I$ (sauf éventuellement en des points isolés où $f'$ s'annule), alors $f$ est \cle{strictement croissante} sur $I$.
\item Si $f'(x) < 0$ pour tout $x \in I$ (sauf éventuellement en des points isolés où $f'$ s'annule), alors $f$ est \cle{strictement décroissante} sur $I$.
\item Si $f'(x) = 0$ pour tout $x \in I$, alors $f$ est \cle{constante} sur $I$.
\end{enumerate}
Plus largement : si $f' \geq 0$ sur $I$, $f$ est croissante (au sens large) sur $I$ ; si $f' \leq 0$ sur $I$, $f$ est décroissante (au sens large) sur $I$.
\end{propriete}

\begin{attention}[Théorème admis]
Ce théorème est \textbf{admis} au programme de Terminale : sa démonstration rigoureuse repose sur le \emph{théorème des accroissements finis}, que tu rencontreras dans l'enseignement supérieur. Tu dois donc savoir l'\textbf{énoncer précisément} et l'\textbf{appliquer}, mais on ne te demandera pas de le démontrer au Baccalauréat.
\end{attention}

\begin{experience}[Justification graphique du théorème]
Pourquoi ce théorème est-il crédible ? Suivons le raisonnement pas à pas.
\begin{enumerate}
\item Fixons un point $a$ de $I$ où $f'(a) > 0$. La tangente en $a$ a un coefficient directeur strictement positif : elle « monte ».
\item Au voisinage de $a$, la courbe de $f$ est très proche de sa tangente (c'est exactement le sens de l'approximation affine $f(a+h) \approx f(a) + h\,f'(a)$). Donc la courbe monte aussi, localement.
\item Si $f'(x) > 0$ en \emph{tout} point de $I$, la courbe monte localement partout : elle ne peut que monter globalement sur tout l'intervalle. D'où la stricte croissance de $f$ sur $I$.
\item Le raisonnement est identique avec $f' < 0$ (la courbe descend partout) et avec $f' = 0$ (tangentes horizontales partout : la courbe ne monte ni ne descend, elle reste à altitude constante).
\end{enumerate}
La subtilité « points isolés où $f'$ s'annule » signifie qu'une tangente horizontale en un point \emph{isolé} ne casse pas la croissance : c'est le cas de $x \mapsto x^3$ en $0$, que nous étudierons plus bas.
\end{experience}

\begin{exemplebox}[Application immédiate]
Soit $f(x) = x^2 - 4x + 3$ sur $\mathbb{R}$. On calcule $f'(x) = 2x - 4 = 2(x-2)$.
\begin{itemize}
\item Si $x < 2$, alors $f'(x) < 0$ : $f$ est strictement décroissante sur $\left]-\infty\,; 2\right]$.
\item Si $x > 2$, alors $f'(x) > 0$ : $f$ est strictement croissante sur $\left[2\,; +\infty\right[$.
\end{itemize}
On retrouve exactement ce que montrait la figure : le sommet de la parabole est atteint en $x = 2$, et $f(2) = 4 - 8 + 3 = -1$ est le minimum de $f$.
\end{exemplebox}

\courssub{Extremum local et annulation de la dérivée}

\begin{definition}[Extremum local]
Soit $f$ une fonction définie sur un intervalle $I$ et $c$ un point \emph{intérieur} à $I$.
\begin{itemize}
\item $f$ admet un \cle{maximum local} en $c$ s'il existe un intervalle ouvert $J$ contenant $c$ tel que, pour tout $x \in J \cap I$, $f(x) \leq f(c)$.
\item $f$ admet un \cle{minimum local} en $c$ s'il existe un intervalle ouvert $J$ contenant $c$ tel que, pour tout $x \in J \cap I$, $f(x) \geq f(c)$.
\end{itemize}
Un extremum local est un maximum local ou un minimum local.
\end{definition}

Intuitivement, un maximum local est un « sommet de colline » : la valeur $f(c)$ dépasse toutes les valeurs voisines. Au sommet d'une colline, la route est momentanément plate : la tangente y est horizontale. C'est l'objet du théorème suivant, dont la démonstration est formatrice et peut être exigée.

\begin{propriete}[Condition nécessaire d'extremum local]
Soit $f$ une fonction définie sur un intervalle $I$, et $c$ un point \textbf{intérieur} à $I$. Si $f$ admet un extremum local en $c$ \textbf{et} si $f$ est dérivable en $c$, alors :
\[ f'(c) = 0. \]
\end{propriete}

\begin{experience}[Démonstration guidée]
Démontrons le théorème dans le cas d'un \emph{maximum local} en $c$ (le cas du minimum est analogue).
\begin{enumerate}
\item Puisque $f$ admet un maximum local en $c$, il existe $h_0 > 0$ tel que, pour tout $h$ avec $|h| < h_0$ :
\[ f(c+h) \leq f(c), \quad \text{donc} \quad f(c+h) - f(c) \leq 0. \]
\item Étudions le taux d'accroissement $\tau(h) = \dfrac{f(c+h)-f(c)}{h}$ en séparant les deux côtés.
\begin{itemize}
\item Pour $0 < h < h_0$ : le numérateur est $\leq 0$ et le dénominateur $> 0$, donc $\tau(h) \leq 0$. En passant à la limite quand $h \to 0^+$ : $f'(c) \leq 0$.
\item Pour $-h_0 < h < 0$ : le numérateur est $\leq 0$ et le dénominateur $< 0$, donc $\tau(h) \geq 0$. En passant à la limite quand $h \to 0^-$ : $f'(c) \geq 0$.
\end{itemize}
\item La fonction $f$ étant dérivable en $c$, les deux limites sont égales à $f'(c)$. On a donc simultanément $f'(c) \leq 0$ et $f'(c) \geq 0$, d'où :
\[ f'(c) = 0. \qquad \blacksquare \]
\end{enumerate}
\end{experience}

\begin{attention}[La réciproque est fausse !]
$f'(c) = 0$ est une condition \textbf{nécessaire} mais \textbf{pas suffisante}. Autrement dit : une tangente horizontale ne garantit pas un extremum. Contre-exemple classique : $f(x) = x^3$ en $c = 0$. On a $f'(x) = 3x^2$, donc $f'(0) = 0$, la tangente en $0$ est bien horizontale... et pourtant $f$ n'a \emph{ni maximum ni minimum} en $0$, car $f$ est strictement croissante sur $\mathbb{R}$ (pour $x < 0$, $f(x) < 0 = f(0)$, et pour $x > 0$, $f(x) > f(0)$). Le point $(0\,;0)$ est un \emph{point d'inflexion à tangente horizontale}. Moralité : pour conclure à un extremum, il faut vérifier que $f'$ \textbf{change de signe} en $c$.
\end{attention}

\begin{popfigure}
\begin{center}
\begin{tikzpicture}
\begin{axis}[
  width=0.8\linewidth, height=6.5cm,
  axis lines=middle,
  xlabel={$x$}, ylabel={$y$},
  xmin=-2.2, xmax=2.2, ymin=-4.5, ymax=4.5,
  xtick={-2,-1,1,2}, ytick={-4,-2,2,4},
  grid=both, grid style={popline!40},
  legend style={at={(0.02,0.98)}, anchor=north west, font=\small}
]
\addplot[popcurve,domain=-1.65:1.65,samples=200]{x^3};
\addlegendentry{$f(x)=x^3$}
\addplot[poporange,thick,domain=-1.5:1.5,samples=2]{0};
\addlegendentry{tangente horizontale en $0$}
\addplot[only marks,mark=*,popdark] coordinates {(0,0)};
\node[popdark,font=\small,anchor=west] at (axis cs:0.15,-1.6) {pas d'extremum !};
\end{axis}
\end{tikzpicture}
\end{center}
\legende{La cubique $f(x)=x^3$ : $f'(0)=0$ et la tangente en $0$ est horizontale, pourtant $f$ reste strictement croissante. La dérivée $f'(x)=3x^2$ est positive de part et d'autre de $0$ : pas de changement de signe, pas d'extremum.}
\end{popfigure}

\begin{propriete}[Extremum et changement de signe de la dérivée]
Soit $f$ dérivable sur un intervalle ouvert contenant $c$. Si $f'$ \cle{s'annule en $c$ en changeant de signe}, alors $f$ admet un extremum local en $c$ :
\begin{itemize}
\item si $f'$ passe du \textbf{positif au négatif} en $c$, alors $f(c)$ est un \textbf{maximum local} (la courbe monte puis descend) ;
\item si $f'$ passe du \textbf{négatif au positif} en $c$, alors $f(c)$ est un \textbf{minimum local} (la courbe descend puis monte).
\end{itemize}
Si $f'$ s'annule en $c$ \emph{sans} changer de signe (cas de $x^3$ en $0$), il n'y a pas d'extremum en $c$.
\end{propriete}

\courssub{Méthode complète d'étude des variations}

\begin{methode}[Les 4 étapes de l'étude des variations d'une fonction]
Pour étudier les variations d'une fonction $f$ sur un intervalle $I$ :
\begin{enumerate}
\item \textbf{Dériver} : calculer $f'(x)$ en utilisant les dérivées usuelles, les opérations (somme, produit, quotient) et, si besoin, la dérivée d'une composée.
\item \textbf{Factoriser} : écrire $f'(x)$ sous forme de produit (ou quotient) de facteurs du premier degré ou de signe constant. C'est l'étape-clé : une dérivée factorisée se discute facilement.
\item \textbf{Étudier le signe} : résoudre $f'(x) = 0$, puis dresser le \emph{tableau de signes} de $f'$ sur $I$.
\item \textbf{Conclure} : là où $f' > 0$, $f$ est strictement croissante ; là où $f' < 0$, $f$ est strictement décroissante. Repérer les points où $f'$ s'annule en changeant de signe : ce sont les extrema locaux. Résumer le tout dans le tableau de variation.
\end{enumerate}
\end{methode}

\begin{exoresolu}[Variations de $f(x) = x^3 - 3x$ sur $\mathbb{R}$]
\textbf{Énoncé.} Soit $f$ la fonction définie sur $\mathbb{R}$ par $f(x) = x^3 - 3x$. Étudier les variations de $f$ et déterminer ses extrema locaux.
\tcblower
\textbf{Solution détaillée.} Suivons la méthode en 4 étapes.

\textbf{Étape 1 : dériver.} $f$ est un polynôme, dérivable sur $\mathbb{R}$, et :
\[ f'(x) = 3x^2 - 3. \]

\textbf{Étape 2 : factoriser.}
\[ f'(x) = 3(x^2 - 1) = 3(x-1)(x+1). \]

\textbf{Étape 3 : signe de $f'$.} Les racines sont $x = -1$ et $x = 1$. Le facteur $3$ est strictement positif ; $(x-1)(x+1)$ est un trinôme du second degré, positif à l'extérieur des racines et négatif entre elles. D'où le tableau de signes :

\begin{center}
\begin{tabular}{|c|ccccccc|}
\hline
$x$ & $-\infty$ & & $-1$ & & $1$ & & $+\infty$ \\
\hline
$f'(x)$ & & $+$ & $0$ & $-$ & $0$ & $+$ & \\
\hline
\end{tabular}
\end{center}

\textbf{Étape 4 : conclusion.} Calculons les valeurs remarquables :
\[ f(-1) = (-1)^3 - 3(-1) = -1 + 3 = 2 \quad ; \quad f(1) = 1 - 3 = -2. \]
\begin{itemize}
\item $f$ est strictement croissante sur $\left]-\infty\,; -1\right]$ ;
\item $f$ est strictement décroissante sur $\left[-1\,; 1\right]$ ;
\item $f$ est strictement croissante sur $\left[1\,; +\infty\right[$.
\end{itemize}
En $x = -1$, $f'$ passe du positif au négatif : $f$ admet un \textbf{maximum local} égal à $2$. En $x = 1$, $f'$ passe du négatif au positif : $f$ admet un \textbf{minimum local} égal à $-2$.
\end{exoresolu}

\begin{attention}[Les trois pièges classiques de la rédaction]
\begin{enumerate}
\item \textbf{Conclure sur le signe de $f$ au lieu de $f'$.} Le tableau de signes porte sur $f'$, le tableau de variations sur $f$. Ne mélange jamais les deux lignes !
\item \textbf{Oublier la factorisation.} Étudier le signe de $f'(x) = 3x^2 - 3$ « directement » est périlleux ; factorise systématiquement.
\item \textbf{Annoncer un extremum dès que $f'(c)=0$.} Vérifie toujours le \emph{changement de signe} de $f'$ en $c$ (souviens-toi de $x^3$ en $0$).
\end{enumerate}
\end{attention}

\begin{savaistu}[Pourquoi ce théorème fait tourner le monde]
L'optimisation par la dérivée est partout autour de toi. Quand un opérateur comme MTN ou Orange dimensionne une antenne pour couvrir un quartier de Yaoundé au moindre coût, quand un transporteur de la ligne Douala--Bafoussam choisit sa vitesse pour minimiser la consommation de carburant, quand un planteur de la plaine des Mbo veut clôturer un champ rectangulaire d'aire maximale avec une longueur de grillage fixée : chaque fois, on modélise par une fonction, on dérive, on cherche où la dérivée s'annule en changeant de signe. La section suivante te donnera tous les outils pour résoudre toi-même ces problèmes, grâce au tableau de variation complet.
\end{savaistu}

\begin{exercice}[À toi de jouer]
Soit $g$ la fonction définie sur $\mathbb{R}$ par $g(x) = 2x^3 + 3x^2 - 12x + 1$.
\begin{enumerate}
\item Calculer $g'(x)$ et le factoriser.
\item Dresser le tableau de signes de $g'$.
\item En déduire les variations de $g$ et préciser ses extrema locaux (valeurs comprises).
\end{enumerate}
\end{exercice}

\begin{corrige}[Exercice : variations de $g$]
1. $g'(x) = 6x^2 + 6x - 12 = 6(x^2 + x - 2) = 6(x-1)(x+2)$.

2. Les racines sont $x = -2$ et $x = 1$. Le trinôme est positif à l'extérieur des racines :

\begin{center}
\begin{tabular}{|c|ccccccc|}
\hline
$x$ & $-\infty$ & & $-2$ & & $1$ & & $+\infty$ \\
\hline
$g'(x)$ & & $+$ & $0$ & $-$ & $0$ & $+$ & \\
\hline
\end{tabular}
\end{center}

3. $g$ est strictement croissante sur $\left]-\infty\,; -2\right]$, strictement décroissante sur $\left[-2\,; 1\right]$, strictement croissante sur $\left[1\,; +\infty\right[$. Comme $g'$ change de signe en $-2$ (de $+$ vers $-$), $g$ admet un maximum local : $g(-2) = 2\times(-8) + 3\times 4 - 12\times(-2) + 1 = -16 + 12 + 24 + 1 = 21$. En $1$ (de $-$ vers $+$), $g$ admet un minimum local : $g(1) = 2 + 3 - 12 + 1 = -6$.
\end{corrige}

\begin{aretenir}[L'essentiel de la section]
\begin{itemize}
\item $f' > 0$ sur $I$ (sauf points isolés) $\Rightarrow$ $f$ strictement croissante sur $I$ ; $f' < 0 \Rightarrow$ strictement décroissante ; $f' = 0 \Rightarrow$ constante. Théorème \textbf{admis}.
\item Extremum local en un point intérieur $c$ avec $f$ dérivable $\Rightarrow$ $f'(c) = 0$ (condition \textbf{nécessaire}, démonstration exigible).
\item $f'(c) = 0$ \textbf{ne suffit pas} : il faut un \textbf{changement de signe} de $f'$ en $c$ (contre-exemple : $x^3$ en $0$).
\item Méthode en 4 étapes : \textbf{dériver, factoriser, étudier le signe, conclure}.
\end{itemize}
\end{aretenir}

\coursec{Tableau de variation et optimisation}

Dans la section précédente, tu as appris à lire le signe de la dérivée pour connaître le sens de variation d'une fonction : $f' > 0$ signifie $f$ croissante, $f' < 0$ signifie $f$ décroissante. Mais un résultat aussi précieux mérite d'être \emph{présenté clairement}. C'est exactement le rôle du \cle{tableau de variation} : il condense en un seul schéma le domaine de définition, le signe de la dérivée, les variations de la fonction et ses extremums. Mieux encore : une fois ce tableau dressé, tu peux résoudre des problèmes concrets fascinants — maximiser le bénéfice d'une vendeuse de beignets, minimiser le coût d'une clôture, trouver la forme la plus économique d'un champ. C'est ce que l'on appelle un \cle{problème d'optimisation}, et c'est un grand classique du Baccalauréat.

\courssub{Construire un tableau de variation complet}

\begin{definition}[Tableau de variation]
Le \cle{tableau de variation} d'une fonction $f$ dérivable sur son domaine $D_f$ est un schéma à deux lignes :
\begin{itemize}
\item la première ligne indique les valeurs remarquables de $x$ : bornes du domaine, valeurs interdites, zéros de $f'$ ;
\item la deuxième ligne indique le signe de $f'(x)$ sur chaque intervalle, puis les variations de $f$ à l'aide de flèches ($\nearrow$ pour croissant, $\searrow$ pour décroissant), avec les valeurs des extremums et les limites aux bornes quand elles sont connues.
\end{itemize}
Une \textbf{double barre} verticale signale une valeur interdite (par exemple un dénominateur nul).
\end{definition}

\begin{methode}[Dresser un tableau de variation complet en 5 étapes]
\begin{enumerate}
\item \textbf{Déterminer le domaine} $D_f$ : repérer les dénominateurs (qui doivent être non nuls) et les racines carrées (dont le contenu doit être positif ou nul).
\item \textbf{Calculer la dérivée} $f'(x)$ en appliquant les formules des sections précédentes.
\item \textbf{Résoudre $f'(x) = 0$} et étudier le \textbf{signe de $f'(x)$} sur chaque intervalle du domaine (tableau de signes si nécessaire).
\item \textbf{En déduire les variations} : $f' > 0 \Rightarrow f \nearrow$ ; $f' < 0 \Rightarrow f \searrow$. Placer les flèches.
\item \textbf{Compléter} par les valeurs des extremums $f(a)$ aux points où $f'$ s'annule, et par les limites ou valeurs aux bornes du domaine.
\end{enumerate}
\end{methode}

\begin{popfigure}
\begin{tikzpicture}[scale=1.05]
% cadre
\draw[popline] (0,0) rectangle (13,4.6);
\draw[popline] (0,3.6) -- (13,3.6);
\draw[popline] (0,1.4) -- (13,1.4);
% vertical lines for x values: -inf, -1, 0, 2, +inf
\draw[popline] (2.2,0) -- (2.2,4.6);
\draw[popline] (4.6,0) -- (4.6,4.6);
\draw[popline] (7.0,0) -- (7.0,4.6);
\draw[popline] (9.4,0) -- (9.4,4.6);
\draw[popline] (11.8,0) -- (11.8,4.6);
% Double bar at x=0 (7.0)
\draw[popline, line width=1.2pt] (6.85,1.4) -- (6.85,3.6);
\draw[popline, line width=1.2pt] (7.15,1.4) -- (7.15,3.6);
% ligne x
\node at (1.1,4.0) {$x$};
\node at (3.4,4.0) {$-\infty$};
\node at (5.8,4.0) {$-1$};
\node at (8.2,4.0) {$0$};
\node at (10.6,4.0) {$2$};
\node at (12.4,4.0) {$+\infty$};
% ligne f'
\node at (1.1,2.5) {$f'(x)$};
\node at (4.0,2.5) {$+$};
\node at (6.2,2.5) {$0$};
\node at (7.0,2.5) {$||$};
\node at (8.2,2.5) {$-$};
\node at (10.0,2.5) {$0$};
\node at (11.2,2.5) {$+$};
% ligne f
\node at (1.1,0.7) {$f(x)$};
\node at (3.4,0.35) {$-\infty$};
\node[popblue] at (5.8,1.05) {$f(-1)$};
\node at (7.0,0.35) {$||$};
\node at (8.2,1.05) {$+\infty$};
\node[popblue] at (10.6,1.05) {$f(2)$};
\node at (12.4,1.05) {$+\infty$};
% arrows f
\draw[-{Stealth[length=3mm]}, popblue, very thick] (3.7,0.5) -- (5.5,1.0);
\draw[-{Stealth[length=3mm]}, popblue, very thick] (6.1,1.0) -- (6.8,0.5);
\draw[-{Stealth[length=3mm]}, popblue, very thick] (7.2,1.0) -- (10.0,0.5);
\draw[-{Stealth[length=3mm]}, popblue, very thick] (10.2,0.5) -- (12.0,1.0);
% annotations
\node[poporange, font=\small, align=center] at (5.8,1.85) {maximum\\ local};
\node[popgreen, font=\small, align=center] at (10.6,1.85) {minimum\\ local};
\node[popmuted, font=\small, align=center] at (8.2,0.9) {valeur\\ interdite\\ (asymptote)};
\end{tikzpicture}
\legende{Modèle de tableau de variation annoté : signe de $f'$, flèches de variation, extremums et double barre pour une valeur interdite (asymptote verticale en $x=0$).}
\end{popfigure}

\begin{attention}[Extremum local ou global ?]
Un point où $f'$ s'annule \textbf{et change de signe} donne un extremum \emph{local}. Pour affirmer qu'il s'agit d'un extremum \emph{sur tout le domaine}, il faut examiner le tableau complet : comparer avec les valeurs aux bornes. Un maximum local n'est pas toujours le maximum global !
\end{attention}

\courssub{Exemples de tableaux de variation}

\textbf{Fonction polynôme de degré 2.} Pour $f(x) = ax^2 + bx + c$ avec $a > 0$, on a $f'(x) = 2ax + b$, qui s'annule en $x_0 = -\dfrac{b}{2a}$. La fonction décroît puis croît : elle admet un \cle{minimum} en $x_0$, de valeur $f(x_0) = c - \dfrac{b^2}{4a}$. C'est la parabole « tournée vers le haut » que tu connais depuis la Seconde — la dérivation en donne une preuve élégante.

\textbf{Fonction polynôme de degré 3.} La dérivée est un trinôme du second degré : selon le signe de son discriminant, la cubique est soit monotone, soit dotée d'un maximum et d'un minimum locaux.

\begin{exemplebox}[{Tableau de variation de $f(x) = x + \dfrac{4}{x}$ sur $]0\,;\,+\infty[$}]
\textbf{Étape 1 — Domaine :} le dénominateur $x$ impose $x \neq 0$ ; on se restreint à $]0\,;\,+\infty[$.

\textbf{Étape 2 — Dérivée :} $f'(x) = 1 - \dfrac{4}{x^2} = \dfrac{x^2 - 4}{x^2}$.

\textbf{Étape 3 — Zéros et signe :} sur $]0\,;\,+\infty[$, $x^2 > 0$, donc $f'(x)$ a le signe de $x^2 - 4 = (x-2)(x+2)$. Comme $x > 0$, $f'(x) = 0 \iff x = 2$, et $f'(x) < 0$ sur $]0\,;\,2[$, $f'(x) > 0$ sur $]2\,;\,+\infty[$.

\textbf{Étape 4 — Variations :} $f$ décroît sur $]0\,;\,2]$, croît sur $[2\,;\,+\infty[$.

\textbf{Étape 5 — Valeurs :} $f(2) = 2 + \dfrac{4}{2} = 4$. De plus $f(x) \to +\infty$ quand $x \to 0^+$ et quand $x \to +\infty$.

\begin{center}
\begin{tabular}{|c|ccccc|}
\hline
$x$ & $0$ & & $2$ & & $+\infty$ \\
\hline
$f'(x)$ & \rule[-2mm]{0.4pt}{5mm}\rule[-2mm]{0.4pt}{5mm} & $-$ & $0$ & $+$ & \\
\hline
$f(x)$ & \rule[-2mm]{0.4pt}{5mm}\rule[-2mm]{0.4pt}{5mm} $+\infty$ & $\searrow$ & $4$ & $\nearrow$ & $+\infty$ \\
\hline
\end{tabular}
\end{center}

\textbf{Conclusion :} $f$ admet un \cle{minimum} égal à $4$, atteint en $x = 2$. Autrement dit, pour tout $x > 0$ : $x + \dfrac{4}{x} \geq 4$.
\end{exemplebox}

\textbf{Fonction avec racine carrée.} Pour $g(x) = \sqrt{x} - x$ sur $[0\,;\,+\infty[$, on a $g'(x) = \dfrac{1}{2\sqrt{x}} - 1 = \dfrac{1 - 2\sqrt{x}}{2\sqrt{x}}$ pour $x > 0$. Le signe dépend de $1 - 2\sqrt{x}$ : $g'(x) > 0 \iff \sqrt{x} < \dfrac{1}{2} \iff x < \dfrac{1}{4}$. Donc $g$ croît sur $\left[0\,;\,\dfrac{1}{4}\right]$ puis décroît, avec un maximum $g\left(\dfrac{1}{4}\right) = \dfrac{1}{2} - \dfrac{1}{4} = \dfrac{1}{4}$.

\courssub{Lire un tableau de variation}

Un tableau de variation bien construit répond à de nombreuses questions sans aucun calcul supplémentaire.

\begin{propriete}[Lectures essentielles d'un tableau de variation]
\begin{itemize}
\item \textbf{Extremums :} les « sommets » et « creux » des flèches donnent les maximums et minimums, avec la valeur de $x$ où ils sont atteints.
\item \textbf{Comparaison d'images :} si $f$ est croissante sur $[a\,;\,b]$ et $a \leq x_1 < x_2 \leq b$, alors $f(x_1) < f(x_2)$.
\item \textbf{Inégalités :} si le minimum de $f$ sur $I$ est $m$, alors $f(x) \geq m$ pour tout $x \in I$.
\item \textbf{Équations $f(x) = k$ (complément utile au Bac) :} en traçant mentalement la droite horizontale $y = k$ sur le tableau, on compte le nombre de solutions. Par exemple, si $f$ décroît de $+\infty$ vers $4$ puis croît vers $+\infty$, l'équation $f(x) = k$ a : deux solutions si $k > 4$, une solution si $k = 4$, aucune si $k < 4$.
\end{itemize}
\end{propriete}

\begin{exoresolu}[Nombre de solutions d'une équation]
Énoncé. On admet que la fonction $f(x) = x + \dfrac{4}{x}$ sur $]0\,;\,+\infty[$ a le tableau de variation établi plus haut. Déterminer, suivant les valeurs du réel $k$, le nombre de solutions de l'équation $f(x) = k$.
\tcblower
Solution. Le tableau montre que $f$ décroît de $+\infty$ vers $4$ sur $]0\,;\,2]$, puis croît de $4$ vers $+\infty$ sur $[2\,;\,+\infty[$.
\begin{itemize}
\item Si $k < 4$ : la valeur $k$ n'est jamais atteinte, l'équation n'a \textbf{aucune solution}.
\item Si $k = 4$ : la valeur est atteinte une seule fois, en $x = 2$ : \textbf{une solution}.
\item Si $k > 4$ : la valeur est atteinte une fois sur chaque branche (théorème des valeurs intermédiaires) : \textbf{deux solutions}.
\end{itemize}
Par exemple, $f(x) = 5$ équivaut à $x^2 - 5x + 4 = 0$, soit $x = 1$ ou $x = 4$ : deux solutions, conformément à la lecture du tableau.
\end{exoresolu}

\courssub{Problèmes d'optimisation}

Voici l'aboutissement de tout ce chapitre. Un \cle{problème d'optimisation} consiste à chercher la \emph{meilleure} situation possible : la plus grande aire, le plus grand bénéfice, le plus petit coût, la plus courte distance. La dérivation est l'outil qui transforme ce rêve en calcul.

\begin{methode}[Résoudre un problème d'optimisation]
\begin{enumerate}
\item \textbf{Choisir la variable} : identifier la grandeur inconnue $x$ et préciser son \textbf{domaine} imposé par le contexte (longueur positive, quantité entière, budget limité...).
\item \textbf{Modéliser} : exprimer la quantité à optimiser en fonction de $x$, en éliminant les autres inconnues grâce aux contraintes de l'énoncé. On obtient une fonction $f(x)$.
\item \textbf{Dériver et étudier} : calculer $f'(x)$, résoudre $f'(x) = 0$, dresser le tableau de variation sur le domaine du problème.
\item \textbf{Identifier l'extremum} cherché et vérifier qu'il appartient bien au domaine.
\item \textbf{Conclure en français, avec les unités} : répondre à la question posée et vérifier la cohérence de la réponse avec le contexte concret.
\end{enumerate}
\end{methode}

\begin{attention}[La valeur optimale doit appartenir au domaine du problème !]
C'est l'erreur la plus fréquente au Bac. Si le sommet de la parabole est en $x = 7{,}3$ mais que le domaine du problème est $[0\,;\,6]$, le maximum n'est \textbf{pas} en $7{,}3$ : il est atteint au bord, en $x = 6$ (car la fonction y est croissante). De même, une longueur négative ou une quantité de beignets non entière doivent te mettre la puce à l'oreille. \textbf{Toujours confronter le résultat mathématique au domaine concret.}
\end{attention}

\begin{exemplebox}[Problème type : le rectangle de périmètre fixe]
Un agriculteur de Bafoussam dispose de \SI{100}{m} de fil barbelé pour clôturer un champ rectangulaire. Quelles dimensions maximisent l'aire ?

\textbf{Modélisation :} soient $L$ et $\ell$ les dimensions. Le périmètre impose $2L + 2\ell = 100$, donc $\ell = 50 - L$. L'aire est $A(L) = L(50 - L) = -L^2 + 50L$, avec $L \in [0\,;\,50]$.

\textbf{Étude :} $A'(L) = -2L + 50$, qui s'annule en $L = 25$, positif avant, négatif après. Le tableau montre un maximum en $L = 25$.

\textbf{Conclusion :} $L = \ell = \SI{25}{m}$ : le champ optimal est un \textbf{carré} de \SI{25}{m} de côté, d'aire maximale $\SI{625}{m^2}$. Cohérent : les dimensions sont positives et le périmètre vaut bien \SI{100}{m}.
\end{exemplebox}

\begin{popfigure}
\begin{tikzpicture}[scale=0.55]
% rectangle
\draw[popblue, very thick, fill=popblueL] (0,0) rectangle (5,5);
\draw[poporange, very thick, dashed] (6.5,0) rectangle (14.5,2.2);
\node[popdark] at (2.5,-0.7) {$x = 25$ m};
\node[popdark, rotate=90] at (-0.6,2.5) {$50 - x = 25$ m};
\node[popdark] at (10.5,-0.7) {$x = 40$ m};
\node[popdark, rotate=90] at (15.1,1.1) {$10$ m};
\node[popblue] at (2.5,2.5) {$\mathcal{A} = 625$ m$^2$};
\node[poporange] at (10.5,1.1) {$\mathcal{A} = 400$ m$^2$};
\node[popmuted, align=center, font=\small] at (7,-2.2) {À périmètre égal (100 m), le carré\\ est le rectangle d'aire maximale.};
\end{tikzpicture}
\legende{Deux rectangles de même périmètre \SI{100}{m} : le carré (à gauche) a une aire supérieure au rectangle allongé (à droite).}
\end{popfigure}

\begin{exoresolu}[Le problème des beignets de l'accroche]
Énoncé. À l'entrée du lycée, une vendeuse de beignets (l'« accroche ») constate que si elle vend le beignet à $x$ francs CFA (avec $50 \leq x \leq 200$), elle en vend $n(x) = 400 - 2x$ par jour. Le coût de fabrication est de $30$ F par beignet. Son bénéfice quotidien, en francs, est donc
\[ B(x) = (x - 30)(400 - 2x). \]
À quel prix doit-elle vendre le beignet pour maximiser son bénéfice ? Quel est ce bénéfice maximal ?
\tcblower
Solution. \textbf{Étape 1 — Domaine :} $x \in [50\,;\,200]$ (prix imposé par le contexte ; au-delà de $200$ F, $n(x) \leq 0$, elle ne vend plus rien).

\textbf{Étape 2 — Développement et dérivée :}
\begin{align*}
B(x) &= (x-30)(400-2x) = -2x^2 + 460x - 12\,000,\\
B'(x) &= -4x + 460.
\end{align*}

\textbf{Étape 3 — Signe :} $B'(x) = 0 \iff x = 115$. De plus $B'(x) > 0$ pour $x < 115$ et $B'(x) < 0$ pour $x > 115$.

\textbf{Étape 4 — Tableau de variation :} $B$ croît sur $[50\,;\,115]$ puis décroît sur $[115\,;\,200]$. Elle admet donc un maximum en $x = 115$, qui appartient bien au domaine $[50\,;\,200]$ \textcolor{popgreen}{\textbf{(vérifié)}}.

\textbf{Étape 5 — Conclusion :}
\[ B(115) = (115 - 30)(400 - 230) = 85 \times 170 = 14\,450. \]
L'accroche doit vendre le beignet à \textbf{115 F} ; elle en vendra $400 - 2 \times 115 = 170$ par jour, pour un bénéfice maximal de \textbf{14\,450 F}. Cohérence : le prix est dans la fourchette observée, la quantité vendue est positive et entière, le bénéfice est supérieur à celui obtenu aux bornes ($B(50) = 6\,000$ F et $B(200) = 0$ F).
\end{exoresolu}

\begin{popfigure}
\begin{tikzpicture}
\begin{axis}[
  width=0.85\linewidth, height=6.2cm,
  axis lines=middle,
  xlabel={Prix $x$ (F)}, ylabel={Bénéfice $B(x)$ (F)},
  xmin=30, xmax=220, ymin=-2000, ymax=16000,
  xtick={50,115,200}, ytick={6000,14450},
  tick label style={font=\small},
  label style={font=\small},
  grid=both, grid style={popline!40},
]
\addplot[popcurve, domain=50:200, samples=100] {-2*x^2 + 460*x - 12000};
\addplot[mark=*, mark size=2.5pt, poporange] coordinates {(115,14450)};
\draw[dashed, poporange] (axis cs:115,0) -- (axis cs:115,14450);
\draw[dashed, poporange] (axis cs:30,14450) -- (axis cs:115,14450);
\node[poporange, font=\small, anchor=south west] at (axis cs:118,14450) {Maximum $(115\,;\,14\,450)$};
\end{axis}
\end{tikzpicture}
\legende{Courbe du bénéfice $B(x) = -2x^2 + 460x - 12\,000$ : parabole tournée vers le bas, dont le sommet donne le prix optimal de 115 F.}
\end{popfigure}

\begin{exercice}[Minimiser un coût de fabrication]
Un artisan de Foumban fabrique des boîtes cylindriques de volume imposé. On admet que le coût en matière première, en francs, d'une boîte de rayon $x$ cm (avec $x > 0$) est $C(x) = 2x^2 + \dfrac{128}{x}$. Déterminer le rayon qui minimise le coût, puis le coût minimal.
\end{exercice}

\begin{corrige}[Exercice 1]
$C'(x) = 4x - \dfrac{128}{x^2} = \dfrac{4x^3 - 128}{x^2}$. Comme $x^2 > 0$, $C'(x)$ a le signe de $4x^3 - 128$, qui s'annule pour $x^3 = 32$, soit $x = \sqrt[3]{32} = 2\sqrt[3]{4} \approx 3{,}17$. La fonction $x \mapsto x^3$ étant croissante, $C' < 0$ avant et $C' > 0$ après : $C$ admet un minimum en $x = 2\sqrt[3]{4}$ cm. Le coût minimal est $C\left(2\sqrt[3]{4}\right) \approx 2(3{,}17)^2 + \dfrac{128}{3{,}17} \approx 20{,}1 + 40{,}4 \approx 60{,}5$ F. La valeur trouvée est bien strictement positive, donc dans le domaine.
\end{corrige}

\begin{savaistu}[L'optimisation, un moteur du monde moderne]
Les problèmes d'optimisation sont au cœur de l'économie, de la logistique et de l'ingénierie. Les compagnies de téléphonie comme MTN ou Orange ajustent leurs tarifs pour maximiser leurs revenus tout en restant attractives ; les transporteurs de la SOTUC ou les agences de voyage entre Douala et Yaoundé cherchent les itinéraires qui minimisent le carburant consommé ; les ingénieurs de CAMTEL optimisent le tracé des câbles à fibre optique pour minimiser la longueur posée. Derrière chacune de ces décisions se cache une fonction à étudier... et un tableau de variation à dresser. En maîtrisant cette leçon, tu tiens l'outil mathématique le plus utilisé de la prise de décision !
\end{savaistu}

\begin{aretenir}[L'essentiel de la section]
\begin{itemize}
\item Le tableau de variation se construit en 5 étapes : domaine, dérivée, zéros et signe de $f'$, flèches, valeurs des extremums et bornes.
\item Il permet de lire les extremums, de comparer des images, de prouver des inégalités et de compter les solutions de $f(x) = k$.
\item Un problème d'optimisation se résout en modélisant par une fonction, en étudiant ses variations, puis en concluant \textbf{dans le domaine du problème} et \textbf{avec les unités}.
\item La cohérence de la réponse avec le contexte concret est une étape obligatoire, pas un luxe.
\end{itemize}
\end{aretenir}

\newpage

\coursec{Activité d'intégration}

\coursec{Dérivées et sens de variation}

\courssub{Rappels : nombre dérivé et tangente}

Soit une fonction $f$ définie sur un intervalle $I$ et $a \in I$. On se rappelle qu'en Première, nous avons introduit le \cle{taux d'accroissement} de $f$ entre $a$ et $x$ ($x \neq a$) :
\[ \tau(a,x) = \frac{f(x) - f(a)}{x - a}. \]
Ce quotient représente géométriquement la pente de la sécante $(AB)$ où $A(a\,;\,f(a))$ et $B(x\,;\,f(x))$ appartiennent à la courbe $\mathcal{C}_f$.

\begin{definition}[Nombre dérivé]
Si la fonction $f$ admet une limite finie lorsque $x$ tend vers $a$ pour le taux d'accroissement, on dit que $f$ est \textbf{dérivable} en $a$. Cette limite est appelée \textbf{nombre dérivé} de $f$ en $a$ et se note $f'(a)$ :
\[ f'(a) = \lim_{x \to a} \frac{f(x) - f(a)}{x - a} = \lim_{h \to 0} \frac{f(a+h) - f(a)}{h}. \]
\end{definition}

\begin{propriete}[Interprétation géométrique : la tangente]
Si $f$ est dérivable en $a$, la courbe $\mathcal{C}_f$ admet une \textbf{tangente} au point $A(a\,;\,f(a))$. Cette tangente est la limite (si elle existe) des sécantes $(AB)$ lorsque $B$ tend vers $A$ sur la courbe. Son équation cartésienne est :
\[ y = f'(a)(x - a) + f(a). \]
Le coefficient directeur de cette tangente est exactement $f'(a)$.
\end{propriete}

\begin{popfigure}
\begin{tikzpicture}[scale=1.2, domain=-0.5:3.5]
\draw[->] (-0.5,0) -- (4,0) node[right] {$x$};
\draw[->] (0,-0.5) -- (0,3.5) node[above] {$y$};
\draw[popcurve, thick, domain=0:3.2,samples=100] plot (\x, {0.5*\x*\x - 0.5*\x + 1});
\node[above right] at (3.2, 4.1) {$\mathcal{C}_f$};
\coordinate (A) at (1,1);
\fill[popblue] (A) circle (2pt);
\node[below left, popblue] at (A) {$A(1\,;\,f(1))$};
% Secantes
\coordinate (B1) at (2.5, 2.375);
\draw[poporange, dashed] (A) -- (B1);
\node[above right, poporange] at (B1) {$B_1$};
\coordinate (B2) at (1.8, 1.52);
\draw[poporange!70!popdark, dashed] (A) -- (B2);
\node[above right, poporange!70!popdark] at (B2) {$B_2$};
\coordinate (B3) at (1.3, 1.145);
\draw[poporange!40!popdark, dashed] (A) -- (B3);
\node[above right, poporange!40!popdark] at (B3) {$B_3$};
% Tangente
\draw[popgreen, thick, domain=0:2.5] plot (\x, {0.5*\x + 0.5});
\node[above left, popgreen] at (2.5, 1.75) {Tangente en $A$};
\end{tikzpicture}
\legende{La tangente est la position limite des sécantes $(AB)$ quand $B \to A$.}
\end{popfigure}

\begin{savaistu}[Le vocabulaire du Bac]
Au Cameroun, on distingue soigneusement :
\begin{itemize}
\item le \textbf{nombre dérivé} $f'(a)$ (un \emph{nombre}, la pente en un point) ;
\item la \textbf{fonction dérivée} $f'$ (une \emph{fonction}, qui associe à chaque $x$ le nombre $f'(x)$).
\end{itemize}
L'examen attend que tu passes sans ambiguïté de l'un à l'autre : « calculer $f'(2)$ » vs « déterminer la fonction dérivée $f'$ sur $I$ ».
\end{savaistu}

\courssub{Dérivées des fonctions usuelles}

Le programme de Terminale A exige la maîtrise des dérivées des fonctions de référence suivantes. Ces formules sont à connaître \textbf{par cœur} et s'appliquent sur les intervalles où les fonctions sont définies et dérivables.

\begin{propriete}[Tableau des dérivées usuelles (à savoir par cœur)]
\begin{center}
\begin{tabularx}{\linewidth}{|l|X|l|}
\hline
\rowcolor{popblueL} \textbf{Fonction $f(x)$} & \textbf{Dérivée $f'(x)$} & \textbf{Domaine de dérivabilité} \\
\hline
$k$ (constante) & $0$ & $\mathbb{R}$ \\
\hline
$x^n$ ($n \in \mathbb{N}^*$) & $n x^{n-1}$ & $\mathbb{R}$ \\
\hline
$x^\alpha$ ($\alpha \in \mathbb{Q}$, $\alpha \notin \mathbb{N}$) & $\alpha x^{\alpha-1}$ & $]0\,;\,+\infty[$ (ou $\mathbb{R}^*$ si dénominateur impair) \\
\hline
$\sqrt{x}$ & $\dfrac{1}{2\sqrt{x}}$ & $]0\,;\,+\infty[$ \\
\hline
$\dfrac{1}{x}$ & $-\dfrac{1}{x^2}$ & $\mathbb{R}^*$ \\
\hline
\end{tabularx}
\end{center}
\emph{Remarque :} Les fonctions exponentielle, logarithme, sinus et cosinus font partie du programme de Terminale D/C. En Série A, on se concentre sur les fonctions algébriques (polynômes, fractions rationnelles, racines).
\end{propriete}

\begin{exemplebox}[Application directe]
Soit $f(x) = 3x^4 - 5x^2 + \dfrac{2}{x} - 7\sqrt{x}$ définie sur $]0\,;\,+\infty[$.
\[ f'(x) = 3(4x^3) - 5(2x) + 2\left(-\frac{1}{x^2}\right) - 7\left(\frac{1}{2\sqrt{x}}\right) = 12x^3 - 10x - \frac{2}{x^2} - \frac{7}{2\sqrt{x}}. \]
\end{exemplebox}

\courssub{Opérations sur les dérivées}

Soient $u$ et $v$ deux fonctions dérivables sur un intervalle $I$, et $k \in \mathbb{R}$.

\begin{propriete}[Règles de dérivation (linéarité, produit, quotient)]
\begin{itemize}
\item \textbf{Somme :} $(u+v)' = u' + v'$.
\item \textbf{Multiple par une constante :} $(ku)' = k\,u'$.
\item \textbf{Produit :} $(uv)' = u'v + uv'$.
\item \textbf{Quotient :} $\left(\dfrac{u}{v}\right)' = \dfrac{u'v - uv'}{v^2}$, pour tout $x$ tel que $v(x) \neq 0$.
\end{itemize}
\end{propriete}

\begin{attention}[Pièges classiques sur les opérations]
\begin{itemize}
\item \textbf{Produit :} $(uv)' \neq u'v'$ ! C'est l'erreur n°1. Pense à la règle : « dérivée du premier $\times$ second \textbf{plus} premier $\times$ dérivée du second ».
\item \textbf{Quotient :} L'ordre au numérateur est crucial : $u'v - uv'$ (et non l'inverse). Le dénominateur est $v^2$ (jamais $v$ seul).
\item \textbf{Domaine :} La dérivée d'un quotient n'existe pas aux zéros du dénominateur $v$. Précise toujours le domaine de $f'$.
\end{itemize}
\end{attention}

\begin{exemplebox}[Dérivée d'une fraction rationnelle]
Soit $f(x) = \dfrac{x^2 + 1}{x - 2}$ sur $\mathbb{R} \setminus \{2\}$.
$u(x) = x^2+1 \Rightarrow u'(x) = 2x$ ; $v(x) = x-2 \Rightarrow v'(x) = 1$.
\[ f'(x) = \frac{2x(x-2) - (x^2+1)\cdot 1}{(x-2)^2} = \frac{2x^2 - 4x - x^2 - 1}{(x-2)^2} = \frac{x^2 - 4x - 1}{(x-2)^2}. \]
\end{exemplebox}

\courssub{Dérivée d'une fonction composée}

C'est un outil puissant pour dériver des fonctions comme $\sqrt{3x-1}$, $(x^2+1)^5$ ou $\dfrac{1}{x^2+3}$.

\begin{propriete}[Formule de dérivation chainée (exigible au Bac)]
Soient $u$ dérivable sur $I$ et $v$ dérivable sur $J$ tel que $u(I) \subset J$.
La fonction composée $f = v \circ u$ (définie par $f(x) = v(u(x))$) est dérivable sur $I$ et :
\[ f'(x) = u'(x) \times v'(u(x)). \]
En mots : \textbf{dérivée de l'intérieur $\times$ dérivée de l'extérieur évaluée en l'intérieur}.
\end{propriete}

\begin{methode}[Dériver une fonction composée : 3 étapes]
\begin{enumerate}
\item \textbf{Identifier} l'intérieur $u(x)$ et l'extérieur $v(X)$ (avec $X = u(x)$).
\item \textbf{Calculer} $u'(x)$ et $v'(X)$ séparément.
\item \textbf{Appliquer} la formule : $f'(x) = u'(x) \cdot v'(u(x))$ et simplifier.
\end{enumerate}
\end{methode}

\begin{exemplebox}[Exemples de fonctions composées]
\begin{enumerate}
\item $f(x) = \sqrt{3x-1} = (3x-1)^{1/2}$ sur $]1/3\,;\,+\infty[$.
$u(x)=3x-1$, $v(X)=\sqrt{X}$. $u'=3$, $v'(X)=\frac{1}{2\sqrt{X}}$.
$f'(x) = 3 \times \frac{1}{2\sqrt{3x-1}} = \frac{3}{2\sqrt{3x-1}}$.
\item $g(x) = (x^2+2x)^4$ sur $\mathbb{R}$.
$u(x)=x^2+2x$, $v(X)=X^4$. $u'=2x+2$, $v'(X)=4X^3$.
$g'(x) = (2x+2) \times 4(x^2+2x)^3 = 8(x+1)(x^2+2x)^3$.
\item $h(x) = \dfrac{1}{x^2+1} = (x^2+1)^{-1}$ sur $\mathbb{R}$.
$u(x)=x^2+1$, $v(X)=X^{-1}$. $u'=2x$, $v'(X)=-X^{-2}$.
$h'(x) = 2x \times (-(x^2+1)^{-2}) = -\frac{2x}{(x^2+1)^2}$.
\end{enumerate}
\end{exemplebox}

\courssub{Signe de la dérivée et sens de variation}

C'est le lien fondamental qui permet d'étudier les variations d'une fonction sans tracer sa courbe.

\begin{propriete}[Théorème fondamental : signe de $f'$ et variations de $f$]
Soit $f$ une fonction dérivable sur un intervalle $I$.
\begin{itemize}
\item Si $f'(x) \geq 0$ sur $I$ (avec un ensemble de zéros \emph{isolés} ou ne contenant pas d'intervalle), alors $f$ est \textbf{croissante} sur $I$. Si $f'(x) > 0$ (sauf points isolés), $f$ est \textbf{strictement croissante}.
\item Si $f'(x) \leq 0$ sur $I$ (même condition sur les zéros), alors $f$ est \textbf{décroissante} sur $I$. Si $f'(x) < 0$ (sauf points isolés), $f$ est \textbf{strictement décroissante}.
\item Si $f'(x) = 0$ sur $I$, alors $f$ est \textbf{constante} sur $I$.
\end{itemize}
\end{propriete}

\begin{aretenir}[Extremums locaux et dérivée]
Soit $f$ dérivable sur $I$ et $a \in I$ tel que $f'(a) = 0$.
\begin{itemize}
\item Si $f'$ change de signe de $+$ en $-$ en $a$ : $f$ admet un \textbf{maximum local} en $a$.
\item Si $f'$ change de signe de $-$ en $+$ en $a$ : $f$ admet un \textbf{minimum local} en $a$.
\item Si $f'$ ne change pas de signe en $a$ : pas d'extremum (point d'inflexion à tangente horizontale).
\end{itemize}
\emph{Attention :} La condition $f'(a)=0$ est \textbf{nécessaire} mais pas \textbf{suffisante} pour avoir un extremum. Le tableau de signes de $f'$ est indispensable pour conclure.
\end{aretenir}

\begin{methode}[Étudier le signe d'une dérivée (affine, trinôme, fraction rationnelle)]
\begin{enumerate}
\item Factoriser $f'(x)$ au maximum.
\item Trouver les zéros de $f'$ (résoudre $f'(x)=0$).
\item Construire un \textbf{tableau de signes} de $f'$ sur l'intervalle d'étude.
\item En déduire les intervalles de positivité/négativité.
\end{enumerate}
\end{methode}

\courssub{Tableau de variation et optimisation}

Le tableau de variation synthétise l'étude complète. C'est l'outil de communication obligatoire au Baccalauréat.

\begin{methode}[Construire un tableau de variation complet]
\begin{enumerate}
\item \textbf{Domaine :} Préciser l'intervalle $I$ d'étude (souvent imposé par le contexte).
\item \textbf{Dérivée :} Calculer $f'(x)$ et simplifier.
\item \textbf{Signe de $f'$ :} Tableau de signes de $f'$ sur $I$ (zéros, bornes).
\item \textbf{Variations :} Reporter les flèches ($\nearrow$ ou $\searrow$) et les valeurs $f(x)$ aux points remarquables (zéros de $f'$, bornes de $I$).
\item \textbf{Limites aux bornes :} Si $I$ est ouvert ou non borné, calculer les limites de $f$ aux bornes.
\end{enumerate}
\end{methode}

\begin{popfigure}
\begin{tikzpicture}[scale=0.8, domain=0:5]
\draw[->] (-0.5,0) -- (5.5,0) node[right] {$x$};
\draw[->] (0,-0.5) -- (0,5.5) node[above] {$y$};
\draw[popcurve, thick, domain=0:5,samples=200] plot (\x, {-0.5*\x*\x + 2.5*\x});
\fill[popblue] (2.5, 3.125) circle (2pt);
\node[above, popblue] at (2.5, 3.125) {Max $A(25)$};
\draw[dashed, popmuted] (2.5,0) node[below] {$25$} -- (2.5, 3.125);
\draw[dashed, popmuted] (0,3.125) node[left] {$1250$} -- (2.5, 3.125);
\node[popdark] at (2.5, -0.8) {Tableau de variation $\leftrightarrow$ Courbe};
\end{tikzpicture}
\legende{Le tableau de variation est la « radiographie » textuelle de la courbe.}
\end{popfigure}

\courssub{Activité d'optimisation : Le champ de l'oncle Mbarga}

\courssub{Objectifs de l'activité}

À l'issue de cette activité, tu seras capable de :

\begin{itemize}
\item modéliser une situation concrète d'optimisation par une fonction numérique d'une variable réelle ;
\item calculer la dérivée d'une fonction polynôme et d'une fonction composée simple ;
\item étudier le signe de la dérivée et dresser le tableau de variation complet d'une fonction ;
\item déterminer l'extremum d'une fonction sur un intervalle et interpréter le résultat dans le contexte ;
\item comparer deux situations d'optimisation et vérifier la cohérence des résultats avec la réalité.
\end{itemize}

\courssub{Situation}

L'oncle Mbarga possède, dans le quartier Nkolbisson à Yaoundé, un terrain en bordure d'un long mur en parpaings qui délimite sa concession. Il souhaite aménager un champ rectangulaire pour y cultiver du manioc et des légumes, afin de ravitailler le marché de Mokolo.

Il dispose exactement de \SI{100}{m} de fil barbelé acheté à la quincaillerie du quartier. Le mur, parfaitement rectiligne et assez long, formera l'un des grands côtés du champ : le fil barbelé ne servira donc que pour \textbf{trois côtés} du rectangle (une longueur parallèle au mur et deux largeurs perpendiculaires au mur).

On note $x$ (en mètres) la largeur du champ, c'est-à-dire la longueur de chacun des deux côtés perpendiculaires au mur, et $L$ la longueur du côté parallèle au mur. On note $A(x)$ l'aire du champ en mètres carrés.

L'oncle Mbarga se pose une question toute simple : \emph{comment choisir les dimensions du champ pour obtenir la plus grande surface cultivable possible avec ses \SI{100}{m} de fil ?}

\begin{popfigure}
\begin{tikzpicture}[scale=0.9]
\fill[poppurpleL] (0,0) rectangle (8,0.25);
\draw[line width=1.5pt, popdark] (0,0.25) -- (8,0.25);
\node[popdark] at (4,0.6) {\textbf{Mur de la concession}};
\draw[line width=1.2pt, popblue] (0.8,0) -- (0.8,-2.5) -- (7.2,-2.5) -- (7.2,0);
\fill[popgreenL, opacity=0.6] (0.8,0) rectangle (7.2,-2.5);
\node at (4,-1.25) {Champ};
\node[left, popblue] at (0.8,-1.25) {$x$};
\node[right, popblue] at (7.2,-1.25) {$x$};
\node[below, popblue] at (4,-2.5) {$L$};
\end{tikzpicture}
\legende{Le champ rectangulaire adossé au mur : le fil barbelé ne sert que pour trois côtés.}
\end{popfigure}

\courssub{Consignes}

\begin{enumerate}
\item \textbf{Modélisation.} Exprime $L$ en fonction de $x$, en utilisant le fait que la longueur totale de fil disponible est \SI{100}{m}. Justifie que $x$ appartient à l'intervalle $]0\,;\,50[$, puis montre que l'aire du champ est donnée par :
\[ A(x) = 100x - 2x^2. \]
\item \textbf{Calcul de dérivée.} Calcule $A'(x)$ en utilisant les dérivées des fonctions usuelles et les opérations sur les dérivées.
\item \textbf{Étude du signe de la dérivée.} Résous l'équation $A'(x) = 0$, puis étudie le signe de $A'(x)$ sur l'intervalle $]0\,;\,50[$.
\item \textbf{Tableau de variation.} Dresse le tableau de variation complet de la fonction $A$ sur $]0\,;\,50[$ (signe de $A'$, sens de variation, valeurs remarquables et limites aux bornes).
\item \textbf{Optimisation.} Déduis-en les dimensions du champ qui rendent l'aire maximale, ainsi que la valeur de cette aire maximale. Rédige une phrase de conclusion à destination de l'oncle Mbarga.
\item \textbf{Prolongement : le cas sans mur.} Supposons maintenant qu'il n'y ait pas de mur et que les \SI{100}{m} de fil servent à clôturer les quatre côtés d'un champ rectangulaire de largeur $x$ et de longueur $\ell$. Montre que l'aire est alors $B(x) = 50x - x^2$, détermine de la même manière les dimensions optimales et l'aire maximale. Compare avec la situation avec mur : le résultat est-il cohérent avec le contexte ?
\end{enumerate}

\courssub{Ressources à mobiliser}

\begin{aretenir}[Boîte à outils de l'activité]
\begin{itemize}
\item \textbf{Dérivées usuelles :} $(x^n)' = n\,x^{n-1}$ pour tout entier $n \geq 1$ ; en particulier $(x)' = 1$ et $(x^2)' = 2x$.
\item \textbf{Opérations :} $(u+v)' = u'+v'$ ; $(k\,u)' = k\,u'$ pour tout réel $k$.
\item \textbf{Lien dérivée/variations :} soit $f$ dérivable sur un intervalle $I$. Si $f'(x) > 0$ sur $I$ (sauf éventuellement en des points isolés), alors $f$ est strictement croissante sur $I$ ; si $f'(x) < 0$, alors $f$ est strictement décroissante sur $I$. Si $f'$ s'annule en $a$ en changeant de signe de $+$ vers $-$, alors $f$ admet un maximum en $a$.
\item \textbf{Tableau de variation :} première ligne pour $x$, deuxième ligne pour le signe de $f'(x)$, troisième ligne pour les flèches de variation de $f$ avec les valeurs aux points remarquables.
\item \textbf{Modélisation :} périmètre d'un rectangle : $2(L + x)$ ; aire d'un rectangle : $L \times x$.
\end{itemize}
\end{aretenir}

\courssub{Démarche et corrigé commenté}

\begin{exoresolu}[Le champ de l'oncle Mbarga : corrigé complet]
\textbf{Énoncé :} On reprend les six questions des consignes.
\tcblower
\textbf{1. Modélisation.} Le fil barbelé sert pour les deux largeurs et la longueur parallèle au mur, donc :
\[ 2x + L = 100 \quad \text{d'où} \quad L = 100 - 2x. \]
Comme $x$ est une longueur, $x > 0$. Comme $L$ est aussi une longueur, $L > 0$, c'est-à-dire $100 - 2x > 0$, soit $x < 50$. Donc $x \in \,]0\,;\,50[$. L'aire du champ est alors :
\[ A(x) = L \times x = (100 - 2x)\,x = 100x - 2x^2. \]

\textbf{2. Calcul de la dérivée.} La fonction $A$ est une fonction polynôme, dérivable sur $\mathbb{R}$, donc sur $]0\,;\,50[$. Par linéarité de la dérivation :
\[ A'(x) = 100 \times (x)' - 2 \times (x^2)' = 100 \times 1 - 2 \times 2x = 100 - 4x. \]

\textbf{3. Signe de la dérivée.} On résout :
\[ A'(x) = 0 \iff 100 - 4x = 0 \iff x = 25. \]
La fonction $A'$ est une fonction affine de coefficient directeur $-4 < 0$, donc strictement décroissante : $A'(x) > 0$ pour $x < 25$ et $A'(x) < 0$ pour $x > 25$. Sur $]0\,;\,50[$ :
\begin{itemize}
\item si $x \in \,]0\,;\,25[$, alors $A'(x) > 0$ : $A$ est strictement croissante ;
\item si $x \in \,]25\,;\,50[$, alors $A'(x) < 0$ : $A$ est strictement décroissante.
\end{itemize}

\textbf{4. Tableau de variation.} On calcule la valeur remarquable :
\[ A(25) = 100 \times 25 - 2 \times 25^2 = 2500 - 1250 = 1250. \]
Aux bornes de l'intervalle ouvert : quand $x \to 0^+$, $A(x) = x(100-2x) \to 0$ ; quand $x \to 50^-$, $A(x) = x(100-2x) \to 0$ (le champ devient un couloir sans surface).

\begin{center}
\begin{tabular}{|c|ccccc|}
\hline
\rowcolor{popblueL} $x$ & $0$ & & $25$ & & $50$ \\
\hline
Signe de $A'(x)$ & & $+$ & $0$ & $-$ & \\
\hline
Variations de $A$ & $0$ & $\nearrow$ & $1250$ & $\searrow$ & $0$ \\
\hline
\end{tabular}
\end{center}

\textbf{5. Optimisation.} La dérivée s'annule en $x = 25$ en passant du signe $+$ au signe $-$ : la fonction $A$ admet donc un \textbf{maximum} en $x = 25$, qui vaut \SI{1250}{m^2}. La longueur correspondante est :
\[ L = 100 - 2 \times 25 = 50 \ \text{m}. \]
\emph{Conclusion pour l'oncle Mbarga :} il doit clôturer un champ de \SI{25}{m} de largeur et de \SI{50}{m} de longueur le long du mur ; il obtiendra ainsi une surface cultivable maximale de \SI{1250}{m^2}, soit un huitième d'hectare.

\textbf{6. Prolongement : le cas sans mur.} Les quatre côtés sont clôturés : $2x + 2\ell = 100$, donc $\ell = 50 - x$, avec $x \in \,]0\,;\,50[$. L'aire est :
\[ B(x) = x(50 - x) = 50x - x^2. \]
On dérive : $B'(x) = 50 - 2x$, qui s'annule en $x = 25$, est positive avant et négative après. Le tableau de variation montre un maximum en $x = 25$, avec $\ell = 50 - 25 = 25$ m et :
\[ B(25) = 25 \times 25 = 625 \ \text{m}^2. \]
Le champ optimal sans mur est donc un \textbf{carré} de \SI{25}{m} de côté, d'aire \SI{625}{m^2}.

\emph{Comparaison et cohérence :} avec le mur, l'aire maximale ($\SI{1250}{m^2}$) est exactement le \textbf{double} de celle obtenue sans mur ($\SI{625}{m^2}$). C'est parfaitement cohérent : le mur remplace gratuitement \SI{50}{m} de clôture, ce qui permet de réinvestir le fil dans les autres côtés. De plus, dans le cas avec mur, la longueur optimale ($L = 50$ m) est le double de la largeur ($x = 25$ m) : le mur « complète » symboliquement le rectangle en un carré de $50 \times 50$ dont le champ serait la moitié. Le résultat mathématique rejoint ainsi l'intuition géométrique.
\end{exoresolu}

\begin{attention}[Pièges fréquents dans ce type de problème]
\begin{itemize}
\item \textbf{Oublier le domaine.} Écrire $A(x) = 100x - 2x^2$ sans préciser que $x \in \,]0\,;\,50[$ est une erreur de modélisation : hors de cet intervalle, la formule n'a plus de sens physique (longueur négative).
\item \textbf{Se tromper de périmètre.} Avec le mur, il n'y a que \emph{trois} côtés de fil : $2x + L = 100$ et non $2x + 2L = 100$. Lis toujours attentivement l'énoncé.
\item \textbf{Confondre $x$ et $A(x)$.} La question « quelle est l'aire maximale ? » demande la valeur $A(25) = 1250$, pas $x = 25$. Distingue toujours la variable, les dimensions et l'extremum.
\item \textbf{Conclure sans justification.} Il ne suffit pas de résoudre $A'(x) = 0$ : il faut vérifier, par le signe de $A'$ (ou le tableau de variation), qu'il s'agit bien d'un \emph{maximum} et non d'un minimum.
\item \textbf{Négliger les bornes.} Sur un intervalle ouvert $]a\,;\,b[$, les valeurs aux bornes sont des \emph{limites}. Vérifie toujours si l'extremum est global (absolu) ou seulement local. Ici, le maximum en $25$ est global car $A(x) \leq 1250$ pour tout $x \in ]0,50[$.
\end{itemize}
\end{attention}

\courssub{Bilan de la leçon}

Cette leçon t'a permis de maîtriser les points essentiels du programme officiel :

\begin{itemize}
\item \textbf{La boîte à outils différentielle :} dérivées des fonctions usuelles (puissances, racine, inverse), règles de dérivation (somme, produit, quotient, composée). Ces formules sont le socle technique.
\item \textbf{Le lien fondamental :} le signe de $f'$ commande les variations de $f$. C'est le théorème à invoquer pour justifier toute monotonie.
\item \textbf{Le tableau de variation :} c'est l'outil de synthèse rigoureuse. Il doit être présenté proprement : ligne $x$, ligne signe $f'$, ligne variations de $f$ avec valeurs exactes (pas d'approximations décimales sauf demande explicite).
\item \textbf{La démarche d'optimisation (compétence clé du Bac) :}
\begin{enumerate}
\item \cle{Modéliser} : variable, fonction objectif, domaine (contraintes réelles).
\item \cle{Dériver} : calcul algébrique rigoureux, simplification.
\item \cle{Étudier le signe} : tableau de signes de $f'$, zéros, intervalles.
\item \cle{Conclure} : tableau de variation, extremum, interprétation dans le contexte (phrase de réponse).
\end{enumerate}
\end{itemize}

Tu es maintenant prêt à traiter tout problème d'optimisation simple au Baccalauréat : rendement d'une parcelle agricole, coût de fabrication d'un objet, recette maximale d'un commerçant à la gare routière, distance minimale pour un câble CAMTEL. La méthode reste toujours la même : \cle{modéliser}, \cle{dériver}, \cle{étudier le signe}, \cle{conclure}.

\newpage

\coursec{Méthodes et exercices résolus}

\coursec{Dérivées et sens de variation}

\courssub{Rappels : nombre dérivé et tangente}

\begin{definition}[Nombre dérivé en un point]
Soit $f$ une fonction définie sur un intervalle $I$ et $a \in I$. Si la limite
\[ \lim_{h \to 0} \frac{f(a+h) - f(a)}{h} \]
existe et est finie, on dit que $f$ est \cle{dérivable} en $a$ et on appelle cette limite le \cle{nombre dérivé} de $f$ en $a$, noté $f'(a)$.
\end{definition}

\begin{propriete}[Équation de la tangente]
Si $f$ est dérivable en $a$, la droite tangente $\mathcal{T}$ à la courbe $\mathcal{C}_f$ au point $A(a\,;\,f(a))$ existe. Son coefficient directeur est $f'(a)$ et son équation réduite s'écrit :
\[ y = f'(a)(x - a) + f(a). \]
\end{propriete}

\begin{exemplebox}[Nombre dérivé et tangente]
Soit $f(x) = x^2$ et $a = 1$. Le taux d'accroissement est $\frac{(1+h)^2 - 1}{h} = \frac{2h+h^2}{h} = 2+h \to 2$. Donc $f'(1)=2$.
La tangente en $A(1\,;\,1)$ a pour équation $y = 2(x-1)+1 = 2x-1$.
\end{exemplebox}

\begin{popfigure}
\begin{tikzpicture}[scale=0.8, >=Stealth]
\def\f{(\x)^2}
\def\tang{2*\x-1}
\draw[->] (-0.5,0) -- (3,0) node[right] {$x$};
\draw[->] (0,-0.5) -- (0,5) node[above] {$y$};
\draw[domain=-0.2:2.2,smooth,variable=\x,popcurve,thick] plot ({\x},{\f});
\draw[domain=0:2.5,smooth,variable=\x,poporange,dashed,thick] plot ({\x},{\tang});
\fill[popblue] (1,1) circle (2pt) node[above left] {$A(1;1)$};
\draw[popmuted] (1,0) node[below] {$1$} -- (1,1);
\draw[popmuted] (0,1) node[left] {$1$} -- (1,1);
\node[poporange, right] at (2.2,3) {Tangente $y=2x-1$};
\node[popblue, right] at (2.2,4) {Courbe $y=x^2$};
\end{tikzpicture}
\legende{Interprétation géométrique du nombre dérivé : pente de la tangente.}
\end{popfigure}

\courssub{Dérivées des fonctions usuelles}

\begin{propriete}[Tableau des dérivées usuelles (à connaître par cœur)]
Pour tout $x$ dans l'ensemble de définition :
\begin{center}
\begin{tabular}{|c|c|c|}
\hline
$f(x)$ & $f'(x)$ & Condition / Domaine \\
\hline
$k$ (constante) & $0$ & $\mathbb{R}$ \\
$x^n$ ($n \in \mathbb{N}$) & $n x^{n-1}$ & $\mathbb{R}$ \\
$x^\alpha$ ($\alpha \in \mathbb{R}$) & $\alpha x^{\alpha-1}$ & $]0\,;\,+\infty[$ (ou $\mathbb{R}^*$ selon $\alpha$) \\
$\sqrt{x}$ & $\dfrac{1}{2\sqrt{x}}$ & $]0\,;\,+\infty[$ \\
$\dfrac{1}{x}$ & $-\dfrac{1}{x^2}$ & $\mathbb{R}^*$ \\
\hline
\end{tabular}
\end{center}
\textbf{Remarque :} En Terminale A, on utilise essentiellement les fonctions polynomiales, rationnelles et la racine carrée. Les fonctions $\ln$, $\exp$, $\sin$, $\cos$ sont au programme de Terminale D/S.
\end{propriete}

\courssub{Opérations sur les dérivées}

\begin{propriete}[Règles de dérivation (somme, produit, quotient)]
Soient $u$ et $v$ deux fonctions dérivables sur un intervalle $I$.
\begin{itemize}
\item \textbf{Somme :} $(u+v)' = u' + v'$.
\item \textbf{Produit par une constante :} $(ku)' = k u'$.
\item \textbf{Produit :} $(uv)' = u'v + uv'$.
\item \textbf{Quotient :} $\left(\dfrac{u}{v}\right)' = \dfrac{u'v - uv'}{v^2}$, sur l'ensemble où $v(x) \neq 0$.
\end{itemize}
\end{propriete}

\courssub{Dérivée d'une fonction composée}

\begin{propriete}[Formule de dérivation chainée]
Soient $u$ dérivable sur $I$ et $v$ dérivable sur $u(I)$. La fonction composée $v \circ u$ est dérivable sur $I$ et :
\[ (v \circ u)'(x) = u'(x) \times v'(u(x)). \]
\textbf{Cas particulier (puissance d'une fonction) :} Pour $n \in \mathbb{N}^*$, $(u^n)' = n\,u'\,u^{n-1}$.
\end{propriete}

\begin{methode}[Calculer une dérivée sans se tromper]
\begin{enumerate}
\item \textbf{Identifier la structure} de $f$ : somme, produit, quotient, ou composée $v\circ u$. C'est l'étape décisive.
\item \textbf{Justifier la dérivabilité} : « $f$ est dérivable sur $I$ comme somme/produit/quotient/composée de fonctions dérivables sur $I$ ».
\item \textbf{Appliquer la formule adaptée} en posant explicitement $u$, $v$, $u'$, $v'$ si nécessaire.
\item \textbf{Simplifier et factoriser} le résultat : une dérivée factorisée rend l'étude du signe immédiate.
\item \textbf{Vérifier mentalement} (degré, coefficient dominant, valeur en un point).
\end{enumerate}
\end{methode}

\begin{exoresolu}[Dérivées de fonctions usuelles et composées]
On considère les fonctions définies par :
\[ f(x) = 3x^3 - 5x^2 + 2x - 7 \ ; \qquad g(x) = \frac{2x-1}{x+3} \ ; \qquad h(x) = (2x-1)^4. \]
\begin{enumerate}
\item Justifier la dérivabilité de $f$ sur $\mathbb{R}$ et calculer $f'(x)$.
\item Justifier la dérivabilité de $g$ sur $]-3\,;\,+\infty[$ et calculer $g'(x)$.
\item Calculer $h'(x)$ pour tout réel $x$.
\end{enumerate}
\tcblower
\textbf{1. Dérivée de $f$.}

$f$ est une fonction \cle{polynôme}, donc $f$ est dérivable sur $\mathbb{R}$ comme somme de fonctions dérivables sur $\mathbb{R}$.

On dérive terme à terme avec $(x^n)' = n x^{n-1}$ et $(ku)' = ku'$ :
\begin{align*}
f'(x) &= 3\times 3x^2 - 5\times 2x + 2\times 1 - 0 \\
&= 9x^2 - 10x + 2.
\end{align*}
\textbf{Conclusion :} pour tout réel $x$, $f'(x) = 9x^2 - 10x + 2$.

\textbf{2. Dérivée de $g$.}

$g$ est un \cle{quotient} de deux fonctions polynômes dérivables, dont le dénominateur $x+3$ ne s'annule pas sur $]-3\,;\,+\infty[$ : $g$ est donc dérivable sur cet intervalle.

Posons $u(x) = 2x-1$ et $v(x) = x+3$. Alors $u'(x) = 2$ et $v'(x) = 1$. La formule du quotient donne :
\begin{align*}
g'(x) &= \frac{u'(x)\,v(x) - u(x)\,v'(x)}{[v(x)]^2} \\
&= \frac{2(x+3) - (2x-1)\times 1}{(x+3)^2} \\
&= \frac{2x + 6 - 2x + 1}{(x+3)^2} \\
&= \frac{7}{(x+3)^2}.
\end{align*}
\textbf{Conclusion :} pour tout $x > -3$, $g'(x) = \dfrac{7}{(x+3)^2}$. On remarque au passage que $g'(x) > 0$ : $g$ est strictement croissante sur $]-3\,;\,+\infty[$.

\textbf{3. Dérivée de $h$.}

$h$ est de la forme $u^n$ avec $u(x) = 2x-1$ (dérivable sur $\mathbb{R}$, $u'(x)=2$) et $n = 4$. Donc $h$ est dérivable sur $\mathbb{R}$ et, par la formule $(u^n)' = n\,u'\,u^{n-1}$ :
\[ h'(x) = 4 \times 2 \times (2x-1)^3 = 8(2x-1)^3. \]
\textbf{Conclusion :} pour tout réel $x$, $h'(x) = 8(2x-1)^3$.
\end{exoresolu}

\courssub{Signe de la dérivée et sens de variation}

\begin{propriete}[Théorème fondamental : lien dérivée / variations]
Soit $f$ une fonction dérivable sur un intervalle $I$.
\begin{itemize}
\item Si $f'(x) \geq 0$ sur $I$, alors $f$ est \cle{croissante} sur $I$. Si de plus $f'(x) > 0$ sauf en un nombre fini de points, $f$ est \cle{strictement croissante}.
\item Si $f'(x) \leq 0$ sur $I$, alors $f$ est \cle{décroissante} sur $I$. Si de plus $f'(x) < 0$ sauf en un nombre fini de points, $f$ est \cle{strictement décroissante}.
\item Si $f'(x) = 0$ sur $I$, alors $f$ est \cle{constante} sur $I$.
\end{itemize}
\end{propriete}

\begin{methode}[Du signe de $f'$ au sens de variation de $f$]
\begin{enumerate}
\item \textbf{Calculer $f'(x)$} et le \textbf{factoriser} au maximum (produit de facteurs du premier degré, discriminant pour un trinôme).
\item \textbf{Résoudre l'équation $f'(x) = 0$} : les solutions sont les « frontières » du tableau de signes.
\item \textbf{Dresser le tableau de signes de $f'(x)$} sur l'intervalle d'étude, en utilisant le signe de chaque facteur (règle des signes) ou la position des racines d'un trinôme (signe de $a$ à l'extérieur des racines).
\item \textbf{Conclure par le théorème fondamental} en reliant explicitement le signe de $f'$ au sens de variation de $f$ : « $f'(x) \geq 0$ sur $[a\,;\,b]$, donc $f$ est croissante sur $[a\,;\,b]$ ».
\end{enumerate}
\textbf{Réflexe :} un carré est toujours positif ou nul ; un dénominateur au carré ne change pas le signe. On ne garde souvent que le numérateur à étudier.
\end{methode}

\begin{exoresolu}[Signe de la dérivée et monotonie]
Soit $f$ la fonction définie sur $\mathbb{R}$ par $f(x) = x^3 - 3x + 1$.
\begin{enumerate}
\item Calculer $f'(x)$ et le factoriser.
\item Étudier le signe de $f'(x)$ sur $\mathbb{R}$.
\item En déduire le sens de variation de $f$ sur $\mathbb{R}$.
\end{enumerate}
\tcblower
\textbf{1. Calcul et factorisation de $f'$.}

$f$ est un polynôme, dérivable sur $\mathbb{R}$, et :
\[ f'(x) = 3x^2 - 3 = 3(x^2 - 1) = 3(x-1)(x+1). \]

\textbf{2. Signe de $f'(x)$.}

$f'(x) = 0 \iff x = 1$ ou $x = -1$. Le facteur $3$ est strictement positif ; le signe de $f'$ est donc celui du produit $(x-1)(x+1)$, trinôme dont les racines sont $-1$ et $1$, positif à l'extérieur des racines, négatif entre les racines :

\begin{center}
\begin{tabular}{|c|ccccccc|}
\hline
$x$ & $-\infty$ & & $-1$ & & $1$ & & $+\infty$ \\
\hline
$f'(x)$ & & $+$ & $0$ & $-$ & $0$ & $+$ & \\
\hline
\end{tabular}
\end{center}

\textbf{3. Sens de variation de $f$.}

D'après le théorème de lien entre signe de la dérivée et monotonie :
\begin{itemize}
\item $f'(x) > 0$ sur $]-\infty\,;\,-1[$ et sur $]1\,;\,+\infty[$, donc $f$ est \textbf{strictement croissante} sur chacun de ces intervalles ;
\item $f'(x) < 0$ sur $]-1\,;\,1[$, donc $f$ est \textbf{strictement décroissante} sur $[-1\,;\,1]$.
\end{itemize}
\textbf{Conclusion :} $f$ croît sur $]-\infty\,;\,-1]$, décroît sur $[-1\,;\,1]$, puis croît à nouveau sur $[1\,;\,+\infty[$. Elle admet un maximum local en $-1$ et un minimum local en $1$.
\end{exoresolu}

\courssub{Tableau de variation et optimisation}

\begin{methode}[Construire un tableau de variation irréprochable]
\begin{enumerate}
\item \textbf{Préciser l'ensemble d'étude} $I$ (première ligne du tableau : bornes de $I$ et valeurs interdites éventuelles, marquées d'une double barre $||$).
\item \textbf{Calculer $f'$}, la factoriser, résoudre $f'(x)=0$ et placer les solutions dans la première ligne, \textbf{dans l'ordre croissant}.
\item \textbf{Ligne du signe de $f'(x)$} : signes $+$ ou $-$ entre les zéros, en cohérence avec l'étude de signe.
\item \textbf{Ligne des variations de $f$} : flèche montante $\nearrow$ sous un $+$, flèche descendante $\searrow$ sous un $-$.
\item \textbf{Compléter par les valeurs} : calculer les images des zéros de $f'$ (extremums locaux) et les limites ou valeurs aux bornes de $I$.
\item \textbf{Vérifier la cohérence} : une flèche qui monte doit partir d'une valeur plus petite que celle où elle arrive.
\end{enumerate}
\textbf{Réflexe Bac :} le tableau de variation résume tout ; il doit pouvoir être lu seul, sans le texte.
\end{methode}

\begin{exoresolu}[Tableau de variation complet]
Soit $f$ la fonction définie sur $[0\,;\,6]$ par $f(x) = -x^2 + 4x + 5$.
\begin{enumerate}
\item Étudier le sens de variation de $f$ sur $[0\,;\,6]$.
\item Dresser le tableau de variation complet de $f$ sur $[0\,;\,6]$.
\item En déduire le maximum de $f$ sur $[0\,;\,6]$ et la valeur en laquelle il est atteint.
\end{enumerate}
\tcblower
\textbf{1. Sens de variation.}

$f$ est un polynôme, dérivable sur $[0\,;\,6]$, et $f'(x) = -2x + 4 = -2(x-2)$.

$f'(x) = 0 \iff x = 2$. Comme le coefficient de $x$ dans $f'$ est $-2 < 0$ :
\begin{itemize}
\item pour $x \in [0\,;\,2[$, $f'(x) > 0$ : $f$ est strictement croissante sur $[0\,;\,2]$ ;
\item pour $x \in ]2\,;\,6]$, $f'(x) < 0$ : $f$ est strictement décroissante sur $[2\,;\,6]$.
\end{itemize}

\textbf{2. Tableau de variation.}

On calcule les valeurs utiles :
\[ f(0) = 5 \ ; \qquad f(2) = -4 + 8 + 5 = 9 \ ; \qquad f(6) = -36 + 24 + 5 = -7. \]

\begin{center}
\begin{tabular}{|c|ccccc|}
\hline
$x$ & $0$ & & $2$ & & $6$ \\
\hline
$f'(x)$ & & $+$ & $0$ & $-$ & \\
\hline
$f(x)$ & $5$ & $\nearrow$ & $9$ & $\searrow$ & $-7$ \\
\hline
\end{tabular}
\end{center}

\textbf{3. Maximum de $f$.}

D'après le tableau de variation, $f$ atteint sur $[0\,;\,6]$ son \textbf{maximum en $x = 2$}, et ce maximum vaut $f(2) = 9$.

\textbf{Conclusion :} pour tout $x \in [0\,;\,6]$, $f(x) \leq 9$, avec égalité si et seulement si $x = 2$.
\end{exoresolu}

\begin{methode}[Optimisation : de la situation concrète à la réponse]
\begin{enumerate}
\item \textbf{Choisir l'inconnue} $x$ (la grandeur que l'on peut faire varier) et \textbf{préciser son ensemble de variation} $I$, dicté par les contraintes concrètes (longueurs positives, budget limité...).
\item \textbf{Traduire la quantité à optimiser} en une fonction $f(x)$ : aire, coût, recette, bénéfice. Utiliser les contraintes de l'énoncé pour éliminer les autres variables.
\item \textbf{Étudier les variations de $f$ sur $I$} : dérivée, signe, tableau de variation.
\item \textbf{Lire l'extremum} dans le tableau : maximum ou minimum cherché, valeur de $x$ qui le réalise.
\item \textbf{Répondre à la question posée} par une phrase en langage courant, avec les unités, et vérifier la vraisemblance du résultat.
\end{enumerate}
\textbf{Réflexe :} au Bac, l'ensemble $I$ est souvent donné ou évident ; l'oublier fait perdre des points car un extremum dépend de l'intervalle d'étude.
\end{methode}

\begin{exoresolu}[Optimisation : le champ de bord de route]
Un agriculteur de la région de l'Ouest dispose de \SI{100}{m} de fil barbelé. Il veut clôturer un champ rectangulaire adossé à un mur existant (le mur forme un des grands côtés, qui n'a donc pas besoin de fil). On note $x$ la largeur du champ (côté perpendiculaire au mur), en mètres.
\begin{enumerate}
\item Montrer que l'aire du champ, en mètres carrés, est donnée par $A(x) = x(100 - 2x)$ avec $x \in [0\,;\,50]$.
\item Étudier les variations de $A$ sur $[0\,;\,50]$.
\item En déduire les dimensions du champ d'aire maximale et la valeur de cette aire maximale.
\end{enumerate}
\tcblower
\textbf{1. Expression de l'aire.}

Le fil sert à clôturer deux largeurs $x$ et une seule longueur $L$ (le mur remplace l'autre longueur). La contrainte s'écrit :
\[ 2x + L = 100 \quad \text{donc} \quad L = 100 - 2x. \]
Les longueurs étant positives, $x \geq 0$ et $100 - 2x \geq 0$, c'est-à-dire $x \leq 50$. Donc $x \in [0\,;\,50]$.

L'aire du rectangle est alors :
\[ A(x) = x \times L = x(100 - 2x) = -2x^2 + 100x. \]

\textbf{2. Variations de $A$.}

$A$ est un polynôme, dérivable sur $[0\,;\,50]$, et :
\[ A'(x) = -4x + 100 = -4(x - 25). \]
$A'(x) = 0 \iff x = 25$. Le coefficient de $x$ étant négatif :
\begin{itemize}
\item $A'(x) > 0$ sur $[0\,;\,25[$ : $A$ est strictement croissante sur $[0\,;\,25]$ ;
\item $A'(x) < 0$ sur $]25\,;\,50]$ : $A$ est strictement décroissante sur $[25\,;\,50]$.
\end{itemize}
Valeurs : $A(0) = 0$, $A(25) = 25 \times 50 = 1250$, $A(50) = 0$.

\begin{center}
\begin{tabular}{|c|ccccc|}
\hline
$x$ & $0$ & & $25$ & & $50$ \\
\hline
$A'(x)$ & & $+$ & $0$ & $-$ & \\
\hline
$A(x)$ & $0$ & $\nearrow$ & $1250$ & $\searrow$ & $0$ \\
\hline
\end{tabular}
\end{center}

\textbf{3. Conclusion.}

D'après le tableau, $A$ admet un maximum en $x = 25$. La longueur correspondante est $L = 100 - 2\times 25 = 50$~m.

\textbf{Le champ d'aire maximale a pour dimensions $\SI{25}{m}$ de large et $\SI{50}{m}$ de long, pour une aire maximale de $\SI{1250}{m^2}$.}
\end{exoresolu}

\begin{methode}[Lire et utiliser un tableau de variation]
\begin{enumerate}
\item \textbf{Encadrer $f(x)$} : sur un intervalle où $f$ est monotone, les images sont comprises entre les valeurs aux bornes (dans le même ordre si $f$ croît, en ordre inversé si $f$ décroît).
\item \textbf{Comparer deux images} : si $f$ est croissante sur $I$ et $a < b$ dans $I$, alors $f(a) < f(b)$ ; si $f$ est décroissante, $f(a) > f(b)$.
\item \textbf{Résoudre $f(x) = k$} : localiser $k$ par rapport aux extremums du tableau ; sur chaque intervalle de monotonie stricte, l'équation a au plus une solution (unicité par stricte monotonie, existence par les valeurs aux bornes).
\item \textbf{Résoudre $f(x) \geq k$} : repérer sur quelles portions du tableau les valeurs dépassent $k$.
\end{enumerate}
\textbf{Réflexe :} toujours justifier par la monotonie : « $f$ est strictement croissante sur $I$, donc... ».
\end{methode}

\begin{exoresolu}[Exploitation d'un tableau de variation]
Soit $g$ la fonction définie sur $[-2\,;\,4]$ par $g(x) = x^2 - 2x - 3$.
\begin{enumerate}
\item Dresser le tableau de variation de $g$ sur $[-2\,;\,4]$.
\item Comparer $g(0,4)$ et $g(0,9)$, puis $g(2,5)$ et $g(3,8)$. Justifier.
\item Montrer que pour tout $x \in [-2\,;\,4]$, $g(x) \geq -4$.
\item Déterminer le nombre de solutions de l'équation $g(x) = 0$ sur $[-2\,;\,4]$ et les donner.
\end{enumerate}
\tcblower
\textbf{1. Tableau de variation.}

$g$ est un polynôme, dérivable sur $[-2\,;\,4]$ : $g'(x) = 2x - 2 = 2(x-1)$. Donc $g'(x) < 0$ sur $[-2\,;\,1[$ et $g'(x) > 0$ sur $]1\,;\,4]$.

Valeurs : $g(-2) = 4 + 4 - 3 = 5$ ; $g(1) = 1 - 2 - 3 = -4$ ; $g(4) = 16 - 8 - 3 = 5$.

\begin{center}
\begin{tabular}{|c|ccccc|}
\hline
$x$ & $-2$ & & $1$ & & $4$ \\
\hline
$g'(x)$ & & $-$ & $0$ & $+$ & \\
\hline
$g(x)$ & $5$ & $\searrow$ & $-4$ & $\nearrow$ & $5$ \\
\hline
\end{tabular}
\end{center}

\textbf{2. Comparaisons.}

Sur $[-2\,;\,1]$, $g$ est strictement décroissante et $0,4 < 0,9$, donc $g(0,4) > g(0,9)$.

Sur $[1\,;\,4]$, $g$ est strictement croissante et $2,5 < 3,8$, donc $g(2,5) < g(3,8)$.

\textbf{3. Minoration.}

D'après le tableau de variation, le minimum de $g$ sur $[-2\,;\,4]$ est $g(1) = -4$. Par définition d'un minimum, pour tout $x \in [-2\,;\,4]$, $g(x) \geq -4$.

\textbf{4. Équation $g(x) = 0$.}

Sur $[-2\,;\,1]$, $g$ est strictement décroissante de $5$ vers $-4$ : elle prend donc exactement une fois la valeur $0$. Sur $[1\,;\,4]$, $g$ est strictement croissante de $-4$ vers $5$ : elle prend aussi exactement une fois la valeur $0$. L'équation a donc \textbf{exactement deux solutions}.

Résolution exacte : $x^2 - 2x - 3 = 0$. Le discriminant est $\Delta = 4 + 12 = 16$, d'où $x = \dfrac{2 - 4}{2} = -1$ ou $x = \dfrac{2+4}{2} = 3$. Les deux solutions, $-1$ et $3$, appartiennent bien à $[-2\,;\,4]$.

\textbf{Conclusion :} $g(x) = 0$ admet exactement deux solutions sur $[-2\,;\,4]$ : $x = -1$ et $x = 3$.
\end{exoresolu}

\begin{methode}[Relier tangente, nombre dérivé et variations]
\begin{enumerate}
\item \textbf{Le nombre dérivé $f'(a)$ est le coefficient directeur} de la tangente à la courbe de $f$ au point d'abscisse $a$.
\item \textbf{Équation de la tangente} en $a$ : $y = f'(a)(x - a) + f(a)$. Calculer séparément $f(a)$ et $f'(a)$ avant de substituer.
\item \textbf{Lecture graphique} : une tangente horizontale signifie $f'(a) = 0$ (extremum local probable) ; une tangente de pente positive signifie $f$ croissante au voisinage de $a$.
\item \textbf{Cohérence} : vérifier que le signe de la pente de la tangente est compatible avec le sens de variation donné par le tableau.
\end{enumerate}
\end{methode}

\begin{exoresolu}[Tangente et sens de variation]
Soit $f$ la fonction définie sur $\mathbb{R}$ par $f(x) = x^2 - 4x + 7$, de courbe représentative $\mathcal{C}$ dans un repère.
\begin{enumerate}
\item Déterminer l'équation réduite de la tangente $T$ à $\mathcal{C}$ au point $A$ d'abscisse $3$.
\item Montrer que $\mathcal{C}$ admet une tangente horizontale en un point $B$ dont on précisera les coordonnées.
\item Que représente le point $B$ pour la fonction $f$ ? Justifier à l'aide de la dérivée.
\end{enumerate}
\tcblower
\textbf{1. Équation de la tangente en $A$.}

$f$ est un polynôme, dérivable sur $\mathbb{R}$ : $f'(x) = 2x - 4$.

On calcule : $f(3) = 9 - 12 + 7 = 4$ et $f'(3) = 6 - 4 = 2$.

L'équation de la tangente en $a = 3$ est $y = f'(3)(x-3) + f(3)$, soit :
\[ y = 2(x-3) + 4 = 2x - 2. \]
\textbf{Conclusion :} $T$ a pour équation réduite $y = 2x - 2$. La pente $2 > 0$ est cohérente avec le fait que $f$ est croissante au voisinage de $3$ (car $3 > 2$, voir question 3).

\textbf{2. Tangente horizontale.}

Une tangente est horizontale si et seulement si son coefficient directeur est nul, c'est-à-dire $f'(x) = 0$ :
\[ 2x - 4 = 0 \iff x = 2. \]
Or $f(2) = 4 - 8 + 7 = 3$. Donc $\mathcal{C}$ admet une tangente horizontale au point $B(2\,;\,3)$.

\textbf{3. Nature du point $B$.}

$f'(x) = 2(x-2)$ : $f'(x) < 0$ pour $x < 2$ et $f'(x) > 0$ pour $x > 2$. Donc $f$ est strictement décroissante sur $]-\infty\,;\,2]$ puis strictement croissante sur $[2\,;\,+\infty[$.

\textbf{Conclusion :} $f$ admet en $x = 2$ un \textbf{minimum}, égal à $f(2) = 3$. Le point $B(2\,;\,3)$ est le point le plus bas de la courbe $\mathcal{C}$ : c'est le sommet de la parabole.
\end{exoresolu}

\begin{popfigure}
\begin{tikzpicture}[scale=0.7, >=Stealth]
\def\f{(\x)^2 - 4*\x + 7}
\def\tangA{2*\x - 2}
\def\tangB{3}
\draw[->] (-1,0) -- (5,0) node[right] {$x$};
\draw[->] (0,-1) -- (0,8) node[above] {$y$};
\draw[domain=-0.5:4.5,smooth,variable=\x,popcurve,thick] plot ({\x},{\f});
\draw[domain=0.5:4,smooth,variable=\x,poporange,dashed,thick] plot ({\x},{\tangA});
\draw[domain=0.5:3.5,smooth,variable=\x,popgreen,dashed,thick] plot ({\x},{\tangB});
\fill[popblue] (3,4) circle (2pt) node[above right] {$A(3;4)$};
\fill[poppurple] (2,3) circle (2pt) node[below left] {$B(2;3)$};
\draw[popmuted] (3,0) node[below] {$3$} -- (3,4);
\draw[popmuted] (2,0) node[below] {$2$} -- (2,3);
\draw[popmuted] (0,3) node[left] {$3$} -- (2,3);
\draw[popmuted] (0,4) node[left] {$4$} -- (3,4);
\node[poporange, right] at (4,5.5) {Tangente en $A$ : $y=2x-2$};
\node[popgreen, right] at (3.5,2.5) {Tangente horizontale en $B$};
\node[popblue, right] at (4.5,7) {$\mathcal{C}_f : y=x^2-4x+7$};
\end{tikzpicture}
\legende{Illustration de l'exercice : tangente en $A(3;4)$ et tangente horizontale au minimum $B(2;3)$.}
\end{popfigure}

\begin{attention}[Les 5 erreurs qui coûtent des points au Bac]
\begin{enumerate}
\item \textbf{Dériver un produit ou un quotient comme un produit ou quotient de dérivées.} $(uv)' \neq u'v'$ et $\left(\dfrac{u}{v}\right)' \neq \dfrac{u'}{v'}$. Les bonnes formules sont $(uv)' = u'v + uv'$ et $\left(\dfrac{u}{v}\right)' = \dfrac{u'v - uv'}{v^2}$, avec le carré au dénominateur.
\item \textbf{Oublier le facteur $u'$ dans une fonction composée.} La dérivée de $(2x-1)^4$ n'est pas $4(2x-1)^3$ mais $8(2x-1)^3$ : on multiplie toujours par la dérivée de « l'intérieur ».
\item \textbf{Conclure sur les variations sans justifier par le signe de $f'$.} Écrire « $f$ est croissante » sans écrire « car $f'(x) \geq 0$ sur l'intervalle » est une réponse non justifiée : le lien signe de $f'$ $\Rightarrow$ monotonie de $f$ doit apparaître explicitement.
\item \textbf{Se tromper dans le tableau de variation} : valeurs de $f'$ dans la ligne de $f$, flèche montante sous un signe $-$, ou images non calculées (un extremum sans sa valeur est incomplet). Toujours calculer $f$ aux zéros de $f'$ et aux bornes.
\item \textbf{Oublier l'ensemble d'étude dans un problème d'optimisation.} Un extremum se cherche sur un intervalle précis (souvent imposé par le concret : $x \geq 0$, $x \leq 50$...). Chercher le maximum sur $\mathbb{R}$ au lieu de l'intervalle de l'énoncé peut donner une réponse fausse, et ne pas répondre en langage courant avec les unités fait perdre le point de conclusion.
\end{enumerate}
\end{attention}

\newpage

\coursec{Exercices}

\coursec{MA03 : Dérivées et sens de variation}

\courssub{Rappels : nombre dérivé et tangente}

\begin{definition}[Nombre dérivé en un point]
Soit $f$ une fonction définie sur un intervalle $I$ et $a \in I$. Si la limite 
\[ \lim_{h \to 0} \frac{f(a+h)-f(a)}{h} \]
existe et est finie, on dit que $f$ est \cle{dérivable} en $a$. Cette limite est appelée \cle{nombre dérivé} de $f$ en $a$ et se note $f'(a)$.
\end{definition}

\begin{propriete}[Interprétation géométrique]
Si $f$ est dérivable en $a$, la courbe $\mathcal{C}_f$ admet en le point $A(a\,;\,f(a))$ une \cle{tangente} non verticale. Son équation cartésienne est :
\[ y = f'(a)(x-a) + f(a). \]
Le coefficient directeur de cette tangente est $f'(a)$.
\end{propriete}

\begin{exemplebox}[Calcul d'un nombre dérivé et tangente]
Soit $f(x)=x^2$ et $a=2$.
\begin{align*}
f'(2) &= \lim_{h \to 0} \frac{(2+h)^2-4}{h} = \lim_{h \to 0} \frac{4+4h+h^2-4}{h} = \lim_{h \to 0} (4+h) = 4.
\end{align*}
La tangente en $A(2;4)$ a pour équation $y = 4(x-2)+4 = 4x-4$.
\end{exemplebox}

\begin{attention}[Dérivabilité et continuité]
Si $f$ est dérivable en $a$, alors $f$ est continue en $a$. La réciproque est fausse : $f(x)=|x|$ est continue en $0$ mais non dérivable en $0$ (tangente verticale / angle).
\end{attention}

\courssub{Dérivées des fonctions usuelles}

\begin{propriete}[Tableau des dérivées usuelles -- \textbf{À connaître par cœur}]
\begin{center}
\begin{tabularx}{\linewidth}{|l|c|c|X|}
\hline
\rowcolor{popblueL}
\textbf{Fonction $f(x)$} & \textbf{Domaine $D_f$} & \textbf{Dérivée $f'(x)$} & \textbf{Domaine de dérivabilité} \\
\hline
$k$ (constante) & $\mathbb{R}$ & $0$ & $\mathbb{R}$ \\
$x^n$ ($n \in \mathbb{N}$) & $\mathbb{R}$ & $n x^{n-1}$ & $\mathbb{R}$ \\
$x^\alpha$ ($\alpha \in \mathbb{R}$) & $]0;+\infty[$ & $\alpha x^{\alpha-1}$ & $]0;+\infty[$ \\
$\dfrac{1}{x} = x^{-1}$ & $\mathbb{R}^*$ & $-\dfrac{1}{x^2}$ & $\mathbb{R}^*$ \\
$\sqrt{x} = x^{1/2}$ & $[0;+\infty[$ & $\dfrac{1}{2\sqrt{x}}$ & $]0;+\infty[$ \\
$\ln x$ & $]0;+\infty[$ & $\dfrac{1}{x}$ & $]0;+\infty[$ \\
$e^x$ & $\mathbb{R}$ & $e^x$ & $\mathbb{R}$ \\
$\cos x$ & $\mathbb{R}$ & $-\sin x$ & $\mathbb{R}$ \\
$\sin x$ & $\mathbb{R}$ & $\cos x$ & $\mathbb{R}$ \\
\hline
\end{tabularx}
\end{center}
\end{propriete}

\begin{aretenir}[Dérivabilité aux bornes]
Pour $\sqrt{x}$ et $\ln x$, la dérivée n'existe pas à la borne gauche du domaine (pente infinie). On étudie la dérivabilité sur l'intérieur de l'intervalle.
\end{aretenir}

\begin{exemplebox}[Application directe]
$f(x) = 3x^4 - 5\sqrt{x} + \dfrac{2}{x}$ sur $]0;+\infty[$.
$f'(x) = 12x^3 - \dfrac{5}{2\sqrt{x}} - \dfrac{2}{x^2}$.
\end{exemplebox}

\courssub{Opérations sur les dérivées}

\begin{propriete}[Règles de dérivation -- \textbf{Démonstration exigible au Bac pour la somme et le produit}]
Soient $u$ et $v$ dérivables sur un intervalle $I$.
\begin{itemize}
\item \textbf{Somme :} $(u+v)' = u' + v'$
\item \textbf{Multiple par une constante :} $(ku)' = k u'$ ($k \in \mathbb{R}$)
\item \textbf{Produit :} $(uv)' = u'v + uv'$
\item \textbf{Quotient :} $\left(\dfrac{u}{v}\right)' = \dfrac{u'v - uv'}{v^2}$ \quad ($v(x) \neq 0$ sur $I$)
\end{itemize}
\end{propriete}

\begin{methode}[Calculer la dérivée d'une combinaison de fonctions usuelles]
\begin{enumerate}
\item Identifier la structure : somme, produit, quotient, ou fonction simple.
\item Appliquer la règle correspondante.
\item Simplifier l'expression obtenue (factoriser, mettre au même dénominateur).
\item Préciser le domaine de dérivabilité (intersection des domaines, en excluant les zéros du dénominateur pour un quotient).
\end{enumerate}
\end{methode}

\begin{exemplebox}[Dérivée d'un quotient]
$f(x) = \dfrac{x^2+1}{x-1}$ sur $\mathbb{R}\setminus\{1\}$.
$f'(x) = \dfrac{2x(x-1) - (x^2+1)\cdot 1}{(x-1)^2} = \dfrac{2x^2-2x-x^2-1}{(x-1)^2} = \dfrac{x^2-2x-1}{(x-1)^2}$.
\end{exemplebox}

\courssub{Dérivée d'une fonction composée}

\begin{propriete}[Formule de la chaîne -- \textbf{Démonstration non exigible, utilisation obligatoire}]
Soient $u$ dérivable sur $I$ et $v$ dérivable sur $u(I)$. La fonction composée $f = v \circ u$ (définie par $f(x)=v(u(x))$) est dérivable sur $I$ et :
\[ f'(x) = u'(x) \times v'(u(x)). \]
En notation de Leibniz : $\dfrac{df}{dx} = \dfrac{dv}{du} \times \dfrac{du}{dx}$.
\end{propriete}

\begin{aretenir}[Cas particuliers fréquents -- \textbf{Formules "boîtes à outils"}]
Si $u$ est dérivable et $u(x) > 0$ si nécessaire :
\begin{align*}
(u^n)' &= n u' u^{n-1} & (\sqrt{u})' &= \frac{u'}{2\sqrt{u}} \\
\left(\frac{1}{u}\right)' &= -\frac{u'}{u^2} & (\ln u)' &= \frac{u'}{u} \\
(e^u)' &= u' e^u & (\cos u)' &= -u' \sin u \\
(\sin u)' &= u' \cos u &
\end{align*}
\end{aretenir}

\begin{methode}[Dériver une fonction composée]
\begin{enumerate}
\item Identifier la fonction "extérieure" $v$ et la fonction "intérieure" $u$ (souvent ce qui est entre parenthèses ou sous le signe racine/ln).
\item Calculer $u'$ et $v'$.
\item Appliquer $f' = u' \times v'(u)$.
\item Simplifier.
\end{enumerate}
\end{methode}

\begin{exemplebox}[Fonction composée]
$f(x) = (3x^2-2x+1)^5$ sur $\mathbb{R}$.
$u(x)=3x^2-2x+1$, $v(X)=X^5$.
$u'(x)=6x-2$, $v'(X)=5X^4$.
$f'(x) = (6x-2) \times 5(3x^2-2x+1)^4 = 10(3x-1)(3x^2-2x+1)^4$.
\end{exemplebox}

\begin{attention}[Piège classique]
Ne pas oublier le facteur $u'(x)$ (la dérivée de l'intérieur).
\textbf{Faux :} $((3x+1)^2)' = 2(3x+1)$.
\textbf{Vrai :} $((3x+1)^2)' = 3 \times 2(3x+1) = 6(3x+1)$.
\end{attention}

\courssub{Signe de la dérivée et sens de variation}

\begin{propriete}[Théorème fondamental -- \textbf{Démonstration exigible au Bac (schéma de principe)}]
Soit $f$ continue sur un intervalle $I$ et dérivable sur l'intérieur de $I$.
\begin{itemize}
\item Si $f'(x) \ge 0$ sur $I$, alors $f$ est \cle{croissante} sur $I$.
\item Si $f'(x) \le 0$ sur $I$, alors $f$ est \cle{décroissante} sur $I$.
\item Si $f'(x) = 0$ sur $I$, alors $f$ est \cle{constante} sur $I$.
\end{itemize}
Réciproquement, si $f$ est strictement monotone, le signe de $f'$ est constant (mais peut s'annuler en des points isolés, ex: $x^3$ en $0$).
\end{propriete}

\begin{methode}[Étudier le signe de $f'$ et en déduire les variations]
\begin{enumerate}
\item Calculer $f'(x)$ et le factoriser au maximum.
\item Résoudre $f'(x)=0$ pour trouver les valeurs critiques.
\item Dresser un \cle{tableau de signes} de $f'$ sur l'intervalle d'étude.
\item En déduire le sens de variation de $f$ (flèches $\nearrow$ ou $\searrow$).
\end{enumerate}
\end{methode}

\begin{exemplebox}[Variations de $f(x)=x^3-3x$ sur $\mathbb{R}$]
$f'(x)=3x^2-3=3(x-1)(x+1)$.
Zéros : $-1$ et $1$.
Signe de $f'$ : $+$ sur $]-\infty;-1]$, $-$ sur $[-1;1]$, $+$ sur $[1;+\infty[$.
$f$ croissante sur $]-\infty;-1]$, décroissante sur $[-1;1]$, croissante sur $[1;+\infty[$.
Maximum local $f(-1)=2$, minimum local $f(1)=-2$.
\end{exemplebox}

\courssub{Tableau de variation et optimisation}

\begin{methode}[Dresser un tableau de variation complet]
\begin{enumerate}
\item Première ligne : variable $x$ (domaine d'étude), valeurs critiques, bornes.
\item Deuxième ligne : signe de $f'$ ($+$, $-$, $0$, $\barre$ pour non défini).
\item Troisième ligne : variations de $f$ ($\nearrow$, $\searrow$) et \textbf{valeurs exactes} aux extrema et aux bornes ($f(x)$).
\item Indiquer les limites aux bornes ouvertes si nécessaire ($+\infty$, $-\infty$, limite finie).
\end{enumerate}
\end{methode}

\begin{propriete}[{Optimisation sur un intervalle fermé $[a;b]$}]
Si $f$ est continue sur $[a;b]$ et dérivable sur $]a;b[$, ses extrema globaux (max/min) sont atteints :
\begin{itemize}
\item soit aux bornes $a$ ou $b$,
\item soit en un point $c \in ]a;b[$ où $f'(c)=0$ (point critique).
\end{itemize}
Il suffit de comparer les valeurs $f(a)$, $f(b)$ et $f(c_i)$.
\end{propriete}

\begin{exemplebox}[Problème d'optimisation : volume d'une boîte]
On découpe 4 carrés de côté $x$ dans une feuille $20\times 30$ cm pour faire une boîte ouverte.
$V(x) = x(20-2x)(30-2x) = 4x^3-100x^2+600x$ sur $[0;10]$.
$V'(x)=12x^2-200x+600=4(3x^2-50x+150)$.
Racines : $x = \frac{50 \pm \sqrt{2500-1800}}{6} = \frac{50 \pm \sqrt{700}}{6} \approx 2,55$ et $13,1$ (hors intervalle).
$V(0)=0$, $V(10)=0$, $V(2,55) \approx 513$ cm$^3$. Maximum pour $x \approx 2,55$ cm.
\end{exemplebox}

\begin{savaistu}[Contexte camerounais : Optimisation des coûts]
À Douala, une usine de transformation de cacao cherche à minimiser son coût moyen $C_M(x) = \frac{C(x)}{x}$ où $C(x)$ est le coût total pour $x$ tonnes. L'étude de la dérivée de $C_M$ permet de trouver la production optimale (point mort économique), évitant le surstockage ou la sous-production. C'est une application directe du programme en économie réelle.
\end{savaistu}

\courssub{Je vérifie mes connaissances}

\begin{exercice}[QCM : dérivées usuelles]
Pour chaque question, une seule réponse est exacte. Entoure-la.
\begin{enumerate}
\item La dérivée de la fonction $f$ définie sur $\mathbb{R}$ par $f(x) = x^5$ est :
\begin{enumerate}
\item[a)] $f'(x) = 4x^4$ \qquad b) $f'(x) = 5x^4$ \qquad c) $f'(x) = 5x^5$
\end{enumerate}
\item La dérivée de la fonction $g$ définie sur $]0\,;\,+\infty[$ par $g(x) = \sqrt{x}$ est :
\begin{enumerate}
\item[a)] $g'(x) = \dfrac{1}{\sqrt{x}}$ \qquad b) $g'(x) = \dfrac{1}{2\sqrt{x}}$ \qquad c) $g'(x) = 2\sqrt{x}$
\end{enumerate}
\item Si $h(x) = \dfrac{1}{x}$ sur $\mathbb{R}^*$, alors $h'(x)$ est :
\begin{enumerate}
\item[a)] $\dfrac{1}{x^2}$ \qquad b) $-\dfrac{1}{x^2}$ \qquad c) $\ln|x|$
\end{enumerate}
\item Si $f'(x) < 0$ pour tout $x$ d'un intervalle $I$, alors $f$ est :
\begin{enumerate}
\item[a)] croissante sur $I$ \qquad b) décroissante sur $I$ \qquad c) constante sur $I$
\end{enumerate}
\end{enumerate}
\end{exercice}

\begin{exercice}[Vrai ou faux, en justifiant]
Pour chaque affirmation, dire si elle est vraie ou fausse, puis justifier la réponse.
\begin{enumerate}
\item La dérivée d'une fonction constante est la fonction nulle.
\item Si $f(x) = (3x+1)^2$, alors $f'(x) = 2(3x+1)$.
\item Si $f'(2) = 0$, alors la courbe de $f$ admet au point d'abscisse $2$ une tangente horizontale.
\item Une fonction décroissante sur $\mathbb{R}$ a une dérivée strictement négative en tout point.
\end{enumerate}
\end{exercice}

\begin{exercice}[Compléter les formules]
Recopie et complète les égalités suivantes, où $u$ et $v$ sont des fonctions dérivables :
\begin{enumerate}
\item $(u + v)' = \ldots$ \qquad $(ku)' = \ldots$ (où $k$ est un réel)
\item $(uv)' = \ldots$ \qquad $\left(\dfrac{u}{v}\right)' = \ldots$ (avec $v \neq 0$)
\item $(u^n)' = \ldots$ \qquad $(\sqrt{u})' = \ldots$
\end{enumerate}
\end{exercice}

\begin{exercice}[Calculs directs]
Calcule la dérivée de chacune des fonctions suivantes, définies sur $\mathbb{R}$ (préciser le domaine si nécessaire) :
\begin{enumerate}
\item $f(x) = 4x^3 - 5x^2 + 7x - 2$
\item $g(x) = (2x - 1)(x + 3)$
\item $h(x) = \dfrac{3x - 1}{x + 2}$
\item $k(x) = x^2 + \dfrac{1}{x}$
\end{enumerate}
\end{exercice}

\courssub{Je m'entraîne}

\begin{exercice}[Dérivées de fonctions composées]
Pour chacune des fonctions suivantes, précise le domaine de dérivabilité, puis calcule la dérivée :
\begin{enumerate}
\item $f(x) = (2x^2 - 3x + 1)^5$
\item $g(x) = \dfrac{x + 1}{x^2 + 2}$
\item $h(x) = \sqrt{3x - 6}$
\end{enumerate}
\end{exercice}

\begin{exercice}[Variations d'une fonction polynôme]
Soit $f$ la fonction définie sur $\mathbb{R}$ par $f(x) = x^3 - 6x^2 + 9x + 1$.
\begin{enumerate}
\item Calcule $f'(x)$ et vérifie que $f'(x) = 3(x - 1)(x - 3)$.
\item Étudie le signe de $f'(x)$ sur $\mathbb{R}$.
\item Dresse le tableau de variation complet de $f$ sur $\mathbb{R}$ (avec les valeurs des extrema).
\end{enumerate}
\end{exercice}

\begin{exercice}[Fonction homographique]
Soit $f$ la fonction définie par $f(x) = \dfrac{x - 2}{x + 1}$.
\begin{enumerate}
\item Détermine le domaine de définition $D_f$ de $f$.
\item Calcule $f'(x)$ et étudie son signe sur $D_f$.
\item Dresse le tableau de variation de $f$ sur $D_f$.
\item En utilisant le tableau de variation, détermine le nombre de solutions de l'équation $f(x) = 3$ sur $D_f$.
\end{enumerate}
\end{exercice}

\begin{exercice}[Tangente et variations]
Soit $f$ la fonction définie sur $\left[-\dfrac{1}{2}\,;\,+\infty\right[$ par $f(x) = \sqrt{2x + 1}$.
\begin{enumerate}
\item Justifie que $f$ est dérivable sur $\left]-\dfrac{1}{2}\,;\,+\infty\right[$ et calcule $f'(x)$.
\item Détermine une équation de la tangente à la courbe de $f$ au point d'abscisse $4$.
\item Étudie le signe de $f'(x)$ et dresse le tableau de variation de $f$.
\end{enumerate}
\end{exercice}

\courssub{Je me prépare au Bac}

\begin{exercice}[Bac blanc : coût moyen d'une entreprise de Douala]
Une entreprise de Douala fabrique des objets artisanaux. Pour une production de $x$ centaines d'objets (avec $x \in [1\,;\,10]$), le coût total de production, exprimé en milliers de FCFA, est donné par :
\[ C(x) = x^3 - 12x^2 + 60x. \]
Le \cle{coût moyen} de production, en milliers de FCFA par centaine d'objets, est défini par $C_M(x) = \dfrac{C(x)}{x}$.
\begin{enumerate}
\item \begin{enumerate}
\item Calcule le coût total et le coût moyen pour une production de $200$ objets.
\item Montre que pour tout $x \in [1\,;\,10]$, $C_M(x) = x^2 - 12x + 60$.
\end{enumerate}
\item \begin{enumerate}
\item Calcule $C_M'(x)$ et étudie son signe sur $[1\,;\,10]$.
\item Dresse le tableau de variation de la fonction $C_M$ sur $[1\,;\,10]$.
\end{enumerate}
\item \begin{enumerate}
\item Déduis-en la production (en nombre d'objets) qui minimise le coût moyen.
\item Quel est alors le coût moyen minimal, en FCFA par objet ?
\end{enumerate}
\end{enumerate}
\end{exercice}

\begin{exercice}[Bac blanc : une inégalité classique]
On considère la fonction $f$ définie sur $]0\,;\,+\infty[$ par :
\[ f(x) = x + \dfrac{1}{x}. \]
\begin{enumerate}
\item Calcule $f'(x)$ et montre que pour tout $x > 0$, $f'(x) = \dfrac{(x-1)(x+1)}{x^2}$.
\item Étudie le signe de $f'(x)$ sur $]0\,;\,+\infty[$.
\item Dresse le tableau de variation de $f$ sur $]0\,;\,+\infty[$ et calcule $f(1)$.
\item Déduis-en que pour tout réel $x > 0$, $x + \dfrac{1}{x} \geq 2$. Dans quel cas y a-t-il égalité ?
\item \textbf{Application.} Un commerçant de Yaoundé veut clôturer un terrain rectangulaire d'aire $\SI{100}{m^2}$. On note $x$ la longueur (en m) d'un côté. Exprime le périmètre du terrain en fonction de $x$, puis détermine les dimensions du terrain de périmètre minimal.
\end{enumerate}
\end{exercice}

\begin{exercice}[Bac blanc : étude complète d'une fonction]
Soit $f$ la fonction définie sur $\mathbb{R}$ par :
\[ f(x) = \dfrac{1}{3}x^3 - x^2 - 3x + 2 \]
et $(\mathcal{C})$ sa courbe représentative dans un repère orthogonal.
\begin{enumerate}
\item Calcule $f'(x)$ et vérifie que $f'(x) = (x + 1)(x - 3)$.
\item Étudie le signe de $f'(x)$ sur $\mathbb{R}$.
\item Dresse le tableau de variation complet de $f$ sur $\mathbb{R}$, en calculant les valeurs des extrema.
\item Détermine une équation de la tangente à $(\mathcal{C})$ au point d'abscisse $0$.
\item \begin{enumerate}
\item Montre que l'équation $f(x) = 0$ admet une unique solution $\alpha$ dans l'intervalle $[4\,;\,5]$.
\item Donne un encadrement de $\alpha$ d'amplitude $0{,}5$.
\end{enumerate}
\item Trace l'allure de la courbe $(\mathcal{C})$ en plaçant les tangentes horizontales et la tangente au point d'abscisse $0$.
\end{enumerate}
\end{exercice}



\newpage

\coursec{Corrigés des exercices}

\begin{corrige}[Exercice 1 — QCM : dérivées usuelles]
\begin{enumerate}
\item \textbf{Réponse b)} $f'(x) = 5x^4$. En effet, pour $f(x)=x^n$ avec $n=5$, la formule $(x^n)' = nx^{n-1}$ donne $f'(x)=5x^{5-1}=5x^4$. La réponse a) correspond à une erreur sur l'exposant, la réponse c) oublie d'abaisser l'exposant.
\item \textbf{Réponse b)} $g'(x) = \dfrac{1}{2\sqrt{x}}$. En effet, $\sqrt{x}=x^{1/2}$, donc $g'(x)=\dfrac{1}{2}x^{-1/2}=\dfrac{1}{2\sqrt{x}}$. La réponse a) oublie le facteur $\frac{1}{2}$ ; la réponse c) confond avec une primitive.
\item \textbf{Réponse b)} $h'(x) = -\dfrac{1}{x^2}$. En effet, $h(x)=x^{-1}$, donc $h'(x)=-1\cdot x^{-2}=-\dfrac{1}{x^2}$. \textbf{Attention au signe moins !}
\item \textbf{Réponse b)} $f$ est décroissante sur $I$. C'est le théorème fondamental : $f'(x)<0$ sur $I$ implique $f$ strictement décroissante sur $I$.
\end{enumerate}
\end{corrige}

\begin{corrige}[Exercice 2 — Vrai ou faux, en justifiant]
\begin{enumerate}
\item \textbf{Vrai.} Si $f(x)=k$ (constante), alors pour tout $a$ et tout $h\neq 0$, le taux d'accroissement vaut $\dfrac{f(a+h)-f(a)}{h}=\dfrac{k-k}{h}=0$, qui tend vers $0$. Donc $f'(x)=0$ pour tout $x$.
\item \textbf{Faux.} C'est le piège classique de la fonction composée : on a oublié la dérivée de l'intérieur. Avec $u(x)=3x+1$ (donc $u'(x)=3$) et $v(X)=X^2$ (donc $v'(X)=2X$) :
\[ f'(x) = u'(x)\times v'(u(x)) = 3\times 2(3x+1) = 6(3x+1). \]
Vérification par développement : $f(x)=9x^2+6x+1$, donc $f'(x)=18x+6=6(3x+1)$. Confirmé.
\item \textbf{Vrai.} L'équation de la tangente au point d'abscisse $2$ est $y=f'(2)(x-2)+f(2)$. Si $f'(2)=0$, cette équation devient $y=f(2)$ : c'est une droite horizontale (coefficient directeur nul).
\item \textbf{Faux.} Contre-exemple : $f(x)=-x^3$ est strictement décroissante sur $\mathbb{R}$, pourtant $f'(x)=-3x^2$ s'annule en $0$. Le bon énoncé est : une fonction décroissante sur un intervalle a une dérivée \emph{négative ou nulle} ($f'\le 0$), la dérivée pouvant s'annuler en des points isolés.
\end{enumerate}
\end{corrige}

\begin{corrige}[Exercice 3 — Compléter les formules]
\begin{enumerate}
\item $(u+v)' = u'+v'$ \quad ; \quad $(ku)' = ku'$ (où $k\in\mathbb{R}$).
\item $(uv)' = u'v+uv'$ \quad ; \quad $\left(\dfrac{u}{v}\right)' = \dfrac{u'v-uv'}{v^2}$ (avec $v(x)\neq 0$).
\item $(u^n)' = n\,u'\,u^{n-1}$ \quad ; \quad $(\sqrt{u})' = \dfrac{u'}{2\sqrt{u}}$ (avec $u(x)>0$).
\end{enumerate}
\textbf{Point de vigilance :} dans la formule du quotient, c'est $u'v-uv'$ au numérateur (dans cet ordre : dérivée du haut fois le bas, \emph{moins} le haut fois dérivée du bas), et le dénominateur est $v^2$, pas $v$.
\end{corrige}

\begin{corrige}[Exercice 4 — Calculs directs]
\begin{enumerate}
\item $f$ est un polynôme, dérivable sur $\mathbb{R}$. On dérive terme à terme :
\[ f'(x) = 4\times 3x^2 - 5\times 2x + 7\times 1 - 0 = 12x^2 - 10x + 7. \]
\item $g$ est définie et dérivable sur $\mathbb{R}$. Deux méthodes :
\begin{itemize}
\item \emph{En développant :} $g(x)=2x^2+6x-x-3=2x^2+5x-3$, donc $g'(x)=4x+5$.
\item \emph{Avec le produit :} $u=2x-1$, $v=x+3$ ; $u'=2$, $v'=1$ :
\[ g'(x)=2(x+3)+(2x-1)\times 1 = 2x+6+2x-1 = 4x+5. \]
\end{itemize}
Les deux méthodes donnent le même résultat : $g'(x)=4x+5$.
\item $h$ est définie et dérivable sur $\mathbb{R}\setminus\{-2\}$. Avec $u=3x-1$, $v=x+2$ ($u'=3$, $v'=1$) :
\[ h'(x) = \frac{3(x+2)-(3x-1)\times 1}{(x+2)^2} = \frac{3x+6-3x+1}{(x+2)^2} = \frac{7}{(x+2)^2}. \]
\item $k$ est définie et dérivable sur $\mathbb{R}^*$ (à cause de $\frac{1}{x}$) :
\[ k'(x) = 2x - \frac{1}{x^2} = \frac{2x^3-1}{x^2}. \]
\end{enumerate}
\end{corrige}

\begin{corrige}[Exercice 5 — Dérivées de fonctions composées]
\begin{enumerate}
\item $f(x)=(2x^2-3x+1)^5$ est un polynôme (composée de deux polynômes), donc dérivable sur $\mathbb{R}$.
On pose $u(x)=2x^2-3x+1$, $u'(x)=4x-3$, et on applique $(u^n)'=nu'u^{n-1}$ avec $n=5$ :
\[ f'(x) = 5(4x-3)(2x^2-3x+1)^4. \]
\item $g(x)=\dfrac{x+1}{x^2+2}$. Le dénominateur $x^2+2\geq 2>0$ ne s'annule jamais : $g$ est définie et dérivable sur $\mathbb{R}$.
Avec $u=x+1$ ($u'=1$) et $v=x^2+2$ ($v'=2x$) :
\begin{align*}
g'(x) &= \frac{1\cdot(x^2+2)-(x+1)\cdot 2x}{(x^2+2)^2} = \frac{x^2+2-2x^2-2x}{(x^2+2)^2} = \frac{-x^2-2x+2}{(x^2+2)^2}.
\end{align*}
\item $h(x)=\sqrt{3x-6}$ est définie lorsque $3x-6\geq 0 \iff x\geq 2$, et dérivable lorsque $3x-6>0$, c'est-à-dire sur $]2\,;\,+\infty[$ (la racine carrée n'est pas dérivable en $0$).
Avec $u(x)=3x-6$, $u'(x)=3$, et $(\sqrt{u})'=\dfrac{u'}{2\sqrt{u}}$ :
\[ h'(x) = \frac{3}{2\sqrt{3x-6}}. \]
(On peut simplifier : $h'(x)=\dfrac{3}{2\sqrt{3}\sqrt{x-2}}=\dfrac{\sqrt{3}}{2\sqrt{x-2}}$.)
\end{enumerate}
\textbf{Point de vigilance :} pour la question 3, ne pas oublier de préciser que la dérivabilité a lieu sur l'intervalle \emph{ouvert} $]2;+\infty[$ : en $x=2$, la dérivée « exploserait » (tangente verticale).
\end{corrige}

\begin{corrige}[Exercice 6 — Variations d'une fonction polynôme]
\begin{enumerate}
\item $f(x)=x^3-6x^2+9x+1$ est dérivable sur $\mathbb{R}$ :
\[ f'(x)=3x^2-12x+9 = 3(x^2-4x+3). \]
Vérifions la factorisation : $3(x-1)(x-3)=3(x^2-4x+3)=3x^2-12x+9$. C'est bien $f'(x)$.
\item $f'(x)=0 \iff x=1$ ou $x=3$. Le trinôme $x^2-4x+3$ a un coefficient dominant positif : il est positif à l'extérieur des racines, négatif entre les racines. Donc :
\begin{itemize}
\item $f'(x)>0$ sur $]-\infty\,;\,1[$ et sur $]3\,;\,+\infty[$ ;
\item $f'(x)<0$ sur $]1\,;\,3[$ ;
\item $f'(1)=f'(3)=0$.
\end{itemize}
\item Valeurs aux extrema :
\[ f(1)=1-6+9+1=5 \quad ; \quad f(3)=27-54+27+1=1. \]
Limites : $\lim_{x\to-\infty}f(x)=-\infty$ et $\lim_{x\to+\infty}f(x)=+\infty$ (terme dominant $x^3$).
\begin{center}
\begin{tabular}{|c|ccccccc|}
\hline
$x$ & $-\infty$ & & $1$ & & $3$ & & $+\infty$ \\
\hline
$f'(x)$ & & $+$ & $0$ & $-$ & $0$ & $+$ & \\
\hline
$f$ & $-\infty$ & $\nearrow$ & $5$ & $\searrow$ & $1$ & $\nearrow$ & $+\infty$ \\
\hline
\end{tabular}
\end{center}
$f$ admet un maximum local $5$ en $x=1$ et un minimum local $1$ en $x=3$.
\end{enumerate}
\end{corrige}

\begin{corrige}[Exercice 7 — Fonction homographique]
\begin{enumerate}
\item $f(x)=\dfrac{x-2}{x+1}$ est définie lorsque $x+1\neq 0$, donc $D_f=\mathbb{R}\setminus\{-1\}=]-\infty\,;\,-1[\,\cup\,]-1\,;\,+\infty[$.
\item Avec $u=x-2$ ($u'=1$) et $v=x+1$ ($v'=1$) :
\[ f'(x)=\frac{1\cdot(x+1)-(x-2)\cdot 1}{(x+1)^2}=\frac{x+1-x+2}{(x+1)^2}=\frac{3}{(x+1)^2}. \]
Pour tout $x\in D_f$, $(x+1)^2>0$, donc $f'(x)>0$ : $f$ est strictement croissante sur \emph{chacun} des intervalles $]-\infty;-1[$ et $]-1;+\infty[$.
\item Limites : $\lim_{x\to-1^-}f(x)=+\infty$ (numérateur vers $-3<0$, dénominateur vers $0^-$), $\lim_{x\to-1^+}f(x)=-\infty$, et $\lim_{x\to\pm\infty}f(x)=1$ (quotient des termes dominants).
\begin{center}
\begin{tabular}{|c|ccc||ccc|}
\hline
$x$ & $-\infty$ & & $-1$ & $-1$ & & $+\infty$ \\
\hline
$f'(x)$ & & $+$ & & & $+$ & \\
\hline
$f$ & $1$ & $\nearrow$ & $+\infty$ & $-\infty$ & $\nearrow$ & $1$ \\
\hline
\end{tabular}
\end{center}
(La double barre en $-1$ matérialise la valeur interdite.)
\item Sur $]-\infty;-1[$, $f$ est continue et strictement croissante, à valeurs dans $]1;+\infty[$. Comme $3\in\,]1;+\infty[$, l'équation $f(x)=3$ admet une unique solution sur cet intervalle (théorème des valeurs intermédiaires).
Sur $]-1;+\infty[$, $f$ prend ses valeurs dans $]-\infty;1[$, qui ne contient pas $3$ : pas de solution.
Vérification par le calcul : $\dfrac{x-2}{x+1}=3 \iff x-2=3x+3 \iff -2x=5 \iff x=-\dfrac{5}{2}$, qui appartient bien à $]-\infty;-1[$.
\textbf{Conclusion :} l'équation admet une unique solution, $x=-\dfrac{5}{2}$.
\end{enumerate}
\textbf{Point de vigilance :} ne jamais écrire « $f$ est croissante sur $D_f$ » tout court : la réunion de deux intervalles de croissance n'est pas un intervalle de croissance (par exemple $f(-2)=4>f(0)=-2$ alors que $-2<0$). On énonce la croissance sur \emph{chaque} intervalle séparément.
\end{corrige}

\begin{corrige}[Exercice 8 — Tangente et variations]
\begin{enumerate}
\item $f(x)=\sqrt{2x+1}=\sqrt{u(x)}$ avec $u(x)=2x+1$. La fonction $u$ est dérivable sur $\mathbb{R}$, et $u(x)>0 \iff x>-\dfrac{1}{2}$. La fonction racine carrée étant dérivable sur $]0;+\infty[$, $f$ est dérivable sur $\left]-\dfrac{1}{2}\,;\,+\infty\right[$ et :
\[ f'(x)=\frac{u'(x)}{2\sqrt{u(x)}}=\frac{2}{2\sqrt{2x+1}}=\frac{1}{\sqrt{2x+1}}. \]
\item $f(4)=\sqrt{9}=3$ et $f'(4)=\dfrac{1}{\sqrt{9}}=\dfrac{1}{3}$. L'équation de la tangente au point d'abscisse $4$ est :
\[ y = f'(4)(x-4)+f(4) = \frac{1}{3}(x-4)+3 = \frac{1}{3}x+\frac{5}{3}. \]
\item Pour tout $x>-\dfrac{1}{2}$, $\sqrt{2x+1}>0$, donc $f'(x)>0$ : $f$ est strictement croissante sur $\left[-\dfrac{1}{2}\,;\,+\infty\right[$. De plus $f\left(-\dfrac{1}{2}\right)=0$ et $\lim_{x\to+\infty}f(x)=+\infty$.
\begin{center}
\begin{tabular}{|c|ccc|}
\hline
$x$ & $-\frac{1}{2}$ & & $+\infty$ \\
\hline
$f'(x)$ & & $+$ & \\
\hline
$f$ & $0$ & $\nearrow$ & $+\infty$ \\
\hline
\end{tabular}
\end{center}
\end{enumerate}
\textbf{Point de vigilance :} en $x=-\dfrac{1}{2}$, la fonction est définie ($f=0$) mais \emph{pas dérivable} : la courbe y admet une demi-tangente verticale. La ligne de $f'$ ne doit donc pas contenir de valeur en $-\frac{1}{2}$.
\end{corrige}

\begin{corrige}[Exercice 9 — Bac blanc : coût moyen d'une entreprise de Douala]
\begin{enumerate}
\item \begin{enumerate}
\item $200$ objets correspondent à $x=2$ (centaines d'objets).
\[ C(2)=2^3-12\times 2^2+60\times 2 = 8-48+120 = 80. \]
Le coût total est de $80$ milliers de FCFA, soit \textbf{80\,000 FCFA}.
\[ C_M(2)=\frac{C(2)}{2}=\frac{80}{2}=40 \text{ milliers de FCFA par centaine}, \]
soit $\dfrac{40\,000}{100}=\textbf{400 FCFA par objet}$.
\item Pour $x\in[1;10]$, $x\neq 0$, donc on peut simplifier par $x$ :
\[ C_M(x)=\frac{C(x)}{x}=\frac{x^3-12x^2+60x}{x}=\frac{x(x^2-12x+60)}{x}=x^2-12x+60. \]
\end{enumerate}
\item \begin{enumerate}
\item $C_M$ est un polynôme, dérivable sur $[1;10]$ : $C_M'(x)=2x-12=2(x-6)$.
$C_M'(x)=0 \iff x=6$ ; $C_M'(x)<0$ sur $[1;6[$ et $C_M'(x)>0$ sur $]6;10]$.
\item Valeurs : $C_M(1)=1-12+60=49$ ; $C_M(6)=36-72+60=24$ ; $C_M(10)=100-120+60=40$.
\begin{center}
\begin{tabular}{|c|ccccc|}
\hline
$x$ & $1$ & & $6$ & & $10$ \\
\hline
$C_M'(x)$ & & $-$ & $0$ & $+$ & \\
\hline
$C_M$ & $49$ & $\searrow$ & $24$ & $\nearrow$ & $40$ \\
\hline
\end{tabular}
\end{center}
\end{enumerate}
\item \begin{enumerate}
\item D'après le tableau, $C_M$ admet son minimum en $x=6$, c'est-à-dire pour une production de \textbf{600 objets}.
\item Le coût moyen minimal est $C_M(6)=24$ milliers de FCFA par centaine d'objets, soit $\dfrac{24\,000}{100}=\textbf{240 FCFA par objet}$.
\end{enumerate}
\end{enumerate}
\textbf{Barème indicatif (sur 5 points) :} 1.a) $1$ pt ; 1.b) $0,5$ pt ; 2.a) $1$ pt ; 2.b) $1$ pt ; 3.a) $0,75$ pt ; 3.b) $0,75$ pt.

\textbf{Points de vigilance :}
\begin{itemize}
\item Bien convertir les unités à chaque étape : $x$ en \emph{centaines}, $C$ en \emph{milliers de FCFA} ; la question 3.b) demande une réponse \emph{par objet} en FCFA.
\item La simplification par $x$ en 1.b) exige de justifier $x\neq 0$ : c'est garanti car $x\in[1;10]$.
\item Pour conclure en 3.a), citer le tableau de variation (minimum global atteint en $x=6$), pas seulement $C_M'(6)=0$ : une dérivée nulle ne suffit pas à garantir un minimum.
\end{itemize}
\end{corrige}

\begin{corrige}[Exercice 10 — Bac blanc : une inégalité classique]
\begin{enumerate}
\item $f(x)=x+\dfrac{1}{x}$ est dérivable sur $]0;+\infty[$ comme somme de fonctions dérivables :
\[ f'(x)=1-\frac{1}{x^2}=\frac{x^2-1}{x^2}=\frac{(x-1)(x+1)}{x^2}. \]
\item Sur $]0;+\infty[$ : $x^2>0$ et $x+1>0$, donc $f'(x)$ a le signe de $x-1$.
Ainsi $f'(x)<0$ sur $]0;1[$, $f'(1)=0$, $f'(x)>0$ sur $]1;+\infty[$.
\item Limites aux bornes : $\lim_{x\to 0^+}f(x)=+\infty$ (car $\frac{1}{x}\to+\infty$) et $\lim_{x\to+\infty}f(x)=+\infty$ (car $x\to+\infty$). Valeur : $f(1)=1+1=2$.
\begin{center}
\begin{tabular}{|c|ccccc|}
\hline
$x$ & $0$ & & $1$ & & $+\infty$ \\
\hline
$f'(x)$ & & $-$ & $0$ & $+$ & \\
\hline
$f$ & $+\infty$ & $\searrow$ & $2$ & $\nearrow$ & $+\infty$ \\
\hline
\end{tabular}
\end{center}
\item D'après le tableau de variation, $f$ admet sur $]0;+\infty[$ un minimum \emph{global} en $x=1$, de valeur $2$. Donc pour tout $x>0$ :
\[ x+\frac{1}{x}=f(x)\geq f(1)=2. \]
L'égalité a lieu si et seulement si le minimum est atteint, c'est-à-dire pour $x=1$.
\item Notons $x$ un côté du rectangle (en m, $x>0$) et $L$ l'autre côté. L'aire vaut $100$ m$^2$ : $x\times L=100$, donc $L=\dfrac{100}{x}$. Le périmètre est :
\[ P(x)=2(x+L)=2x+\frac{200}{x}=2\left(x+\frac{100}{x}\right). \]
$P$ est dérivable sur $]0;+\infty[$ et :
\[ P'(x)=2-\frac{200}{x^2}=\frac{2x^2-200}{x^2}=\frac{2(x-10)(x+10)}{x^2}. \]
Pour $x>0$, $P'(x)$ a le signe de $x-10$ : $P$ est décroissante sur $]0;10]$ et croissante sur $[10;+\infty[$. Elle admet donc un minimum en $x=10$.
Alors $L=\dfrac{100}{10}=10$ m : le terrain de périmètre minimal est le \textbf{carré de côté $10$ m}, de périmètre $P(10)=20+20=\textbf{40 m}$.
\end{enumerate}
\textbf{Barème indicatif (sur 6 points) :} 1) $1$ pt ; 2) $1$ pt ; 3) $1,5$ pt ; 4) $1$ pt ; 5) $1,5$ pt.

\textbf{Points de vigilance :}
\begin{itemize}
\item En 2), la justification du signe doit mentionner explicitement que $x^2>0$ et $x+1>0$ \emph{parce que} $x>0$ : c'est l'hypothèse du domaine qui fait tout.
\item En 4), il faut dire « minimum \emph{global} » et conclure par une inégalité valable \emph{pour tout} $x>0$, avec le cas d'égalité.
\item En 5), ne pas oublier le facteur $2$ dans le périmètre, et vérifier que la solution trouvée donne bien un rectangle d'aire $100$ m$^2$.
\end{itemize}
\end{corrige}

\begin{corrige}[Exercice 11 — Bac blanc : étude complète d'une fonction]
\begin{enumerate}
\item $f(x)=\dfrac{1}{3}x^3-x^2-3x+2$ est un polynôme, dérivable sur $\mathbb{R}$ :
\[ f'(x)=\frac{1}{3}\times 3x^2-2x-3=x^2-2x-3. \]
Vérification : $(x+1)(x-3)=x^2-3x+x-3=x^2-2x-3$. C'est bien $f'(x)$.
\item $f'(x)=0 \iff x=-1$ ou $x=3$. Trinôme de coefficient dominant positif : positif à l'extérieur des racines, négatif entre elles.
$f'(x)>0$ sur $]-\infty;-1[\,\cup\,]3;+\infty[$ ; $f'(x)<0$ sur $]-1;3[$.
\item Valeurs aux extrema :
\[ f(-1)=-\frac{1}{3}-1+3+2=\frac{11}{3} \quad ; \quad f(3)=\frac{27}{3}-9-9+2=9-16=-7. \]
Limites : $\lim_{x\to-\infty}f(x)=-\infty$, $\lim_{x\to+\infty}f(x)=+\infty$ (terme dominant $\frac{1}{3}x^3$).
\begin{center}
\begin{tabular}{|c|ccccccc|}
\hline
$x$ & $-\infty$ & & $-1$ & & $3$ & & $+\infty$ \\
\hline
$f'(x)$ & & $+$ & $0$ & $-$ & $0$ & $+$ & \\
\hline
$f$ & $-\infty$ & $\nearrow$ & $\frac{11}{3}$ & $\searrow$ & $-7$ & $\nearrow$ & $+\infty$ \\
\hline
\end{tabular}
\end{center}
Maximum local $\dfrac{11}{3}$ en $x=-1$ ; minimum local $-7$ en $x=3$.
\item $f(0)=2$ et $f'(0)=-3$. Tangente au point d'abscisse $0$ :
\[ y=f'(0)(x-0)+f(0) \iff y=-3x+2. \]
\item \begin{enumerate}
\item Sur $[4;5]$, $f$ est continue (polynôme) et strictement croissante (car $f'>0$ sur $]3;+\infty[$, donc sur $]4;5[$). De plus :
\[ f(4)=\frac{64}{3}-16-12+2=\frac{64-48-36+6}{3}=-\frac{14}{3}<0, \]
\[ f(5)=\frac{125}{3}-25-15+2=\frac{125-75-45+6}{3}=\frac{11}{3}>0. \]
Comme $0\in\left[-\dfrac{14}{3}\,;\,\dfrac{11}{3}\right]$, le théorème des valeurs intermédiaires (version « bijection » : $f$ continue et strictement monotone) garantit que l'équation $f(x)=0$ admet une \textbf{unique} solution $\alpha$ dans $[4;5]$.
\item Testons le milieu : $f(4{,}5)=\dfrac{91{,}125}{3}-20{,}25-13{,}5+2=30{,}375-31{,}75=-1{,}375<0$.
Comme $f(4{,}5)<0<f(5)$, on obtient l'encadrement d'amplitude $0{,}5$ :
\[ 4{,}5 < \alpha < 5. \]
\end{enumerate}
\item Allure de la courbe : on place le maximum $M\left(-1\,;\,\frac{11}{3}\right)$ et le minimum $m(3\,;\,-7)$ avec leurs tangentes horizontales, le point $A(0;2)$ avec la tangente $y=-3x+2$, et la courbe traverse l'axe des abscisses une fois entre $4{,}5$ et $5$ (ainsi que deux fois avant $-1$, d'après le tableau).
\begin{popfigure}
\begin{tikzpicture}[scale=0.75]
\draw[popline, ->] (-3.5,0) -- (5.5,0) node[right] {$x$};
\draw[popline, ->] (0,-8) -- (0,5) node[above] {$y$};
\foreach \x in {-3,-2,-1,1,2,3,4,5} \draw (\x,0.08) -- (\x,-0.08) node[below] {\small $\x$};
\foreach \y in {-7,-4,-1,2} \draw (0.08,\y) -- (-0.08,\y) node[left] {\small $\y$};
\draw[popblue, thick, domain=-2.6:5, samples=200] plot (\x, {\x*\x*\x/3 - \x*\x - 3*\x + 2});
\draw[poporange, dashed] (-1.8, 3.667) -- (0.6, 3.667);
\draw[poporange, dashed] (2, -7) -- (4.2, -7);
\draw[popgreen, dashed, domain=-1:1.6] plot (\x, {-3*\x + 2});
\fill[popblue] (-1, 3.667) circle (2.5pt) node[above left] {\small $M$};
\fill[popblue] (3, -7) circle (2.5pt) node[below right] {\small $m$};
\fill[popblue] (0, 2) circle (2.5pt) node[above right] {\small $A$};
\end{tikzpicture}
\legende{Courbe $\mathcal{C}$, tangentes horizontales en $M$ et $m$, tangente en $A(0;2)$.}
\end{popfigure}
\end{enumerate}
\textbf{Barème indicatif (sur 8 points) :} 1) $1$ pt ; 2) $1$ pt ; 3) $1,5$ pt ; 4) $1$ pt ; 5.a) $1,5$ pt ; 5.b) $1$ pt ; 6) $1$ pt.

\textbf{Points de vigilance :}
\begin{itemize}
\item En 5.a), les \textbf{trois hypothèses} doivent être citées : continuité, stricte monotonie (pour l'\emph{unicité}), et le changement de signe $f(4)<0<f(5)$ (pour l'\emph{existence}). Oublier la stricte monotonie ne donne que l'existence, pas l'unicité : c'est une erreur classique sanctionnée au Bac.
\item Les calculs de $f(4)$ et $f(5)$ doivent être présentés proprement, avec mise au même dénominateur.
\item En 6), les tangentes horizontales se placent aux extrema ($x=-1$ et $x=3$), pas ailleurs ; la tangente en $0$ doit avoir un coefficient directeur \emph{négatif} ($-3$), donc être descendante.
\end{itemize}
\end{corrige}

\newpage

\coursec{Fiche bilan}

\coursec{Dérivées et sens de variation}

\begin{popfigure}
\begin{tikzpicture}[mindmap, grow cyclic,
  every node/.style={concept, text=white, font=\small\bfseries},
  root concept/.append style={concept color=popink, font=\bfseries, minimum size=2.8cm},
  level 1 concept/.append style={level distance=3.8cm, sibling angle=60, font=\footnotesize\bfseries},
  level 2 concept/.append style={level distance=2.6cm, font=\scriptsize\bfseries}]
\node[root concept]{Dérivées et variations}
  child[concept color=popblue]{ node {Nombre dérivé \& Tangente}
    child { node[concept color=popblue!70] {Limite du taux d'accroissement} }
    child { node[concept color=popblue!70] {Équation : $y=f'(a)(x-a)+f(a)$} }
  }
  child[concept color=poporange]{ node {Fonctions usuelles}
    child { node[concept color=poporange!70] {$x^n,\ \frac{1}{x},\ \sqrt{x}$} }
    child { node[concept color=poporange!70] {Ensembles de dérivabilité} }
  }
  child[concept color=popgreen]{ node {Opérations algébriques}
    child { node[concept color=popgreen!70] {Somme, Produit, Quotient} }
    child { node[concept color=popgreen!70] {Règle : $(uv)'=u'v+uv'$} }
  }
  child[concept color=poppurple]{ node {Fonctions composées}
    child { node[concept color=poppurple!70] {Chaîne : $(u^n)'=nu'u^{n-1}$} }
    child { node[concept color=poppurple!70] {$(\sqrt{u})'=\frac{u'}{2\sqrt{u}}$} }
  }
  child[concept color=poppink]{ node {Signe de $f'$ \& Variations}
    child { node[concept color=poppink!70] {$f'\ge 0 \Rightarrow f$ croissante} }
    child { node[concept color=poppink!70] {Extremums locaux} }
  }
  child[concept color=popgold]{ node {Tableau de variation \& Optimisation}
    child { node[concept color=popgold!70] {Valeurs exactes aux bornes} }
    child { node[concept color=popgold!70] {Problèmes concrets (coût, volume)} }
  };
\end{tikzpicture}
\legende{Carte mentale de la leçon : six compétences clés autour de la dérivation (programme Terminale A).}
\end{popfigure}

\courssub{Rappels : nombre dérivé et tangente}

\begin{definition}[Nombre dérivé en un point]
Soit $f$ une fonction définie sur un intervalle $I$ et $a \in I$. Si la limite 
\[
\lim_{h \to 0} \frac{f(a+h)-f(a)}{h}
\]
existe et est finie, on dit que $f$ est \cle{dérivable} en $a$. Cette limite est appelée \cle{nombre dérivé} de $f$ en $a$ et se note $f'(a)$.
\end{definition}

\begin{definition}[Fonction dérivée]
Si $f$ est dérivable en tout point d'un intervalle $I$, la fonction $f'$ qui à tout $x \in I$ associe $f'(x)$ est la \cle{fonction dérivée} de $f$ sur $I$. L'ensemble des points où $f$ est dérivable s'appelle l'\cle{ensemble de dérivabilité} de $f$.
\end{definition}

\begin{propriete}[Tangente à une courbe]
Soit $f$ dérivable en $a$. La tangente à la courbe $\mathcal{C}_f$ au point $A(a\,;\,f(a))$ est la droite d'équation :
\[
y = f'(a)(x-a) + f(a).
\]
Le coefficient directeur de cette tangente est $f'(a)$.
\end{propriete}

\courssub{Dérivées des fonctions usuelles et opérations}

\begin{propriete}[Tableau des dérivées usuelles -- À connaître par cœur]
\begin{center}
\begin{tabularx}{\linewidth}{|l|X|l|X|}
\hline
\rowcolor{popblueL} \textbf{Fonction $f(x)$} & \textbf{Ensemble de dérivabilité} & \textbf{Dérivée $f'(x)$} & \textbf{Remarque} \\
\hline
$k$ (constante) & $\mathbb{R}$ & $0$ &  \\
\hline
$x^n\ (n \in \mathbb{N}^*)$ & $\mathbb{R}$ & $n x^{n-1}$ & Valable pour $n \in \mathbb{Z}$ sur $\mathbb{R}^*$ \\
\hline
$\dfrac{1}{x} = x^{-1}$ & $\mathbb{R}^*$ & $-\dfrac{1}{x^2}$ & Cas particulier $n=-1$ \\
\hline
$\sqrt{x} = x^{1/2}$ & $\mathbb{R}_+^*$ & $\dfrac{1}{2\sqrt{x}}$ & Cas particulier $n=1/2$ \\
\hline
\end{tabularx}
\end{center}
\end{propriete}

\begin{propriete}[Règles de dérivation -- Opérations sur les fonctions]
Soient $u$ et $v$ deux fonctions dérivables sur un intervalle $I$.
\begin{center}
\begin{tabularx}{\linewidth}{|l|X|}
\hline
\rowcolor{popgreenL} \textbf{Opération} & \textbf{Dérivée sur $I$} \\
\hline
Somme : $u+v$ & $u'+v'$ \\
\hline
Produit par un scalaire : $k\,u\ (k \in \mathbb{R})$ & $k\,u'$ \\
\hline
Produit : $u\,v$ & $u'v + u v'$ \\
\hline
Quotient : $\dfrac{u}{v}$ & $\dfrac{u'v - u v'}{v^2}$ \quad ($v(x) \neq 0$ sur $I$) \\
\hline
\end{tabularx}
\end{center}
\end{propriete}

\begin{propriete}[Dérivée d'une fonction composée -- Règle de la chaîne]
Soient $u$ dérivable sur $I$ et $g$ dérivable sur $u(I)$. La composée $f = g \circ u$ est dérivable sur $I$ et :
\[
f'(x) = g'(u(x)) \cdot u'(x).
\]
\textbf{Cas particuliers fondamentaux (à connaître par cœur) :}
\begin{center}
\begin{tabularx}{\linewidth}{|l|l|X|}
\hline
\rowcolor{poppurpleL} \textbf{Fonction composée} & \textbf{Dérivée} & \textbf{Condition} \\
\hline
$[u(x)]^n\ (n \in \mathbb{Z})$ & $n\,u'(x)\,[u(x)]^{n-1}$ & $u(x) \neq 0$ si $n < 0$ \\
\hline
$\sqrt{u(x)}$ & $\dfrac{u'(x)}{2\sqrt{u(x)}}$ & $u(x) > 0$ \\
\hline
$\dfrac{1}{u(x)}$ & $-\dfrac{u'(x)}{[u(x)]^2}$ & $u(x) \neq 0$ \\
\hline
$u(ax+b)\ (a,b \in \mathbb{R})$ & $a\,u'(ax+b)$ &  \\
\hline
\end{tabularx}
\end{center}
\end{propriete}

\courssub{Signe de la dérivée et sens de variation}

\begin{propriete}[Théorème fondamental : lien signe de $f'$ -- variations de $f$]
Soit $f$ une fonction dérivable sur un intervalle $I$.
\begin{itemize}
  \item Si $f'(x) \geq 0$ sur $I$ (et $f'$ ne s'annule pas sur un intervalle non vide), alors $f$ est \cle{croissante} sur $I$.
  \item Si $f'(x) \leq 0$ sur $I$ (et $f'$ ne s'annule pas sur un intervalle non vide), alors $f$ est \cle{décroissante} sur $I$.
  \item Si $f'$ s'annule en $c \in I$ en \cle{changeant de signe} :
  \begin{itemize}
    \item de $+$ en $-$ : $f$ admet un \cle{maximum local} en $c$ ;
    \item de $-$ en $+$ : $f$ admet un \cle{minimum local} en $c$.
  \end{itemize}
\end{itemize}
\emph{(Propriété admise en série A : on l'utilise sans la démontrer.)}
\end{propriete}

\begin{experience}[Pourquoi le signe de $f'$ donne le sens de variation ?]
\textbf{L'idée (sans démonstration, la propriété est admise) :} en chaque point, $f'(x)$ est le coefficient directeur de la tangente à la courbe.
\begin{enumerate}
  \item Si $f'(x) > 0$ sur un intervalle, toutes les tangentes « montent » : la courbe monte, $f$ est croissante.
  \item Si $f'(x) < 0$ sur un intervalle, toutes les tangentes « descendent » : la courbe descend, $f$ est décroissante.
  \item Là où $f'$ s'annule en changeant de signe, la tangente est horizontale : c'est un sommet (maximum) ou un creux (minimum).
\end{enumerate}
\textbf{Remarque :} si $f'(x) > 0$ sauf en quelques points isolés où elle s'annule, $f$ reste strictement croissante.
\end{experience}

\courssub{Tableau de variation et optimisation}

\begin{methode}[Dresser un tableau de variation complet]
\textbf{Étape 1 :} Déterminer l'ensemble de définition $D_f$ et l'ensemble de dérivabilité $D_{f'}$.
\textbf{Étape 2 :} Calculer $f'(x)$ et le \textbf{simplifier} (mise au même dénominateur, \textbf{factorisation}).
\textbf{Étape 3 :} Résoudre $f'(x) = 0$ sur $D_{f'}$ pour trouver les valeurs critiques.
\textbf{Étape 4 :} Construire le \textbf{tableau de signes de $f'$} sur $D_{f'}$ (lignes : facteurs, $f'$, variations de $f$).
\textbf{Étape 5 :} Calculer les valeurs exactes (ou limites) de $f$ aux bornes de $D_f$ et aux points critiques.
\textbf{Étape 6 :} Remplir le \textbf{tableau de variations de $f$} : flèches $\nearrow$ (croissante) ou $\searrow$ (décroissante), valeurs $f(x)$ aux extremums et bornes.
\end{methode}

\begin{methode}[Résoudre un problème d'optimisation]
\textbf{Étape 1 :} Modéliser la grandeur à optimiser (surface, volume, coût, bénéfice...) par une fonction $f(x)$. Préciser l'intervalle $I$ (contraintes du problème : $x>0$, $x<L$, etc.).
\textbf{Étape 2 :} Étudier les variations de $f$ sur $I$ (tableau de variation).
\textbf{Étape 3 :} Identifier le maximum ou minimum global sur $I$ (aux points critiques ou aux bornes).
\textbf{Étape 4 :} \textbf{Conclure par une phrase} dans le contexte : « Le volume maximal est de ... pour $x = ...$ » ou « Le coût minimal est atteint pour une production de ... unités ».
\end{methode}

\begin{exemplebox}[Optimisation : Volume d'un colis pour CAMPOST]
\textbf{Énoncé :} On souhaite fabriquer un colis rectangulaire à base carrée de côté $x$ (en dm) et de hauteur $h$ (en dm) pour l'acheminer par CAMPOST. La somme de la hauteur et du périmètre de la base ne doit pas dépasser $100$ cm ($10$ dm). On veut maximiser le volume $V$.
\begin{enumerate}
  \item Exprimer $h$ en fonction de $x$. En déduire $V(x)$.
  \item Déterminer l'intervalle d'étude de $V$.
  \item Étudier les variations de $V$ et trouver les dimensions du colis de volume maximal.
\end{enumerate}
\textbf{Solution :}
\begin{enumerate}
  \item Contrainte : $h + 4x = 10 \Rightarrow h = 10 - 4x$. Volume $V(x) = x^2 h = x^2(10-4x) = 10x^2 - 4x^3$.
  \item $x > 0$ et $h > 0 \Rightarrow 10-4x > 0 \Rightarrow x < 2,5$. Intervalle : $I = ]0\,;\,2,5[$ (on pourra étudier sur $[0\,;\,2,5]$ par continuité).
  \item $V'(x) = 20x - 12x^2 = 4x(5 - 3x)$.
  \item $V'(x)=0 \Rightarrow x=0$ (borne) ou $x = \frac{5}{3} \approx 1,67$ dm.
  \item Tableau de signes : $V'(x) > 0$ sur $]0\,;\frac{5}{3}[$ et $V'(x) < 0$ sur $]\frac{5}{3}\,;\,2,5]$.
  \item $V$ croît puis décroît : maximum en $x = \frac{5}{3}$ dm. $h = 10 - 4\times\frac{5}{3} = \frac{10}{3}$ dm.
  \item $V_{\text{max}} = V(\frac{5}{3}) = \frac{25}{9} \times \frac{10}{3} = \frac{250}{27} \approx 9,26$ dm$^3$.
\end{enumerate}
\textbf{Conclusion :} Le volume maximal est $\frac{250}{27}$ dm$^3$ pour une base de côté $\frac{5}{3}$ dm et une hauteur de $\frac{10}{3}$ dm.
\end{exemplebox}

\begin{attention}[Pièges fréquents et points de vigilance -- \cle{Checklist erreurs Bac}]
\begin{itemize}
  \item \textbf{Quotient :} Numérateur $u'v - uv'$ (ordre !) et dénominateur $v^2$. Ne pas écrire $u'v'$.
  \item \textbf{Chaîne :} Ne jamais oublier le facteur $u'(x)$ : $(u^n)' = n u' u^{n-1}$, $(\sqrt{u})' = \frac{u'}{2\sqrt{u}}$.
  \item \textbf{Domaine :} $\sqrt{u}$ dérivable seulement si $u(x) > 0$ (strict). $\frac{u}{v}$ dérivable si $v(x) \neq 0$.
  \item \textbf{Factorisation :} Factoriser $f'(x)$ \textbf{systématiquement} pour faire le tableau de signes. Forme développée = piège.
  \item \textbf{Tableau de signes :} Une ligne par facteur. Signe de $f'$ = produit des signes. Zéros de $f'$ = valeurs critiques.
  \item \textbf{Tableau de variations :} Flèches $\nearrow$ si $f'>0$, $\searrow$ si $f'<0$. \textbf{Valeurs exactes} $f(x)$ aux extremums et bornes (pas d'approximations décimales sauf demande).
  \item \textbf{Redondance :} Ne pas confondre « $f$ croissante » ($f' \ge 0$) et « $f'$ croissante » ($f'' \ge 0$).
  \item \textbf{Optimisation :} Vérifier les bornes de l'intervalle (maximum/minimum global possible aux bornes même si $f' \neq 0$).
  \item \textbf{Rédaction :} « $f$ est croissante sur $I$ \textbf{car} $f'(x) \geq 0$ sur $I$ ». Justifier toujours par le signe de la dérivée.
\end{itemize}
\end{attention}

\begin{aretenir}[Checklist finale avant le Baccalauréat]
\begin{itemize}
  \item[$\square$] J'écris correctement la définition du nombre dérivé : $f'(a)=\lim_{h\to 0}\frac{f(a+h)-f(a)}{h}$.
  \item[$\square$] Je connais par cœur les dérivées de $x^n$, $\frac{1}{x}$, $\sqrt{x}$ \textbf{et leurs ensembles de dérivabilité}.
  \item[$\square$] J'applique sans erreur : $(u+v)'=u'+v'$, $(uv)'=u'v+uv'$, $\left(\frac{u}{v}\right)'=\frac{u'v-uv'}{v^2}$.
  \item[$\square$] Je gère les fonctions composées : $(u^n)'=n u' u^{n-1}$, $(\sqrt{u})'=\frac{u'}{2\sqrt{u}}$, $\left(\frac{1}{u}\right)'=-\frac{u'}{u^2}$.
  \item[$\square$] Je vérifie les conditions d'existence ($u>0$ pour $\sqrt{u}$, $v\neq 0$ pour $\frac{u}{v}$).
  \item[$\square$] Je factorise $f'(x)$ complètement avant d'étudier son signe.
  \item[$\square$] Je construis un tableau de signes de $f'$ rigoureux (lignes facteurs, double trait pour valeurs interdites).
  \item[$\square$] Je dresse le tableau de variations de $f$ avec flèches, valeurs exactes aux extremums et bornes.
  \item[$\square$] J'écris l'équation de la tangente : $y = f'(a)(x-a) + f(a)$.
  \item[$\square$] Pour un problème d'optimisation : je définis la fonction, son intervalle, j'étudie les variations, je conclus par une phrase contextualisée.
  \item[$\square$] Ma rédaction est structurée : « On a $f'(x) = \dots$, donc $f'(x) \geq 0$ sur $\dots$, donc $f$ est croissante sur $\dots$ ».
\end{itemize}
\end{aretenir}

\findecours

\end{document}
